% Sensibilisation sur la nécessité du renouvellement de l'ATP lors des exercices musculaires — SVT Tle D — Cours signé OnBuch+
% Programme officiel MINESEC. Compiler avec : tectonic N02-restauration-atp-exercice-musculaire.tex
\newcommand{\DOCMATIERE}{SVT}
\newcommand{\DOCNIVEAU}{Tle D}
\newcommand{\DOCMODULE}{Module I : Le Monde Vivant — Métabolisme énergétique des cellules musculaires}
\newcommand{\DOCLECON}{N02}
\newcommand{\DOCTITRE}{Sensibilisation sur la nécessité du renouvellement de l'ATP lors des exercices musculaires}
% OnBuch+ — préambule « cours signés » ENRICHI (compilé avec tectonic / XeLaTeX)
% Système visuel « Pop » d'OnBuch+ : fond crème (jamais blanc), bordures encre
% épaisses, ombres « dures » décalées, titres Archivo Black en capitales.
% Version MATHÉMATIQUES : graphes (pgfplots/tikz), boîtes pédagogiques
% (définition, propriété, méthode, attention, le savais-tu, exercice résolu,
% exercice, TP, bilan), page de garde et signature OnBuch+ ancrée en bas.
%
% Macros attendues dans main.tex AVANT \input{preamble} :
%   \DOCMATIERE  \DOCNIVEAU  \DOCMODULE  \DOCLECON (numéro)  \DOCTITRE
\documentclass[11pt]{article}

\usepackage{fontspec}
\usepackage[french]{babel}
\usepackage[a4paper,margin=2cm,top=3.4cm,bottom=2.8cm,headheight=1.4cm,footskip=1.3cm]{geometry}
\usepackage{amsmath,amssymb,amsfonts}
\usepackage{mathtools} % \xmapsto, \coloneqq... souvent utilisés par les modèles
\usepackage{cancel} % simplifications barrées (bilans de forces)
% \abs{z} (module, valeur absolue) et \norm{u} (norme), utilisés par les modèles
\providecommand{\abs}{}\let\abs\relax\DeclarePairedDelimiter{\abs}{\lvert}{\rvert}
\providecommand{\norm}{}\let\norm\relax\DeclarePairedDelimiter{\norm}{\lVert}{\rVert}
\usepackage[table]{xcolor} % [table] : fournit aussi \rowcolors (alternance de lignes)
\usepackage{pagecolor}
\usepackage{tikz}
\usetikzlibrary{arrows.meta,positioning,calc,shapes.geometric,shapes.misc,shapes.arrows,decorations.pathmorphing,decorations.pathreplacing,decorations.markings,patterns,shadows,mindmap,trees,fit,backgrounds,matrix,intersections,angles,quotes}
\usepackage{pgfplots}
\usepackage[version=4]{mhchem} % équations nucléaires \ce{^{235}_{92}U}
\usepackage[european,siunitx]{circuitikz} % schémas électriques
\pgfplotsset{compat=1.18}
\usepgfplotslibrary{fillbetween,groupplots} % aires entre courbes (calcul intégral)
\usepackage{enumitem}
\usepackage{fancyhdr}
% [most] charge toutes les bibliothèques tcolorbox (listings, minted, raster,
% documentation...) et gonfle inutilement la mémoire TeX sur un document long
% (~300 boîtes) : on ne charge que ce qui est réellement utilisé.
\usepackage[skins,breakable]{tcolorbox}
\usepackage{graphicx}
\usepackage{colortbl}
\usepackage{booktabs}
\usepackage{tabularx}
\usepackage{diagbox} % case d'en-tête diagonale des tableaux à double entrée
% Colonnes extensibles centrée / gauche / droite (C, L, R), souvent utilisées par les modèles
\newcolumntype{C}{>{\centering\arraybackslash}X}
\newcolumntype{L}{>{\raggedright\arraybackslash}X}
\newcolumntype{R}{>{\raggedleft\arraybackslash}X}
\usepackage{array}
\usepackage{multicol}
\usepackage{siunitx}
% parse-numbers=false : accepte aussi \SI{2,5 \times 10^{-3}}{...}
\sisetup{locale=FR,output-decimal-marker={,},group-minimum-digits=4,parse-numbers=false}
\DeclareSIUnit\ohm{\ensuremath{\symup{\Omega}}} % U+2126 absent de la police
\DeclareSIUnit\division{div} % réglages d'oscilloscope (ms/div, V/div)
\DeclareSIUnit\tr{tr} % tours (vitesse de rotation, ex. \SI{1500}{\tr\per\minute})
\DeclareSIUnit\molar{M} % concentration molaire (ex. \SI{3}{\molar}, \SI{1}{\micro\molar})
\DeclareSIUnit\an{an} % année (le modèle confond parfois avec \year, un compteur TeX) — ex. \SI{5}{\kilo\meter\per\an}
\DeclareSIUnit\FCFA{FCFA} % franc CFA (ex. \SI{500}{\FCFA\per\kilo\gram})
\DeclareSIUnit\degreeBrix{\textdegree Brix} % teneur en sucre d'un sirop/jus (réfractométrie)
\DeclareSIUnit\month{mois} % le modèle utilise parfois \month, un compteur TeX, comme unité de durée
\usepackage{qrcode}
% \degree : commande fréquemment utilisée par les modèles pour un angle ou une
% température en dehors de \SI{}{} (non fournie par les paquets chargés).
\providecommand{\degree}{^{\circ}}
% \qed (fin de démonstration), utilisé par les modèles sans amsthm
\providecommand{\qed}{\hfill$\square$}
% \barre : double barre des valeurs interdites dans un tableau de variations (tkz-tab)
\providecommand{\barre}{\ensuremath{\|}}
% environnement proof (amsthm non chargé) utilisé par les modèles
\ifdefined\proof\else\newenvironment{proof}[1][Démonstration]{\par\noindent\textit{#1.}\ }{\qed\par}\fi
\usepackage{lastpage}
\usepackage{hyperref}
\hypersetup{hidelinks}

% ============================================================
% Palette « Pop » OnBuch+ — mêmes hex que la classe P (plus_ui.dart)
% ============================================================
\definecolor{poporange}{HTML}{F59321}
\definecolor{popdark}{HTML}{E07A0C}
\definecolor{popcream}{HTML}{FFF6E8}
\definecolor{popink}{HTML}{1C1714}
\definecolor{poppurple}{HTML}{7A5AE0}
\definecolor{popgreen}{HTML}{2E9E6A}
\definecolor{poppink}{HTML}{E0457B}
\definecolor{popblue}{HTML}{2F7DE1}
\definecolor{popgold}{HTML}{C98A12}
\definecolor{popmuted}{HTML}{8A7F76}
\definecolor{popline}{HTML}{EADFD2}
\definecolor{popgray}{HTML}{8A7F76}
\colorlet{popgrey}{popgray}
% Alias tolérants (le modèle nomme parfois les couleurs différemment)
\colorlet{popwhite}{white}
\colorlet{popblack}{popink}
\colorlet{popred}{poppink}
\colorlet{popyellow}{popgold}
% Teintes claires (fonds de tableaux/encadrés) — le modèle nomme parfois
% les couleurs avec un suffixe L (ex. popblueL) sans les avoir définies.
\colorlet{poporangeL}{poporange!20}
\colorlet{popdarkL}{popdark!20}
\colorlet{poppurpleL}{poppurple!20}
\colorlet{popgreenL}{popgreen!20}
\colorlet{poppinkL}{poppink!20}
\colorlet{popblueL}{popblue!20}
\colorlet{popgoldL}{popgold!20}
\colorlet{popredL}{popred!20}
\colorlet{popgrayL}{popgray!20}
% Teintes claires pour fonds de tableaux / zones
\colorlet{popblueL}{popblue!10!white}
\colorlet{poporangeL}{poporange!14!white}
\colorlet{popgreenL}{popgreen!12!white}
\colorlet{poppurpleL}{poppurple!12!white}
\colorlet{poppinkL}{poppink!12!white}
\colorlet{popgoldL}{popgold!14!white}

\pagecolor{popcream}

% ============================================================
% Typographie
% ============================================================
\defaultfontfeatures{Ligatures=TeX}
\setmainfont{PlusJakartaSans-Medium.ttf}[Path=../../../fonts/,BoldFont=PlusJakartaSans-Bold.ttf]
% Maths en sans-serif (Fira Math) avec chiffres et lettres droites de Plus
% Jakarta, pour que les formules se fondent dans le texte.
\usepackage[math-style=ISO,bold-style=ISO]{unicode-math}
\setmathfont{FiraMath-Regular.otf}[Path=../../../fonts/]
\setmathfont{PlusJakartaSans-Medium.ttf}[Path=../../../fonts/,range={up/num,up/{Latin,latin},\mathup}]
\usepackage{icomma} % virgule décimale sans espace en mode math
% Caractères mathématiques Unicode absents des polices de texte (Plus Jakarta,
% Archivo Black) : rendus par leur équivalent mathématique, en texte comme en
% formule — sinon ils s'affichent en carré vide (ex. « Arithmétique dans ℤ »).
\usepackage{newunicodechar}
\newunicodechar{ℝ}{\ensuremath{\mathbb{R}}}
\newunicodechar{ℤ}{\ensuremath{\mathbb{Z}}}
\newunicodechar{ℂ}{\ensuremath{\mathbb{C}}}
\newunicodechar{ℕ}{\ensuremath{\mathbb{N}}}
\newunicodechar{ℚ}{\ensuremath{\mathbb{Q}}}
\newunicodechar{𝔻}{\ensuremath{\mathbb{D}}}
\newunicodechar{α}{\ensuremath{\alpha}}
\newunicodechar{β}{\ensuremath{\beta}}
\newunicodechar{γ}{\ensuremath{\gamma}}
\newunicodechar{δ}{\ensuremath{\delta}}
\newunicodechar{ε}{\ensuremath{\varepsilon}}
\newunicodechar{θ}{\ensuremath{\theta}}
\newunicodechar{λ}{\ensuremath{\lambda}}
\newunicodechar{μ}{\ensuremath{\mu}}
\newunicodechar{σ}{\ensuremath{\sigma}}
\newunicodechar{τ}{\ensuremath{\tau}}
\newunicodechar{φ}{\ensuremath{\varphi}}
\newunicodechar{ψ}{\ensuremath{\psi}}
\newunicodechar{ω}{\ensuremath{\omega}}
\newunicodechar{Σ}{\ensuremath{\Sigma}}
% majuscules produites par \MakeUppercase dans les titres (α-aminés -> Α)
\newunicodechar{Α}{\ensuremath{\alpha}}
\newunicodechar{Β}{\ensuremath{\beta}}
\newunicodechar{Γ}{\ensuremath{\Gamma}}
\newunicodechar{Δ}{\ensuremath{\Delta}}
\newunicodechar{Ε}{\ensuremath{\varepsilon}}
\newunicodechar{Θ}{\ensuremath{\Theta}}
\newunicodechar{Λ}{\ensuremath{\Lambda}}
\newunicodechar{Μ}{\ensuremath{\mu}}
\newunicodechar{Τ}{\ensuremath{\tau}}
\newunicodechar{Φ}{\ensuremath{\Phi}}
\newunicodechar{Ψ}{\ensuremath{\Psi}}
\newunicodechar{Ω}{\ensuremath{\Omega}}
\newunicodechar{↦}{\ensuremath{\mapsto}}
\newunicodechar{∈}{\ensuremath{\in}}
\newunicodechar{∉}{\ensuremath{\notin}}
\newunicodechar{∧}{\ensuremath{\wedge}}
\newunicodechar{⇔}{\ensuremath{\Leftrightarrow}}
\newunicodechar{⟺}{\ensuremath{\Longleftrightarrow}}
\newunicodechar{⇒}{\ensuremath{\Rightarrow}}
\newunicodechar{⟹}{\ensuremath{\Longrightarrow}}
\newunicodechar{∘}{\ensuremath{\circ}}
\newunicodechar{∪}{\ensuremath{\cup}}
\newunicodechar{∩}{\ensuremath{\cap}}
\newunicodechar{≡}{\ensuremath{\equiv}}
\newunicodechar{⊂}{\ensuremath{\subset}}
\newunicodechar{⊥}{\ensuremath{\perp}}
\newunicodechar{∃}{\ensuremath{\exists}}
\newunicodechar{∀}{\ensuremath{\forall}}
\newunicodechar{∛}{\ensuremath{\sqrt[3]{\,}}}
\newunicodechar{∜}{\ensuremath{\sqrt[4]{\,}}}
\newunicodechar{⃗}{\ensuremath{{}^{\to}}}
\newunicodechar{ⁿ}{\ifmmode^{n}\else\textsuperscript{\ensuremath{n}}\fi}
\newunicodechar{ˣ}{\ifmmode^{x}\else\textsuperscript{\ensuremath{x}}\fi}
\newunicodechar{ᵇ}{\ifmmode^{b}\else\textsuperscript{\ensuremath{b}}\fi}
\newunicodechar{ʳ}{\ifmmode^{r}\else\textsuperscript{\ensuremath{r}}\fi}
\newunicodechar{ᵘ}{\ifmmode^{u}\else\textsuperscript{\ensuremath{u}}\fi}
\newunicodechar{ʸ}{\ifmmode^{y}\else\textsuperscript{\ensuremath{y}}\fi}
\newunicodechar{ᵃ}{\ifmmode^{a}\else\textsuperscript{\ensuremath{a}}\fi}
\newunicodechar{ᵉ}{\ifmmode^{e}\else\textsuperscript{\ensuremath{e}}\fi}
\newunicodechar{ᵏ}{\ifmmode^{k}\else\textsuperscript{\ensuremath{k}}\fi}
\newunicodechar{ᵖ}{\ifmmode^{p}\else\textsuperscript{\ensuremath{p}}\fi}
\newunicodechar{ᴺ}{\ifmmode^{N}\else\textsuperscript{\ensuremath{N}}\fi}
\newunicodechar{ᶜ}{\ifmmode^{c}\else\textsuperscript{\ensuremath{c}}\fi}
\newunicodechar{ʰ}{\ifmmode^{h}\else\textsuperscript{\ensuremath{h}}\fi}
\newunicodechar{ᵠ}{\ifmmode^{q}\else\textsuperscript{\ensuremath{q}}\fi}
\newunicodechar{ₙ}{\ifmmode_{n}\else\textsubscript{\ensuremath{n}}\fi}
\newunicodechar{ₐ}{\ifmmode_{a}\else\textsubscript{\ensuremath{a}}\fi}
\newunicodechar{ᵢ}{\ifmmode_{i}\else\textsubscript{\ensuremath{i}}\fi}
\newunicodechar{ᵣ}{\ifmmode_{r}\else\textsubscript{\ensuremath{r}}\fi}
\newunicodechar{ₖ}{\ifmmode_{k}\else\textsubscript{\ensuremath{k}}\fi}
\newunicodechar{ᵦ}{\ifmmode_{\beta}\else\textsubscript{\ensuremath{\beta}}\fi}
\newunicodechar{ₓ}{\ifmmode_{x}\else\textsubscript{\ensuremath{x}}\fi}
\newfontfamily\popdisplay{ArchivoBlack-Regular.ttf}[Path=../../../fonts/]
\newfontfamily\poplogo{SpaceGrotesk-Bold.ttf}[Path=../../../fonts/]
\newfontfamily\popsemi{PlusJakartaSans-SemiBold.ttf}[Path=../../../fonts/]
\newfontfamily\popxbold{PlusJakartaSans-ExtraBold.ttf}[Path=../../../fonts/]
\newfontfamily\popxboldls{PlusJakartaSans-ExtraBold.ttf}[Path=../../../fonts/,LetterSpace=3.0]

\setlist[enumerate]{leftmargin=1.7em,itemsep=5pt,topsep=4pt}
\setlist[itemize]{leftmargin=1.7em,itemsep=4pt,topsep=4pt,label={\color{poporange}\textbullet}}
\setlength{\parindent}{0pt}
\setlength{\parskip}{5pt}
\renewcommand{\arraystretch}{1.25}

% pgfplots : style Pop par défaut
\pgfplotsset{
  every axis/.append style={axis line style={popink,line width=1pt},tick style={popink},
    grid style={popline,line width=0.5pt},label style={font=\small},tick label style={font=\footnotesize}},
  popcurve/.style={poporange,line width=1.8pt,no markers,smooth},
}
% Même style disponible dans un \draw TikZ simple (hors pgfplots)
\tikzset{popcurve/.style={poporange,line width=1.8pt}}
\tikzset{popgrid/.style={popline,very thin}}
\colorlet{popcurve}{poporange} % le nom du style est parfois utilisé comme couleur

% ============================================================
% Pill / squiggle
% ============================================================
\newcommand{\poppill}[3]{%
  \tikz[baseline=-0.55ex]\node[draw=popink,line width=1.2pt,rounded corners=2.6pt,%
    fill=#2,inner xsep=6.5pt,inner ysep=3pt]{%
    \popxboldls\fontsize{7}{7}\selectfont\color{#3}\MakeUppercase{#1}};%
}
\newcommand{\popsquiggle}[1]{%
  \tikz[baseline=-0.4ex]{
    \draw[#1,line width=2.6pt,line cap=round]
      plot [smooth] coordinates {(0,0.09) (0.34,0.01) (0.68,0.21) (1.02,0.01) (1.36,0.10)};
  }%
}
% Mot-clé surligné (marqueur orange sous le texte)
\newcommand{\cle}[1]{{\popxbold\color{popdark}#1}}

% ============================================================
% En-tête / pied
% ============================================================
\pagestyle{fancy}
\fancyhf{}
\renewcommand{\headrulewidth}{0pt}
\renewcommand{\footrulewidth}{0pt}
\lhead{\raisebox{-0.15ex}{\poplogo\fontsize{13}{13}\selectfont\color{popink}OnBuch\color{poporange}+}}
\rhead{\poppill{\DOCMATIERE\ \textbullet\ \DOCNIVEAU\ \textbullet\ Leçon \DOCLECON}{white}{popink}}
\lfoot{\popxbold\fontsize{7.6}{7.6}\selectfont\color{popmuted}\MakeUppercase{Cours signé OnBuch+}\ \textbullet\ {\popsemi\DOCTITRE}}
\rfoot{\tikz[baseline=-0.6ex]\node[rounded corners=8.5pt,fill=popink,minimum height=17pt,minimum width=17pt,inner xsep=5pt,inner sep=0]{\popxbold\fontsize{8}{8}\selectfont\color{popcream}\thepage\,/\,\pageref*{LastPage}};}

% ============================================================
% Titres
% ============================================================
\newcommand{\chaptitre}[2]{%
  \begin{tcolorbox}[enhanced,colback=poporange,colframe=popink,boxrule=2.6pt,arc=11pt,
      drop shadow=popink,left=15pt,right=15pt,top=12pt,bottom=11pt,before skip=3pt,after skip=18pt]
    \poppill{#2}{popink}{popcream}\par\vspace{9pt}
    {\raggedright\hyphenpenalty=10000\exhyphenpenalty=10000\popdisplay\fontsize{22}{24}\selectfont\color{popcream}\MakeUppercase{#1}\par}
  \end{tcolorbox}
}
% Section numérotée : pastille orange avec le numéro + titre Archivo
\newcounter{popsec}
\newcommand{\coursec}[1]{%
  \stepcounter{popsec}\setcounter{popsub}{0}%
  \par\vspace{16pt}\noindent
  \begin{minipage}{\linewidth}
  \tikz[baseline=-0.7ex]\node[draw=popink,line width=1.6pt,fill=poporange,rounded corners=4pt,minimum size=22pt,inner sep=0]{\popdisplay\fontsize{12}{12}\selectfont\color{popcream}\thepopsec};%
  \hspace{8pt}{\popdisplay\fontsize{15}{17}\selectfont\color{popink}\MakeUppercase{#1}}\par
  \vspace{4pt}\noindent{\color{poporange}\rule{36pt}{3.6pt}}
  \end{minipage}\par\nopagebreak\vspace{8pt}
}
\newcounter{popsub}
\newcommand{\courssub}[1]{%
  \stepcounter{popsub}%
  \par\vspace{9pt}\noindent{\popxbold\fontsize{11.5}{13}\selectfont\color{popink}{\color{poporange}\thepopsec.\thepopsub}\hspace{6pt}#1}\par\nopagebreak\vspace{4pt}
}

% ============================================================
% Boîtes pédagogiques — même recette : fond blanc, bordure épaisse colorée,
% ombre dure encre, badge plein. Titre optionnel [#1].
% ============================================================
\tcbset{popbase/.style={breakable,enhanced,colback=white,boxrule=2.3pt,arc=9pt,
  shadow={3pt}{-3pt}{0pt}{popink},left=14pt,right=14pt,top=11pt,bottom=11pt,before skip=11pt,after skip=15pt,
  fontupper=\popsemi\fontsize{10.6}{14.5}\selectfont\color{popink},
  fontlower=\popsemi\fontsize{10.6}{14.5}\selectfont\color{popink}}}
% Coche dessinée (le glyphe ✓ n'existe pas dans les polices du document)
\newcommand{\popcheck}[1]{\tikz[baseline=-0.5ex]\draw[#1,line width=1.6pt,line cap=round,line join=round](-0.13,0.0)--(-0.04,-0.09)--(0.13,0.1);}
\providecommand{\coche}{\popcheck{popgreen}}
\providecommand{\resultat}[1]{\par\smallskip\noindent\fcolorbox{popgreen}{popgreenL}{\parbox{\dimexpr\linewidth-2\fboxsep-2\fboxrule\relax}{\textbf{Résultat.} #1}}\par\smallskip}
\newcommand{\popboxhead}[3]{\poppill{#1}{#2}{white}\ifx\relax#3\relax\else\hspace{6pt}{\popxbold\fontsize{10.6}{12}\selectfont\color{popink}#3}\fi\par\vspace{7pt}}

\newtcolorbox{aretenir}[1][]{popbase,colframe=popgreen,before upper={\popboxhead{À retenir}{popgreen}{#1}}}
\newtcolorbox{exemplebox}[1][]{popbase,colframe=poporange,before upper={\popboxhead{Exemple}{poporange}{#1}}}
\newtcolorbox{definition}[1][]{popbase,colframe=popblue,before upper={\popboxhead{Définition}{popblue}{#1}}}
\newtcolorbox{propriete}[1][]{popbase,colframe=poppurple,before upper={\popboxhead{Propriété}{poppurple}{#1}}}
\newtcolorbox{methode}[1][]{popbase,colframe=popgold,before upper={\popboxhead{Méthode}{popgold}{#1}}}
\newtcolorbox{attention}[1][]{popbase,colframe=poppink,before upper={\popboxhead{Attention}{poppink}{#1}}}
\newtcolorbox{experience}[1][]{popbase,colframe=popblue,colback=popblueL,before upper={\popboxhead{Expérience}{popblue}{#1}}}
% « Le savais-tu ? » : carte violette pleine (contexte, culture, Cameroun)
\newtcolorbox{savaistu}[1][]{popbase,colback=poppurple,colframe=popink,
  fontupper=\popsemi\fontsize{10.4}{14.2}\selectfont\color{white},
  before upper={\poppill{Le savais-tu\,?}{white}{poppurple}\ifx\relax#1\relax\else\hspace{6pt}{\popxbold\color{white}#1}\fi\par\vspace{7pt}}}
% Exercice résolu : énoncé en haut, \tcblower puis la solution sur fond crème
\newcounter{popexo}\newcounter{popexr}
\newtcolorbox{exoresolu}[1][]{popbase,colframe=poporange,colbacklower=popcream,
  segmentation style={popink,line width=1.2pt,dashed},
  before upper={\stepcounter{popexr}\popboxhead{Exercice résolu \thepopexr}{poporange}{#1}},
  before lower={\poppill{Solution}{popink}{popcream}\par\vspace{6pt}}}
% Exercice (non résolu en place) — la correction est dans la partie Corrigés
\newtcolorbox{exercice}[1][]{popbase,colframe=popink,
  before upper={\stepcounter{popexo}\popboxhead{Exercice \thepopexo}{popink}{#1}}}
% Corrigé
\newtcolorbox{corrige}[1][]{popbase,colframe=popgreen,colback=popgreenL,
  before upper={\popboxhead{Corrigé}{popgreen}{#1}}}
% Objectifs / prérequis / bilan
\newtcolorbox{objectifs}{popbase,colframe=popink,colback=poporangeL,
  before upper={\popboxhead{Objectifs de la leçon}{popink}{}}}
\newtcolorbox{prerequis}{popbase,colframe=popink,colback=popblueL,
  before upper={\popboxhead{Prérequis}{popblue}{}}}
\newtcolorbox{bilan}{popbase,colframe=popink,colback=white,boxrule=2.8pt,
  before upper={\popboxhead{Fiche bilan}{poporange}{L'essentiel à connaître pour l'examen}}}
% Figure encadrée avec légende
\newtcolorbox{popfigure}[1][]{enhanced,breakable=false,colback=white,colframe=popline,boxrule=1.4pt,arc=8pt,
  left=10pt,right=10pt,top=10pt,bottom=8pt,before skip=10pt,after skip=12pt,halign=center,
  lower separated=false,halign lower=center,
  fontlower=\popsemi\fontsize{9}{11.5}\selectfont\color{popmuted},
  #1}
\newcommand{\legende}[1]{\par\vspace{6pt}{\popsemi\fontsize{9}{11.5}\selectfont\color{popmuted}#1}}

% ============================================================
% Signature OnBuch+
% ============================================================
\newcommand{\poplogoinline}[1]{{\poplogo\fontsize{#1}{#1}\selectfont\color{popink}OnBuch\color{poporange}+}}
\newcommand{\signatureonbuch}{%
  \begin{tcolorbox}[enhanced,colback=popink,colframe=popink,boxrule=2.2pt,arc=10pt,
      drop shadow=poporange,left=14pt,right=14pt,top=10pt,bottom=10pt,before skip=0pt,after skip=0pt]
    \tikz[baseline=-0.7ex]\node[circle,fill=popgreen,draw=popcream,line width=1.3pt,minimum size=22pt,inner sep=0]{\popcheck{white}};%
    \hspace{9pt}\begin{minipage}[c]{0.62\linewidth}
      {\popxbold\fontsize{12}{13}\selectfont\color{popcream}Signé\ \textbullet\ {\poplogo OnBuch\color{poporange}+}}\par\vspace{2pt}
      {\raggedright\popsemi\fontsize{8}{10.5}\selectfont\color{popline}\MakeUppercase{Équipe pédagogique OnBuch+}\\\MakeUppercase{Conforme au programme officiel MINESEC}\par}
    \end{minipage}\hfill
    {\poplogo\fontsize{22}{22}\selectfont\color{popcream}OnBuch\color{poporange}+}
  \end{tcolorbox}%
}

% ============================================================
% Page de garde
% #1 titre · #2 matière · #3 niveau · #4 module · #5 n° leçon · #6 durée ·
% #7 description / objectifs (texte ou itemize)
% ============================================================
\newcommand{\pagegarde}[7]{%
  \thispagestyle{empty}
  \newgeometry{a4paper,margin=1.7cm,top=1.6cm,bottom=1.4cm}%
  \begin{tikzpicture}[remember picture,overlay]
    \fill[poporange!18] ([xshift=-3.2cm,yshift=-3.6cm]current page.north east) circle (4.6cm);
    \fill[poppurple!14] ([xshift=2.4cm,yshift=7.6cm]current page.south west) circle (3.8cm);
  \end{tikzpicture}%
  {\poplogo\fontsize{30}{30}\selectfont\color{popink}OnBuch\color{poporange}+}\hfill
  \raisebox{0.6ex}{\poppill{Cours signé \textbullet\ Premium}{popink}{popcream}}\par
  \vspace{1pt}\popsquiggle{poppurple}\par
  \vspace{8pt}
  \begin{tcolorbox}[enhanced,colback=poporange,colframe=popink,boxrule=2.8pt,arc=13pt,
      drop shadow=popink,left=18pt,right=18pt,top=12pt,bottom=13pt,after skip=10pt]
    \poppill{#2 \textbullet\ #3}{popink}{popcream}\hspace{5pt}\poppill{#4}{white}{popink}\par\vspace{8pt}
    {\popdisplay\fontsize{14}{14}\selectfont\color{popink}LEÇON #5}\par\vspace{4pt}
    {\raggedright\hyphenpenalty=10000\exhyphenpenalty=10000\popdisplay\fontsize{25}{27}\selectfont\color{popcream}\MakeUppercase{#1}\par}
    \vspace{8pt}
    {\popxbold\fontsize{9.5}{9.5}\selectfont\color{popink}\MakeUppercase{Durée conseillée : #6}}
  \end{tcolorbox}
  \begin{tcolorbox}[enhanced,colback=white,colframe=popink,boxrule=2.2pt,arc=10pt,
      drop shadow=popink,left=16pt,right=16pt,top=10pt,bottom=10pt,after skip=8pt]
    \poppill{Dans cette leçon}{poppurple}{white}\par\vspace{5pt}
    \popsemi\fontsize{9.3}{12.6}\selectfont\color{popink}#7
  \end{tcolorbox}
  \noindent\begin{minipage}[t]{0.487\linewidth}
    \begin{tcolorbox}[enhanced,colback=popgreen,colframe=popink,boxrule=2.2pt,arc=10pt,
        drop shadow=popink,left=11pt,right=11pt,top=10pt,bottom=10pt,height=3.1cm,valign=center]
      \tcbox[colback=white,colframe=popink,boxrule=1.3pt,arc=5pt,boxsep=0pt,nobeforeafter,
        left=3pt,right=3pt,top=3pt,bottom=3pt,valign=center]{\qrcode[height=1.9cm]{https://play.google.com/store/apps/details?id=com.luvvix.onbuch&hl=fr}}%
      \hspace{8pt}\begin{minipage}[c]{0.5\linewidth}
        {\popdisplay\fontsize{11}{12.5}\selectfont\color{white}\MakeUppercase{Télécharge OnBuch}}\par\vspace{4pt}
        {\raggedright\popsemi\fontsize{8.4}{11}\selectfont\color{white}Gratuit \textbullet\ Google Play\\Tuteur IA, annales, concours, résultats\par}
      \end{minipage}
    \end{tcolorbox}
  \end{minipage}\hfill
  \begin{minipage}[t]{0.487\linewidth}
    \begin{tcolorbox}[enhanced,colback=poppurple,colframe=popink,boxrule=2.2pt,arc=10pt,
        drop shadow=popink,left=11pt,right=11pt,top=10pt,bottom=10pt,height=3.1cm,valign=center]
      \tikz[baseline=-0.7ex]\node[circle,fill=white,draw=popink,line width=1.3pt,minimum size=1.6cm,inner sep=0]{\popdisplay\fontsize{24}{24}\selectfont\color{poppurple}+};%
      \hspace{8pt}\begin{minipage}[c]{0.6\linewidth}
        {\popdisplay\fontsize{11}{12.5}\selectfont\color{white}\MakeUppercase{Passe à OnBuch+}}\par\vspace{4pt}
        {\raggedright\popsemi\fontsize{8.4}{11}\selectfont\color{white}Tous les cours signés, Tuteur IA illimité, fiches PDF — onglet OnBuch+ de l'app\par}
      \end{minipage}
    \end{tcolorbox}
  \end{minipage}
  \vspace{8pt}
  \popcontenu
  \vfill
  \signatureonbuch
  \restoregeometry
  \newpage
}

% Rangée « Ce cours contient » (page de garde)
\newcommand{\poptile}[3]{% #1 couleur #2 gros texte #3 légende
  \begin{tcolorbox}[enhanced,colback=#1,colframe=popink,boxrule=2pt,arc=9pt,shadow={2.5pt}{-2.5pt}{0pt}{popink},
      left=6pt,right=6pt,top=9pt,bottom=9pt,halign=center,valign=center,height=1.75cm,width=\linewidth,nobeforeafter]
    {\popdisplay\fontsize{12.5}{13}\selectfont\color{white}\MakeUppercase{#2}}\par\vspace{3pt}
    {\centering\popsemi\fontsize{7.8}{9.5}\selectfont\color{white}#3\par}
  \end{tcolorbox}}
\newcommand{\popcontenu}{%
  \par\noindent{\popxboldls\fontsize{7.5}{7.5}\selectfont\color{popmuted}CE COURS SIGNÉ CONTIENT}\par\vspace{7pt}
  \noindent\begin{minipage}[t]{0.235\linewidth}\poptile{popblue}{Cours}{complet, illustré, pas à pas}\end{minipage}\hfill
  \begin{minipage}[t]{0.235\linewidth}\poptile{poporange}{Méthodes}{et exercices résolus}\end{minipage}\hfill
  \begin{minipage}[t]{0.235\linewidth}\poptile{poppink}{Exercices}{type examen + corrigés}\end{minipage}\hfill
  \begin{minipage}[t]{0.235\linewidth}\poptile{popgreen}{Fiche bilan}{l'essentiel à retenir}\end{minipage}\par}

% Sommaire
\newtcolorbox{sommaire}{enhanced,colback=white,colframe=popink,boxrule=2.2pt,arc=10pt,
  drop shadow=popink,left=18pt,right=18pt,top=13pt,bottom=13pt,before skip=6pt,after skip=20pt,
  before upper={\poppill{Sommaire}{poppurple}{white}\par\vspace{9pt}\popsemi\fontsize{10.6}{14.5}\selectfont\color{popink}}}

% Signature de fin de cours — poussée tout en bas de la dernière page
\newcommand{\findecours}{%
  \par\vfill\nopagebreak
  \begin{center}\popsquiggle{poporange}\end{center}\vspace{4pt}
  \signatureonbuch
}
\AtBeginDocument{\def\llbracket{\mathopen{[\mkern-3.5mu[}}\def\rrbracket{\mathclose{]\mkern-3.5mu]}}} % intervalles d'entiers (glyphe ⟦ absent de la police)
% Glyphes absents de Fira Math : redéfinis à partir de glyphes disponibles
\AtBeginDocument{%
  \def\setminus{\mathbin{\backslash}}\let\smallsetminus\setminus
  \def\vdots{\mathord{\vbox{\baselineskip3.2pt\lineskiplimit0pt\kern4pt\hbox{.}\hbox{.}\hbox{.}}}}%
  \def\ddots{\mathinner{\mkern1mu\raise6.5pt\hbox{.}\mkern2mu\raise3.6pt\hbox{.}\mkern2mu\raise0.7pt\hbox{.}\mkern1mu}}%
  \def\triangle{\mathord{\scalebox{0.9}{$\symup{\Delta}$}}}%
  \def\nabla{\mathord{\raisebox{\depth}{\scalebox{1}[-1]{$\symup{\Delta}$}}}}%
  \def\checkmark{\popcheck{popdark}}%
  \def\star{\mathbin{\ast}}\def\bigstar{\mathord{\ast}}%
  \def\diamond{\mathbin{\scalebox{0.6}{$\Diamond$}}}%
  \def\bigcup{\mathop{\mathchoice{\raisebox{-0.3em}{\scalebox{1.7}{$\cup$}}}{\scalebox{1.25}{$\cup$}}{\cup}{\cup}}\displaylimits}%
  \def\bigcap{\mathop{\mathchoice{\raisebox{-0.3em}{\scalebox{1.7}{$\cap$}}}{\scalebox{1.25}{$\cap$}}{\cap}{\cap}}\displaylimits}%
  \def\lesseqqgtr{\mathrel{\lessgtr}}\def\bigoplus{\mathop{\oplus}\displaylimits}%
}
\providecommand{\ssi}{si et seulement si} % abréviation fréquente
\AtBeginDocument{\providecommand{\wideparen}{\overparen}} % arc de cercle

%% FIGFIX-BEGIN v3 — correctifs globaux des figures (docs/onbuch-generation/tools/figfix)
\usepackage{adjustbox}\usepackage{etoolbox}
% toute figure TikZ/pgfplots plus large que la ligne est réduite à la largeur disponible
\BeforeBeginEnvironment{tikzpicture}{\begin{adjustbox}{max width=\linewidth}}
\AfterEndEnvironment{tikzpicture}{\end{adjustbox}}
% légendes pgfplots lisibles par-dessus les courbes
\pgfplotsset{every axis/.append style={legend style={fill=white,fill opacity=0.92,text opacity=1,draw=black!15,inner sep=3pt}}}
% étiquettes d'arêtes (syntaxe edge["..."]) avec fond blanc
\tikzset{every edge quotes/.append style={fill=white,inner sep=1.2pt}}
% bulles de cartes mentales : texte à la ligne dans la bulle, bulle un peu plus grande
\tikzset{every concept/.append style={text width=2.0cm,align=center,minimum size=2.3cm,inner sep=2pt}}
%% FIGFIX-END
\begin{document}

\pagegarde{Sensibilisation sur la nécessité du renouvellement de l'ATP lors des exercices musculaires}{SVT}{Tle D}{Module I : Le Monde Vivant — Métabolisme énergétique des cellules musculaires}{N02}{8 h}{Cette leçon étudie la structure de la fibre musculaire striée squelettique, le mécanisme de sa contraction et les voies métaboliques (anaérobie alactique, anaérobie lactique, aérobie) qui assurent le renouvellement permanent de l'ATP selon le type d'effort fourni. Elle s'appuie sur l'interprétation de résultats expérimentaux (biopsies, courbes physiologiques) pour établir l'adaptation du métabolisme musculaire à chaque type d'exercice.
\vspace{4pt}
\begin{multicols}{2}\small
\begin{itemize}
\item À la fin de cette leçon, je sais décrire la structure d'une fibre musculaire striée squelettique et citer les principaux composants chimiques du muscle
\item À la fin de cette leçon, je sais distinguer les principaux types de fibres musculaires et leurs aptitudes
\item À la fin de cette leçon, je sais illustrer par des schémas une portion de fibre musculaire au repos et en contraction
\item À la fin de cette leçon, je sais réaliser le schéma d'interprétation du mécanisme de la contraction musculaire
\item À la fin de cette leçon, je sais interpréter les variations de la composition chimique du muscle en fonction de l'activité
\item À la fin de cette leçon, je sais interpréter les données de biopsies relatives aux taux d'ATP et de phosphocréatine lors d'exercices de durées différentes
\end{itemize}\end{multicols}}

\begin{sommaire}
\begin{multicols}{2}
\begin{enumerate}
\item I. Structure et composition chimique de la fibre musculaire striée squelettique
\item II. Les principaux types de fibres musculaires
\item III. Le mécanisme de la contraction musculaire
\item IV. La nécessité du renouvellement de l'ATP : les voies de restauration
\item V. Adaptation du métabolisme au type d'effort
\item Activité d'intégration
\item Méthodes et exercices résolus
\item Exercices
\item Corrigés des exercices
\item Fiche bilan
\end{enumerate}
\end{multicols}
\end{sommaire}

\begin{objectifs}
\begin{itemize}
\item À la fin de cette leçon, je sais décrire la structure d'une fibre musculaire striée squelettique et citer les principaux composants chimiques du muscle
\item À la fin de cette leçon, je sais distinguer les principaux types de fibres musculaires et leurs aptitudes
\item À la fin de cette leçon, je sais illustrer par des schémas une portion de fibre musculaire au repos et en contraction
\item À la fin de cette leçon, je sais réaliser le schéma d'interprétation du mécanisme de la contraction musculaire
\item À la fin de cette leçon, je sais interpréter les variations de la composition chimique du muscle en fonction de l'activité
\item À la fin de cette leçon, je sais interpréter les données de biopsies relatives aux taux d'ATP et de phosphocréatine lors d'exercices de durées différentes
\item À la fin de cette leçon, je sais décrire les trois voies de restauration de l'ATP et les comparer (vitesse, rendement, durée)
\item À la fin de cette leçon, je sais analyser et interpréter des documents liés aux exigences métaboliques de différents sports
\item À la fin de cette leçon, je sais interpréter l'évolution des paramètres physiologiques (ATP, phosphocréatine, glycogène) au cours de l'exercice musculaire
\end{itemize}
\end{objectifs}

\begin{prerequis}
\begin{itemize}
\item Notion de métabolisme cellulaire et rôle de l'ATP comme molécule énergétique (Première)
\item Structure générale d'une cellule animale et rôle des mitochondries (Seconde/Première)
\item Notion d'enzyme et de réaction chimique catalysée (Première)
\item Lecture et exploitation de courbes et de tableaux de données expérimentales
\item Notions sur le système nerveux et le message nerveux moteur (début de Terminale)
\end{itemize}
\end{prerequis}

\newpage

\coursec{I. Structure et composition chimique de la fibre musculaire striée squelettique}

\textbf{Situation d'accroche.} Lors des grandes vacances, Junior, élève en Terminale D à Yaoundé, participe à un match de football inter-quartiers. Après un sprint de 30 mètres pour marquer, il se sent en pleine forme ; mais après 20 minutes de jeu intense, ses cuisses « brûlent » et il doit ralentir. Son grand frère, marathonien amateur, court pourtant pendant deux heures sans s'arrêter. D'où vient l'énergie qui fait contracter les muscles ? Pourquoi cette énergie semble-t-elle s'épuiser différemment selon l'effort ? Comment le muscle fabrique-t-il et renouvelle-t-il son ATP ?

\textbf{Problématique.} Pour répondre à ces questions, il faut d'abord connaître le matériel biologique concerné : de quoi est fait un muscle squelettique ? Quelle est l'organisation de ses cellules ? Quelles molécules y trouve-t-on et en quelles quantités ? C'est l'objet de cette première partie : décrire la \cle{structure} de la fibre musculaire striée squelettique et établir sa \cle{composition chimique}, avant d'étudier, dans les parties suivantes, comment cette machine contractile fonctionne et se réapprovisionne en énergie.

\courssub{De l'observation macroscopique à l'observation microscopique : l'organisation du muscle}

\textbf{Observation.} Prends un muscle squelettique, par exemple le biceps d'un bœuf acheté au marché (ou une cuisse de poulet). À l'œil nu, on distingue des faisceaux charnus entourés d'une gaine conjonctive. En disséquant délicatement, chaque faisceau se révèle constitué de filaments fins, parallèles : ce sont les \cle{fibres musculaires}. Plaçons maintenant un fragment de fibre entre lame et lamelle, colorons-le et observons-le au microscope optique : la fibre apparaît comme un long cylindre présentant des \emph{striations transversales} alternées claires et sombres — d'où le nom de muscle \emph{strié}. En microscopie électronique, on découvre que chaque fibre contient elle-même des centaines de filaments contractiles, les \cle{myofibrilles}, dont l'organisation régulière explique les stries.

\begin{experience}[Observation multi-échelle d'un muscle squelettique]
\textbf{Matériel :} fragment de muscle squelettique de bœuf ou de grenouille, aiguilles de dissection, lames et lamelles, colorant (bleu de méthylène ou hématoxyline), microscope optique, coupes histologiques au microscope électronique (documents).

\textbf{Protocole :}
\begin{enumerate}
\item À l'œil nu et à la loupe binoculaire : identifier les faisceaux de fibres et les effilocher avec les aiguilles.
\item Au microscope optique (grossissement $\times 400$) : observer une fibre isolée ; repérer les striations transversales et les noyaux situés en périphérie.
\item Sur document (microscopie électronique) : observer l'intérieur d'une fibre et identifier les myofibrilles et leurs bandes.
\end{enumerate}

\textbf{Observations :} la fibre est une cellule cylindrique, très allongée (jusqu'à plusieurs centimètres), de diamètre compris entre \num{10} et \SI{100}{\micro\metre}, contenant de nombreux noyaux repoussés à la périphérie. L'intérieur est rempli de myofibrilles alignées, chacune montrant une succession régulière de bandes claires et sombres parfaitement alignées d'une myofibrille à l'autre, ce qui produit la striation visible à l'échelle de la fibre entière.

\textbf{Interprétation :} le muscle est une \emph{hiérarchie de structures emboîtées} : muscle $\rightarrow$ faisceaux $\rightarrow$ fibres $\rightarrow$ myofibrilles $\rightarrow$ filaments protéiques. La régularité de cette organisation à tous les niveaux est la signature d'une machine contractile ordonnée.
\end{experience}

\begin{popfigure}
\begin{center}
\begin{tikzpicture}[scale=1, every node/.style={font=\small}]
% Muscle
\draw[fill=poporangeL, draw=popdark, rounded corners=2pt] (0,0) ellipse (1.6 and 0.55);
\node[below] at (0,-0.7) {\textbf{Muscle}};
% Faisceau
\draw[fill=popgoldL, draw=popdark, rounded corners=2pt] (4.2,0) ellipse (1.1 and 0.5);
\foreach \y in {-0.25,-0.08,0.09,0.26} \draw[popdark!60] (3.2,\y) -- (5.2,\y);
\node[below] at (4.2,-0.7) {\textbf{Faisceau de fibres}};
% Fibre
\draw[fill=popgreenL, draw=popdark, rounded corners=6pt] (7.6,-0.35) rectangle (10.4,0.35);
\foreach \x in {7.9,8.4,8.9,9.4,9.9} \draw[popdark!50] (\x,-0.35) -- (\x,0.35);
\fill[popdark] (7.75,0.42) circle (0.05); \fill[popdark] (9.9,0.42) circle (0.05);
\node[below] at (9,-0.7) {\textbf{Fibre musculaire (cellule)}};
\node[above, font=\scriptsize] at (9,0.5) {noyaux périphériques};
% Myofibrille
\draw[fill=popblueL, draw=popdark] (11.6,-0.18) rectangle (14.6,0.18);
\foreach \x in {12.2,12.8,13.4,14.0} \draw[popdark, thick] (\x,-0.18) -- (\x,0.18);
\node[below] at (13.1,-0.7) {\textbf{Myofibrille (organite)}};
\node[above, font=\scriptsize] at (13.1,0.28) {stries Z};
% flèches
\draw[-{Stealth}, thick, poporange] (1.7,0) -- (3.0,0);
\draw[-{Stealth}, thick, poporange] (5.4,0) -- (7.4,0);
\draw[-{Stealth}, thick, poporange] (10.5,0) -- (11.5,0);
\end{tikzpicture}
\end{center}
\legende{Organisation emboîtée du muscle squelettique, de l'échelle macroscopique à l'échelle subcellulaire : le muscle est formé de faisceaux de fibres ; chaque fibre (une cellule plurinucléée) contient des centaines de myofibrilles striées, elles-mêmes constituées de myofilaments protéiques.}
\end{popfigure}

\begin{savaistu}[Une cellule géante issue de fusions]
Une fibre musculaire peut mesurer plusieurs centimètres de long (jusqu'à \SI{30}{cm} dans les muscles de la cuisse !) et contenir des centaines de noyaux. Cette particularité étonnante s'explique par son origine embryonnaire : la fibre résulte de la \emph{fusion} de nombreuses cellules embryonnaires appelées \emph{myoblastes}. C'est donc un \emph{syncytium} : une seule membrane, un seul cytoplasme, mais de nombreux noyaux, ce qui permet de coordonner la synthèse des protéines sur toute la longueur de la cellule.
\end{savaistu}

\courssub{Structure fine de la fibre musculaire}

\begin{definition}[La fibre musculaire striée squelettique]
La \cle{fibre musculaire} striée squelettique est une \textbf{cellule} cylindrique, allongée, plurinucléée, entourée d'une membrane appelée \cle{sarcolemme} et dont le cytoplasme, le \cle{sarcoplasme}, contient des organites spécialisés dans la contraction (myofibrilles) et dans la production d'énergie (mitochondries).
\end{definition}

Détaillons les constituants de cette cellule remarquable :

\begin{itemize}
\item Le \cle{sarcolemme} : membrane plasmique de la fibre. Elle conduit l'influx nerveux (potentiel d'action) qui déclenche la contraction, et s'invagine vers l'intérieur de la cellule en formant les \emph{tubules T} (ou tubules transverses), qui transmettent le signal au cœur de la fibre au plus près des myofibrilles.
\item Le \cle{sarcoplasme} : cytoplasme de la fibre, riche en glycogène (réserve glucidique), en enzymes et en \cle{myoglobine}, protéine rouge qui fixe et stocke le dioxygène (elle est apparentée à l'hémoglobine).
\item Les \cle{noyaux périphériques} : nombreux, plaqués contre le sarcolemme, ils dirigent la synthèse des protéines de la cellule géante.
\item Le \cle{réticulum sarcoplasmique} : réseau de saccules membranaires entourant chaque myofibrille. C'est le \textbf{réservoir de calcium} de la fibre : il stocke les ions \ce{Ca^{2+}} au repos et les libère massivement lors de l'excitation. Nous verrons en partie III que le calcium est la clé qui déclenche la contraction.
\item Les \cle{mitochondries} : très nombreuses, serrées entre les myofibrilles, elles réalisent la respiration cellulaire qui produit l'ATP en présence de dioxygène. Leur abondance témoigne des besoins énergétiques considérables du muscle.
\item Les \cle{myofibrilles} : organites contractiles, au nombre de plusieurs centaines à plus d'un millier par fibre, disposés parallèlement et occupant l'essentiel du volume cellulaire.
\end{itemize}

\begin{attention}[Fibre ou myofibrille : ne pas confondre !]
Erreur très fréquente au Baccalauréat : la \textbf{fibre musculaire} est une \emph{cellule} (elle possède une membrane, des noyaux, des mitochondries), tandis que la \textbf{myofibrille} est un \emph{organite contractile intracellulaire}, contenu dans le sarcoplasme de la fibre. Une fibre contient des centaines de myofibrilles. De même, ne confonds pas \emph{myofibrille} et \emph{myofilament} : le myofilament (d'actine ou de myosine) est une molécule protéique assemblée en filament, contenu dans la myofibrille. Retiens la hiérarchie : fibre $\supset$ myofibrille $\supset$ myofilaments.
\end{attention}

\courssub{La myofibrille et le sarcomère : l'unité contractile}

\textbf{Observation.} Au microscope électronique, chaque myofibrille montre une succession périodique de \emph{bandes sombres} (bandes A) et de \emph{bandes claires} (bandes I). Chaque bande claire est traversée en son milieu par un trait sombre très fin : la \cle{strie Z}. Cette répétition rigoureuse suggère l'existence d'une unité élémentaire répétée des milliers de fois le long de la myofibrille.

\begin{definition}[Le sarcomère]
Le \cle{sarcomère} est l'\textbf{unité structurale et fonctionnelle} de la myofibrille. Il est compris entre \textbf{deux stries Z consécutives} et mesure environ \SI{2.5}{\micro\metre} au repos. Il comprend une bande A (sombre) centrale, encadrée de deux demi-bandes I (claires).
\end{definition}

L'examen à fort grossissement révèle que ces bandes correspondent à deux types de filaments protéiques, les \cle{myofilaments}, partiellement imbriqués :

\begin{itemize}
\item Les \textbf{filaments épais}, constitués de \cle{myosine}, occupent le centre du sarcomère (bande A). Chaque molécule de myosine possède une \emph{tête} globulaire qui fait saillie latéralement : ces \emph{têtes de myosine} pourront s'accrocher à l'actine et hydrolyser l'ATP.
\item Les \textbf{filaments fins}, constitués d'\cle{actine}, sont ancrés sur les stries Z et s'étendent vers le centre du sarcomère, où ils s'intercalent entre les filaments de myosine. Ils sont associés à deux protéines régulatrices : la \cle{tropomyosine}, qui recouvre les sites de fixation de la myosine au repos, et la \cle{troponine}, capable de fixer les ions \ce{Ca^{2+}}.
\end{itemize}

L'aspect des bandes s'explique alors simplement : la \textbf{bande A} (sombre) correspond à la zone occupée par les filaments épais de myosine (avec ou sans recouvrement par l'actine) ; la \textbf{bande I} (claire) correspond aux portions des filaments fins seuls, de part et d'autre de la strie Z ; au centre de la bande A, la \textbf{zone H}, plus claire, correspond à la région où la myosine n'est pas recouverte par l'actine.

\begin{popfigure}
\begin{center}
\begin{tikzpicture}[scale=1.05, every node/.style={font=\small}]
% Stries Z
\draw[very thick, popdark] (0,-1.1) -- (0,1.1);
\draw[very thick, popdark] (10,-1.1) -- (10,1.1);
\node[above] at (0,1.15) {strie Z};
\node[above] at (10,1.15) {strie Z};
% Myosine (filaments épais) avec têtes
\foreach \y in {-0.55,0,0.55}{
  \draw[line width=2.2pt, poporange] (3.2,\y) -- (6.8,\y);
  \foreach \x in {3.5,4.1,4.7,5.3,5.9,6.5}{
    \draw[poporange, thick] (\x,\y) -- (\x,\y+0.28);
    \draw[poporange, thick] (\x,\y) -- (\x,\y-0.28);
  }
}
% Actine (filaments fins)
\foreach \y in {-0.85,-0.28,0.28,0.85}{
  \draw[line width=1.1pt, popblue] (0,\y) -- (4.4,\y);
  \draw[line width=1.1pt, popblue] (5.6,\y) -- (10,\y);
}
% Légendes de bandes
\draw[<->, popmuted] (0,-1.5) -- (3.2,-1.5); \node[below, font=\scriptsize] at (1.6,-1.5) {bande I};
\draw[<->, popmuted] (3.2,-1.5) -- (6.8,-1.5); \node[below, font=\scriptsize] at (5,-1.5) {bande A};
\draw[<->, popmuted] (6.8,-1.5) -- (10,-1.5); \node[below, font=\scriptsize] at (8.4,-1.5) {bande I};
\draw[<->, popmuted] (4.4,1.5) -- (5.6,1.5); \node[above, font=\scriptsize] at (5,1.5) {zone H};
% Accolade sarcomère
\draw[decorate, decoration={brace, amplitude=5pt}, popgreen] (0,-2.1) -- (10,-2.1);
\node[below, popgreen, font=\small\bfseries] at (5,-2.35) {Sarcomère ($\approx$ \SI{2.5}{\micro\metre} au repos)};
% Légendes filaments
\node[poporange, font=\scriptsize, anchor=west] at (10.3,0.55) {filament épais de myosine (têtes latérales)};
\node[popblue, font=\scriptsize, anchor=west] at (10.3,-0.55) {filament fin d'actine (+ tropomyosine, troponine)};
\end{tikzpicture}
\end{center}
\legende{Schéma d'un sarcomère au repos, délimité par deux stries Z. La bande A (sombre) correspond aux filaments épais de myosine ; les bandes I (claires) aux filaments fins d'actine seuls ; la zone H, au centre de la bande A, est la région de myosine non recouverte par l'actine. Les têtes de myosine font saillie vers les filaments d'actine.}
\end{popfigure}

\begin{aretenir}[L'essentiel de la structure]
\begin{itemize}
\item Muscle $\rightarrow$ faisceaux $\rightarrow$ \textbf{fibres} (cellules plurinucléées) $\rightarrow$ \textbf{myofibrilles} $\rightarrow$ \textbf{myofilaments} (actine et myosine).
\item Le \textbf{sarcomère}, entre deux stries Z, est l'unité contractile : bande A centrale (myosine), demi-bandes I latérales (actine seule), zone H centrale.
\item Le \textbf{réticulum sarcoplasmique} stocke le \ce{Ca^{2+}} ; les \textbf{mitochondries} produisent l'ATP ; la \textbf{myoglobine} stocke le \ce{O2}.
\end{itemize}
\end{aretenir}

\courssub{Composition chimique de la fibre musculaire : mise en évidence expérimentale}

\textbf{Question.} De quoi est chimiquement constitué un muscle au repos ? Pour le savoir, les physiologistes réalisent des \emph{biopsies musculaires} : on prélève, sous anesthésie locale, un fragment de quelques milligrammes de muscle (par exemple le vaste externe de la cuisse) à l'aide d'une aiguille spéciale, puis on dose ses constituants par des techniques biochimiques.

\begin{experience}[Dosage des constituants d'un muscle au repos]
\textbf{Protocole :} sur un échantillon de \SI{1}{kg} de muscle frais prélevé au repos, on réalise : une dessiccation (mesure de l'eau), une extraction et un dosage des protéines, une hydrolyse enzymatique du glycogène suivie du dosage du glucose, et des dosages spectrophotométriques de l'ATP et de la phosphocréatine.

\textbf{Résultats (valeurs moyennes pour \SI{1}{kg} de muscle frais au repos) :}

\begin{center}
\begin{tabularx}{\linewidth}{|l|X|X|}\hline
\rowcolor{popblueL} \textbf{Constituant} & \textbf{Proportion ou teneur} & \textbf{Rôle} \\ \hline
Eau & environ 75 \% de la masse & milieu des réactions biochimiques \\ \hline
Protéines & environ 20 \% de la masse & protéines contractiles (actine, myosine), enzymes, myoglobine \\ \hline
Glycogène & environ 1 \% (\SIrange{10}{20}{g.kg^{-1}}) & réserve glucidique \\ \hline
Lipides, sels minéraux & quelques pourcents & membranes, équilibre ionique (\ce{Ca^{2+}}, \ce{K+}, \ce{Na+}) \\ \hline
ATP & environ \SI{5}{mmol.kg^{-1}} & source d'énergie immédiate \\ \hline
Phosphocréatine (PC) & environ \SI{20}{mmol.kg^{-1}} & réserve énergétique rapide \\ \hline
\end{tabularx}
\end{center}

\textbf{Interprétation :} le muscle est essentiellement fait d'eau et de protéines. Les protéines contractiles (actine et myosine) représentent la machine ; le glycogène, l'ATP et la phosphocréatine constituent le « carburant » et les réserves d'énergie ; la myoglobine stocke le comburant (\ce{O2}).
\end{experience}

\begin{definition}[La phosphocréatine]
La \cle{phosphocréatine} (PC) est une molécule de \textbf{réserve énergétique} propre au muscle. Elle possède une liaison phosphate riche en énergie et peut céder \emph{très rapidement} son groupement phosphate à l'ADP pour régénérer l'ATP selon la réaction réversible :
\[ \text{phosphocréatine} + \text{ADP} \rightleftharpoons \text{créatine} + \text{ATP} \]
Cette réaction, catalysée par l'enzyme \emph{créatine kinase}, est la voie la plus rapide de régénération de l'ATP (voie anaérobie alactique, étudiée en partie IV).
\end{definition}

\begin{exemplebox}[Des réserves d'ATP étonnamment faibles]
Le muscle au repos contient environ \SI{5}{mmol} d'ATP et \SI{20}{mmol} de phosphocréatine par kilogramme de muscle frais. Or un effort intense consomme plusieurs mmol d'ATP par kilogramme et par seconde ! Le stock d'ATP ne permet donc que \textbf{quelques secondes} d'effort maximal. C'est précisément pourquoi le muscle doit \emph{renouveler en permanence} son ATP : c'est toute la problématique de cette leçon, annoncée par la situation de Junior.
\end{exemplebox}

\begin{exoresolu}[Exploitation d'un tableau de composition]
\textbf{Énoncé.} Une biopsie du muscle vaste externe d'un sujet au repos donne les résultats suivants (par kg de muscle frais) : eau : \SI{750}{g} ; protéines : \SI{200}{g} ; glycogène : \SI{15}{g} ; ATP : \SI{5}{mmol} ; phosphocréatine : \SI{20}{mmol} ; lipides et sels minéraux : le reste.
\begin{enumerate}
\item Calculer le pourcentage massique des protéines dans ce muscle.
\item Comparer les quantités d'ATP et de phosphocréatine et proposer une interprétation.
\item Sachant qu'un sprint maximal consomme environ \SI{2}{mmol} d'ATP par kg de muscle et par seconde, estimer la durée pendant laquelle le seul stock d'ATP pourrait alimenter l'effort.
\end{enumerate}
\tcblower
\textbf{Solution détaillée.}
\begin{enumerate}
\item Pourcentage de protéines : $\dfrac{200}{1000} \times 100 = 20\,\%$. Les protéines représentent environ un cinquième de la masse du muscle.
\item La phosphocréatine est quatre fois plus abondante que l'ATP ($20 = 4 \times 5$ mmol). Interprétation : l'ATP n'est pas stocké en grande quantité ; la PC constitue une réserve tampon permettant de régénérer rapidement l'ATP consommé au début de l'effort.
\item Durée théorique d'utilisation du stock d'ATP : $t = \dfrac{5\ \text{mmol}}{2\ \text{mmol.s}^{-1}} = 2,5\ \text{s}$. Le stock d'ATP ne permettrait qu'environ \SI{2.5}{s} de sprint maximal ! Conclusion : le muscle doit impérativement \emph{resynthétiser} son ATP pendant l'effort, ce qui justifie l'existence des voies de restauration étudiées dans la suite de la leçon.
\end{enumerate}
\end{exoresolu}

\begin{attention}[Pièges de rédaction au Baccalauréat]
\begin{itemize}
\item Ne dis pas que le muscle « contient de l'énergie » : il contient des \emph{molécules} (ATP, PC, glycogène) dont l'hydrolyse ou la dégradation \emph{libère} de l'énergie utilisable.
\item L'unité des dosages est la mmol par kg de muscle \emph{frais} : précise toujours l'unité et le référentiel.
\item La myoglobine n'est pas de l'hémoglobine : elle stocke le \ce{O2} \emph{dans le muscle}, alors que l'hémoglobine le \emph{transporte} dans le sang.
\item Attention à l'écriture des unités : on écrit \SI{5}{mmol.kg^{-1}} (ou \SI{5}{mmol/kg}) et non « mmol/kg de muscle » dans une formule ; en texte, on précise « par kilogramme de muscle frais ».
\end{itemize}
\end{attention}

\begin{exercice}[Pour t'entraîner]
On prélève par biopsie un fragment de muscle de \SI{50}{mg} chez un athlète au repos. On y dose \SI{0.25}{\micro mol} d'ATP.
\begin{enumerate}
\item Vérifier que ce résultat est conforme à la valeur moyenne de \SI{5}{mmol.kg^{-1}}.
\item Citer deux autres constituants que l'on pourrait doser dans cet échantillon et préciser leur rôle.
\item Expliquer pourquoi les mitochondries sont particulièrement nombreuses dans les fibres musculaires.
\end{enumerate}
\end{exercice}

\begin{corrige}[Exercice]
\begin{enumerate}
\item $\SI{50}{mg} = \SI{5e-5}{kg}$. Quantité attendue : $5 \times \num{5e-5} = \num{2.5e-4}$ mmol $= \SI{0.25}{\micro mol}$. Le résultat est conforme.
\item Le glycogène (réserve glucidique pour la production d'ATP) et la phosphocréatine (réserve énergétique rapide régénérant l'ATP) ; on peut aussi citer la myoglobine (réserve de \ce{O2}).
\item La contraction musculaire consomme beaucoup d'ATP. Les mitochondries réalisent la respiration cellulaire ($\ce{C6H12O6 + 6O2 -> 6CO2 + 6H2O}$ avec production d'ATP) : leur abondance permet de couvrir les besoins énergétiques élevés et continus de la fibre.
\end{enumerate}
\end{corrige}

\textbf{Transition.} Nous savons désormais de quoi est faite la fibre musculaire : une cellule géante organisée en sarcomères, riche en protéines contractiles et dotée de réserves énergétiques faibles mais renouvelables. Mais toutes les fibres se ressemblent-elles ? Le sprint de Junior et le marathon de son frère ne sollicitent-ils pas des fibres différentes ? C'est ce que nous découvrirons dans la partie suivante : les principaux types de fibres musculaires.

\coursec{II. Les principaux types de fibres musculaires}

Dans la section précédente, nous avons décrit la fibre musculaire striée squelettique comme une cellule géante, plurinucléée, remplie de myofibrilles et riche en protéines contractiles (actine, myosine), en glycogène et en enzymes. Cette description, cependant, laisserait croire que toutes les fibres musculaires se ressemblent. Or, une simple observation de boucherie contredit cette idée : la viande de poulet est pâle, celle du canard est foncée. De même, au stade omnisports de Yaoundé, le sprinteur qui explose sur 100 m et le maratonien qui court 42 km n'ont visiblement pas les mêmes « moteurs » musculaires. D'où vient cette différence ? Existe-t-il plusieurs \emph{types} de fibres musculaires ? C'est la question que nous allons traiter ici, en suivant la démarche expérimentale qui a permis aux histologistes de répondre.

\courssub{II.1. Mise en évidence expérimentale : une mosaïque de fibres}

\textbf{Observation et problème.} À l'œil nu, certains muscles paraissent rouge sombre (muscles des membres du canard, muscles du maintien de la posture chez l'Homme comme le soléaire du mollet), d'autres pâles (muscles pectoraux du poulet). \emph{Problème :} cette différence de couleur traduit-elle une différence de nature des fibres musculaires ?

\begin{experience}[Colorations histochimiques des coupes de muscle]
\textbf{Matériel :} muscle squelettique de mammifère (par exemple le gastrocnémien de rat), congélateur, microtome, microscopes optiques, révélateurs enzymatiques.

\textbf{Protocole :}
\begin{enumerate}
\item On réalise des \cle{coupes transversales} fines du muscle.
\item Sur une série de coupes, on révèle l'activité des \cle{ATPases myosiniques} (enzymes des ponts de myosine qui hydrolysent l'ATP) : plus l'activité ATPasique est forte, plus la fibre se colore intensément (coloration à pH basique usuelle).
\item Sur une série voisine, on révèle les \cle{enzymes oxydatives} (par exemple la succino-déshydrogénase, enzyme mitochondriale de la respiration cellulaire) : la coloration est d'autant plus foncée que la fibre est riche en mitochondries.
\end{enumerate}

\textbf{Observations :} la coupe apparaît comme une \emph{mosaïque} : côte à côte, on distingue des fibres foncées et des fibres claires. De plus, la comparaison des deux séries montre que les fibres \emph{faiblement} colorées par les ATPases sont \emph{fortement} colorées par les enzymes oxydatives, et inversement.

\textbf{Interprétation :} il existe au sein d'un même muscle des fibres de nature différente :
\begin{itemize}
\item des fibres à \emph{faible} activité ATPasique (contraction lente) mais \emph{riches} en mitochondries (métabolisme aérobie) ;
\item des fibres à \emph{forte} activité ATPasique (contraction rapide) mais \emph{pauvres} en mitochondries (métabolisme essentiellement anaérobie).
\end{itemize}

\textbf{Conclusion :} un muscle squelettique n'est pas homogène : c'est un assemblage de fibres de \emph{types différents}, en proportions variables.
\end{experience}

\begin{popfigure}
\begin{center}
\begin{tikzpicture}[scale=0.9]
% Fibre rouge (Type I)
\draw[fill=poppink!30,draw=poppink!80!black,thick] (0,0) ellipse (1.3 and 1.1);
% Mitochondries (nombreuses, périphériques et centrales)
\foreach \a in {0,60,120,180,240,300}{
  \draw[fill=popblue!60,draw=popblue!80!black] ({0.55*cos(\a)},{0.5*sin(\a)}) ellipse (0.16 and 0.10);}
\draw[fill=popblue!60,draw=popblue!80!black] (0,0) ellipse (0.16 and 0.10);
% Myoglobine (petits points rouges/orangés diffus)
\foreach \a in {30,90,150,210,270,330}{
  \fill[popgold!90!black] ({0.9*cos(\a)},{0.8*sin(\a)}) circle (0.06);}
\fill[popgold!90!black] (0,0) circle (0.06);
\node[below] at (0,-1.35) {\textbf{Fibre rouge (type I)}};
\node[below,align=center,font=\small] at (0,-1.8) {diamètre plus faible\\ nombreuses mitochondries\\ riche en myoglobine};
% Fibre blanche (Type II)
\draw[fill=popcream,draw=popmuted,thick] (5.5,0) ellipse (1.7 and 1.4);
% Mitochondries (rares)
\draw[fill=popblue!60,draw=popblue!80!black] (5.0,0.3) ellipse (0.16 and 0.10);
\draw[fill=popblue!60,draw=popblue!80!black] (6.1,-0.4) ellipse (0.16 and 0.10);
% Glycogène (nombreux granules verts)
\foreach \a in {0,45,90,135,180,225,270,315}{
  \fill[popgreen!70!black] ({5.5+1.2*cos(\a)},{1.0*sin(\a)}) circle (0.05);}
\foreach \a in {20,110,200,290}{
  \fill[popgreen!70!black] ({5.5+0.6*cos(\a)},{0.5*sin(\a)}) circle (0.05);}
\node[below] at (5.5,-1.65) {\textbf{Fibre blanche (type II)}};
\node[below,align=center,font=\small] at (5.5,-2.1) {diamètre plus grand\\ rares mitochondries\\ riche en glycogène};
% Légende des symboles
\draw[fill=popblue!60,draw=popblue!80!black] (-2.2,2.2) ellipse (0.16 and 0.10);
\node[right,font=\small] at (-1.95,2.2) {mitochondrie};
\fill[popgold!90!black] (-2.2,1.75) circle (0.06);
\node[right,font=\small] at (-1.95,1.75) {myoglobine (fixe \ce{O2})};
\fill[popgreen!70!black] (-2.2,1.3) circle (0.05);
\node[right,font=\small] at (-1.95,1.3) {granules de glycogène};
\end{tikzpicture}
\end{center}
\legende{Comparaison schématique d'une fibre rouge (type I) et d'une fibre blanche (type II) vues en coupe transversale. La fibre rouge, plus fine, est riche en mitochondries et en myoglobine ; la fibre blanche, plus grosse, est riche en glycogène et pauvre en mitochondries.}
\end{popfigure}

\courssub{II.2. Les fibres rouges (type I) : des fibres lentes et endurantes}

\begin{definition}[Fibres rouges ou fibres de type I]
Les \cle{fibres rouges} (ou fibres de \cle{type I}, encore appelées fibres lentes oxydatives) sont des fibres musculaires de diamètre relativement faible, riches en \cle{myoglobine} et en \cle{mitochondries}, abondamment vascularisées, dont le métabolisme est essentiellement \cle{aérobie}. Leur contraction est \emph{lente} mais elles sont très \emph{résistantes à la fatigue}.
\end{definition}

\begin{definition}[La myoglobine]
La \cle{myoglobine} est une \emph{protéine musculaire} capable de \emph{fixer réversiblement le dioxygène} (comme l'hémoglobine du sang, mais avec une affinité plus forte). Elle constitue une \emph{réserve de \ce{O2}} au sein de la fibre et facilite la diffusion de l'oxygène du sang vers les mitochondries. C'est elle, pigment rouge, qui donne leur \emph{couleur rouge} aux fibres oxydatives.
\end{definition}

\textbf{Pourquoi ces fibres sont-elles endurantes ?} Le raisonnement enchaîne logiquement les faits :
\begin{enumerate}
\item richesse en mitochondries et en myoglobine, vascularisation dense $\Rightarrow$ la fibre dispose en permanence de \ce{O2} ;
\item présence de \ce{O2} $\Rightarrow$ la fibre dégrade complètement le glucose (et les lipides) par la \emph{respiration cellulaire} : \ce{C6H12O6 + 6O2 -> 6CO2 + 6H2O}, bilan très rentable en ATP (environ 30 à 32 ATP par glucose) ;
\item l'ATP est produite \emph{lentement} mais \emph{longtemps}, sans accumulation de déchet acide $\Rightarrow$ contraction soutenue sans fatigue précoce.
\end{enumerate}
En contrepartie, l'activité ATPasique des ponts de myosine y est faible : le cycle d'attachement-détachement des ponts est lent, donc la contraction est lente et peu puissante. Les fibres rouges sont les fibres de l'\emph{endurance} et du \emph{maintien postural}.

\courssub{II.3. Les fibres blanches (type II) : des fibres rapides et puissantes}

\begin{definition}[Fibres blanches ou fibres de type II]
Les \cle{fibres blanches} (ou fibres de \cle{type II}, encore appelées fibres rapides glycolytiques) sont des fibres musculaires de \emph{grand diamètre}, \emph{pauvres en myoglobine} et en mitochondries, mais \emph{riches en glycogène} et en enzymes de la glycolyse. Leur métabolisme est essentiellement \cle{anaérobie}. Leur contraction est \emph{rapide et puissante} mais elles sont rapidement \emph{fatigables}.
\end{definition}

\textbf{Pourquoi sont-elles rapides mais fatigables ?}
\begin{enumerate}
\item forte activité ATPasique myosinique et grand diamètre (nombreuses myofibrilles en parallèle) $\Rightarrow$ contraction \emph{rapide} et développant une \emph{grande force} ;
\item pauvreté en mitochondries et en myoglobine $\Rightarrow$ l'ATP est régénérée surtout par la \emph{glycolyse anaérobie} (fermentation lactique) : \ce{C6H12O6 -> 2 C3H6O3} (lactate + H+), bilan peu rentable (2 ATP par glucose, soit environ 15 à 18 fois moins que la respiration cellulaire complète) ;
\item accumulation de lactate et d'ions H+ (acidose) et épuisement rapide du glycogène $\Rightarrow$ \emph{fatigue} rapide.
\end{enumerate}
Les fibres blanches sont les fibres de l'\emph{explosivité} : sprint, saut, lancer, haltérophilie.

\begin{attention}[Un muscle n'est jamais « tout rouge » ou « tout blanc »]
Ne jamais écrire qu'un muscle est « rouge » ou « blanc » dans son ensemble. \emph{Tout} muscle squelettique contient un \emph{mélange} de fibres de type I et de type II, en \emph{proportions variables} selon le muscle et selon l'individu. On dit seulement qu'un muscle est « à prédominance de fibres rouges » (comme le soléaire) ou « à prédominance de fibres blanches » (comme les pectoraux du poulet). De même, il existe des fibres intermédiaires (type IIa, oxydo-glycolytiques) : la classification en deux types est une simplification exigible au Baccalauréat.
\end{attention}

\begin{methode}[Identifier un type de fibre sur une coupe histologique]
Face à une coupe de muscle colorée ou à un schéma de fibre, procéder par étapes :
\begin{enumerate}
\item Observer la \textbf{couleur} (ou l'intensité de la coloration des enzymes oxydatives) : foncée/rouge $\Rightarrow$ type I ; claire/blanche $\Rightarrow$ type II.
\item Mesurer ou comparer le \textbf{diamètre} de la fibre : fin $\Rightarrow$ type I ; gros $\Rightarrow$ type II.
\item Évaluer l'abondance des \textbf{mitochondries} (et, si le document le montre, de la myoglobine et des capillaires sanguins) : nombreuses $\Rightarrow$ type I ; rares $\Rightarrow$ type II.
\item Conclure en reliant le type de fibre à ses propriétés fonctionnelles : type I = lente, aérobie, endurante ; type II = rapide, anaérobie, puissante mais fatigable.
\end{enumerate}
\end{methode}

\courssub{II.4. Interprétation de biopsies : le sprinteur et le marathonien}

La \cle{biopsie musculaire} consiste à prélever, à l'aide d'une aiguille spéciale, un fragment de quelques milligrammes de muscle (généralement le vaste externe de la cuisse) chez un sujet vivant, puis à analyser histochimiquement les types de fibres. Les résultats obtenus chez des athlètes de disciplines différentes sont éloquents.

\begin{popfigure}
\begin{center}
\begin{tikzpicture}
\begin{axis}[
  ybar, bar width=0.55cm,
  width=0.85\linewidth, height=6.2cm,
  ymin=0, ymax=100,
  ylabel={Proportion de fibres (\%)},
  symbolic x coords={Sprinteur, Sédentaire, Marathonien},
  xtick=data,
  legend style={at={(0.5,-0.2)},anchor=north,legend columns=2},
  nodes near coords, every node near coord/.append style={font=\small},
  ymajorgrids=true, grid style={popline!40},
]
\addplot[fill=popblue!70,draw=popblue!80!black] coordinates {(Sprinteur,25) (Sédentaire,50) (Marathonien,80)};
\addplot[fill=poporange!75,draw=poporange!90!black] coordinates {(Sprinteur,75) (Sédentaire,50) (Marathonien,20)};
\legend{Fibres rouges (type I), Fibres blanches (type II)}
\end{axis}
\end{tikzpicture}
\end{center}
\legende{Proportions moyennes de fibres de type I et de type II dans le vaste externe (muscle de la cuisse) chez un sprinteur, un sujet sédentaire et un marathonien (valeurs indicatives issues de biopsies musculaires).}
\end{popfigure}

\textbf{Interprétation du document.} Le diagramme montre que :
\begin{itemize}
\item le \emph{sprinteur} possède une forte majorité de fibres blanches (environ \SI{75}{\%} de type II) : sa musculature est adaptée aux efforts brefs, violents et explosifs ;
\item le \emph{marathonien} possède une forte majorité de fibres rouges (environ \SI{80}{\%} de type I) : sa musculature est adaptée aux efforts prolongés d'intensité modérée ;
\item le \emph{sédentaire} présente des proportions intermédiaires, voisines de \num{50}/\num{50}.
\end{itemize}

\textbf{Conclusion :} il existe une \emph{corrélation} entre la composition en types de fibres des muscles et l'aptitude sportive : chacun est, pour ainsi dire, « taillé » musculairement pour un type d'effort.

\begin{aretenir}[Déterminisme des proportions de fibres]
Les proportions de fibres rouges et blanches sont \emph{largement déterminées génétiquement} (rôle de l'\cle{hérédité}) : on naît plutôt sprinteur ou plutôt endurant. L'\cle{entraînement} ne transforme pas massivement un type de fibre en un autre ; il améliore surtout les \emph{performances} des fibres existantes (hypertrophie, enrichissement en mitochondries et en capillaires pour l'endurance, en glycogène et en enzymes glycolytiques pour la force). C'est pourquoi l'entraînement ne peut pas faire d'un marathonien de naissance un champion du 100 m, et inversement.
\end{aretenir}

\begin{savaistu}[Le poulet et le canard : deux stratégies musculaires]
Les muscles pectoraux du \emph{poulet}, oiseau qui ne vole que par courtes envolées paniquées, sont majoritairement constitués de fibres \emph{blanches} : d'où sa chair pâle. Chez le \emph{canard migrateur}, capable de voler des centaines de kilomètres sans repos, ces mêmes muscles sont majoritairement \emph{rouges}, gorgés de myoglobine et de mitochondries : d'où sa chair foncée. La nature illustre exactement le principe que nous venons d'établir chez l'athlète : le type de fibres est adapté au type d'effort habituel de l'animal. On retrouve la même logique chez les poissons : le thon, nageur infatigable, a une chair rouge, alors que le poisson de fond, qui ne bouge que par à-coups, a une chair blanche.
\end{savaistu}

\begin{exoresolu}[Exploitation d'une biopsie musculaire]
\textbf{Énoncé.} On réalise des biopsies du muscle vaste externe chez deux athlètes A et B. Les colorations histochimiques donnent les résultats suivants : chez A, \SI{78}{\%} des fibres sont fortement colorées par le révélateur des ATPases myosiniques et faiblement colorées par le révélateur des enzymes oxydatives ; chez B, \SI{82}{\%} des fibres sont faiblement colorées par les ATPases et fortement colorées par les enzymes oxydatives.
\begin{enumerate}
\item Identifier le type de fibres dominant chez chaque athlète, en justifiant.
\item Indiquer, en le justifiant, lequel des deux athlètes est le plus probablement un spécialiste du 100 m.
\end{enumerate}
\tcblower
\textbf{Solution détaillée.}
\begin{enumerate}
\item Chez \textbf{A}, la forte activité ATPasique (hydrolyse rapide de l'ATP par la myosine) et la faible teneur en enzymes oxydatives (donc en mitochondries) caractérisent les \textbf{fibres blanches de type II} : contraction rapide, métabolisme anaérobie. Chez \textbf{B}, la faible activité ATPasique et l'abondance des enzymes oxydatives (mitochondries nombreuses) caractérisent les \textbf{fibres rouges de type I} : contraction lente, métabolisme aérobie.
\item Le 100 m est un effort \emph{bref, violent et explosif}, qui exige une contraction rapide et puissante : il sollicite les fibres blanches. C'est donc l'athlète \textbf{A} (\SI{78}{\%} de fibres de type II) qui est le plus probablement le sprinteur. L'athlète B, riche en fibres rouges endurantes, serait plutôt adapté aux courses de longue durée (marathon).
\end{enumerate}
\end{exoresolu}

\begin{attention}[Pièges de rédaction au Baccalauréat]
\begin{itemize}
\item Ne pas confondre \emph{myoglobine} (protéine musculaire fixant \ce{O2}, dans les fibres) et \emph{hémoglobine} (protéine sanguine transportant \ce{O2}, dans les globules rouges).
\item Ne pas écrire que les fibres blanches « n'ont pas de mitochondries » : elles en ont, mais \emph{peu}. De même, les fibres rouges contiennent du glycogène, mais moins que les blanches.
\item Ne pas affirmer que l'entraînement « transforme » les fibres blanches en fibres rouges : les proportions sont essentiellement \emph{héréditaires} ; l'entraînement ne fait qu'optimiser les fibres présentes.
\item Toujours relier la \emph{structure} (mitochondries, myoglobine, glycogène) à la \emph{fonction} (type de métabolisme, vitesse de contraction, résistance à la fatigue) : c'est cette relation structure-fonction que le correcteur attend.
\end{itemize}
\end{attention}

En résumé, les muscles squelettiques sont des mosaïques de fibres rouges (lentes, oxydatives, endurantes) et de fibres blanches (rapides, glycolytiques, puissantes mais fatigables), dont les proportions, largement héréditaires, conditionnent les aptitudes sportives. Reste à comprendre \emph{comment} ces fibres se contractent : c'est l'objet de la section suivante, consacrée au mécanisme de la contraction musculaire.

\coursec{III. Le mécanisme de la contraction musculaire}

Dans la section précédente, nous avons vu que les fibres musculaires striées squelettiques ne sont pas toutes identiques : fibres rapides glycolytiques, fibres lentes oxydatives, fibres intermédiaires. Mais qu'elle soit rapide ou lente, toute fibre musculaire possède une propriété fondamentale : elle peut se \cle{raccourcir} activement et développer une force. Comment un ensemble de protéines peut-il produire un mouvement ? C'est toute la question du \cle{mécanisme de la contraction musculaire}, que nous allons élucider en suivant la démarche expérimentale qui a conduit à la \emph{théorie du glissement des filaments}.

\courssub{Les faits expérimentaux : de la stimulation au raccourcissement}

\textbf{Observation de départ.} Au laboratoire, on prélève un muscle de grenouille (classiquement le muscle gastrocnémien, l'équivalent du mollet) maintenu en survie dans du liquide de Ringer, une solution physiologique qui préserve la viabilité des tissus. Ce muscle est fixé par une extrémité à un support et par l'autre à un myographe, appareil qui enregistre les variations de longueur.

\begin{experience}[Stimulation électrique d'un muscle isolé de grenouille]
\textbf{Matériel :} muscle gastrocnémien de grenouille, liquide de Ringer, électrodes reliées à un stimulateur, myographe enregistreur.

\textbf{Protocole :} on applique au muscle isolé un choc électrique bref et on enregistre sa longueur au cours du temps. On observe ensuite le muscle au microscope optique avant et pendant la contraction.

\textbf{Observations :}
\begin{itemize}
\item à chaque stimulation électrique efficace, le muscle se \textbf{raccourcit} et développe une tension (secousse musculaire) ; puis il se relâche et retrouve sa longueur initiale ;
\item l'observation microscopique montre que ce raccourcissement macroscopique s'accompagne du \textbf{raccourcissement de chaque sarcomère} : les lignes Z se rapprochent ;
\item l'examen fin des stries révèle que les \textbf{bandes I} (claires) et la \textbf{zone H} (au centre de la bande A) \emph{rétrécissent}, tandis que la \textbf{bande A} (sombre) conserve une longueur constante.
\end{itemize}

\textbf{Interprétation :} la bande A correspond à la zone occupée par les filaments épais de myosine ; sa longueur constante signifie que les filaments de myosine ne se raccourcissent pas. La bande I correspond aux filaments fins d'actine non recouverts par la myosine, et la zone H à la partie centrale de la bande A où seule la myosine est présente : leur rétrécissement signifie que les filaments d'actine \emph{pénètrent davantage} entre les filaments de myosine.

\textbf{Conclusion :} la contraction résulte du \textbf{glissement des filaments fins d'actine entre les filaments épais de myosine}, et non du raccourcissement des filaments eux-mêmes.
\end{experience}

\begin{propriete}[Théorie du glissement des filaments (Huxley, 1954)]
Lors de la contraction musculaire, la \textbf{longueur des filaments d'actine et de myosine ne change pas} : seul leur \textbf{glissement relatif} raccourcit le sarcomère. Les lignes Z sont tractées l'une vers l'autre par les filaments d'actine, entraînés par les têtes des molécules de myosine. Le raccourcissement de la fibre (et donc du muscle entier) est la somme des raccourcissements de tous les sarcomères disposés en série.
\end{propriete}

\begin{exemplebox}[Un ordre de grandeur à connaître]
Un sarcomère de muscle squelettique mesure environ \SI{2,5}{\micro\meter} au repos. En contraction maximale, il peut descendre jusqu'à environ \SI{1,5}{\micro\meter}, soit un raccourcissement de l'ordre de \SI{40}{\percent}. Comme une myofibrille contient des milliers de sarcomères en série, le muscle entier peut se raccourcir d'un tiers de sa longueur, ce qui suffit à rapprocher ses deux points d'insertion sur les os et à produire le mouvement (par exemple la flexion de l'avant-bras lorsque le biceps se contracte).
\end{exemplebox}

\begin{popfigure}
\begin{center}
\begin{tikzpicture}[scale=1.0]
% ===== Sarcomère au repos (Longueur 5 unités = 2.5 µm) =====
% Lignes Z
\draw[line width=2pt,popdark] (0,0) -- (0,2.4);
\draw[line width=2pt,popdark] (5,0) -- (5,2.4);
\node[below,font=\small\bfseries] at (0,-0.1) {Z};
\node[below,font=\small\bfseries] at (5,-0.1) {Z};
% Filaments Myosine (Épais, Bleu) : Centrés, longueur 1.8 unités (0.9 µm)
\foreach \y in {0.5,1.2,1.9} {\draw[line width=2.5pt,popblue] (1.6,\y) -- (3.4,\y);}
% Filaments Actine (Fins, Orange) : Longueur 2.0 unités (1.0 µm), chevauchement 0.4 unités sur myosine
\foreach \y in {0.8,1.5} {\draw[line width=1.2pt,poporange] (0,\y) -- (2.0,\y);} % Gauche : 0 à 2.0
\foreach \y in {0.8,1.5} {\draw[line width=1.2pt,poporange] (5,\y) -- (3.0,\y);} % Droite : 5 à 3.0
% Annotations Bandes
\draw[<->,popmuted] (0,2.7) -- (1.6,2.7); \node[above,font=\scriptsize] at (0.8,2.7) {Bande I};
\draw[<->,popmuted] (1.6,2.7) -- (3.4,2.7); \node[above,font=\scriptsize] at (2.5,2.7) {Bande A};
\draw[<->,popmuted] (2.0,-0.55) -- (3.0,-0.55); \node[below,font=\scriptsize] at (2.5,-0.55) {Zone H};
\draw[<->,popdark] (0,-1.0) -- (5,-1.0); \node[below,font=\small] at (2.5,-1.0) {\SI{2,5}{\micro\meter} (repos)};
\node[font=\small\bfseries,popdark] at (2.5,3.3) {Sarcomère au repos};
% ===== Sarcomère contracté (Longueur 3.5 unités = 1.75 µm ~ 1.5 µm) =====
\begin{scope}[xshift=8cm]
% Lignes Z
\draw[line width=2pt,popdark] (0,0) -- (0,2.4);
\draw[line width=2pt,popdark] (3.5,0) -- (3.5,2.4);
\node[below,font=\small\bfseries] at (0,-0.1) {Z};
\node[below,font=\small\bfseries] at (3.5,-0.1) {Z};
% Filaments Myosine (Même longueur 1.8, centrés sur 1.75) : 0.85 à 2.65
\foreach \y in {0.5,1.2,1.9} {\draw[line width=2.5pt,popblue] (0.85,\y) -- (2.65,\y);}
% Filaments Actine (Glissés vers le centre) : Gauche 0 à ~1.55, Droite 3.5 à ~1.95
\foreach \y in {0.8,1.5} {\draw[line width=1.2pt,poporange] (0,\y) -- (1.55,\y);}
\foreach \y in {0.8,1.5} {\draw[line width=1.2pt,poporange] (3.5,\y) -- (1.95,\y);}
% Annotations Bandes
\draw[<->,popmuted] (0,2.7) -- (0.85,2.7); \node[above,font=\scriptsize] at (0.42,2.7) {I};
\draw[<->,popmuted] (0.85,2.7) -- (2.65,2.7); \node[above,font=\scriptsize] at (1.75,2.7) {A};
% Zone H quasi nulle (actine se chevauche au centre), on ne la trace pas ou trait minime
\draw[<->,popdark] (0,-1.0) -- (3.5,-1.0); \node[below,font=\small] at (1.75,-1.0) {\SI{1,5}{\micro\meter} (contracté)};
\node[font=\small\bfseries,popdark] at (1.75,3.3) {Sarcomère contracté};
\end{scope}
\end{tikzpicture}
\end{center}
\legende{Comparaison d'un sarcomère au repos et en contraction. Les filaments épais de myosine (bleu) gardent une longueur constante : la bande A ne change pas. Les filaments fins d'actine (orange) glissent vers le centre : les bandes I et la zone H rétrécissent, les lignes Z se rapprochent.}
\end{popfigure}

\courssub{Le signal déclencheur : le message nerveux et l'ion calcium}

Un muscle squelettique ne se contracte que sur \textbf{ordre nerveux}. Le motoneurone qui commande la fibre musculaire libère un neurotransmetteur, l'\cle{acétylcholine}, au niveau de la \cle{plaque motrice} (jonction neuromusculaire). Cette libération provoque une dépolarisation de la membrane de la fibre musculaire (le \cle{sarcolemme}) : c'est le \emph{potentiel d'action musculaire}.

Ce signal électrique se propage à la surface de la fibre et pénètre en profondeur grâce aux \textbf{tubules T} (ou tubules transverses), invaginations du sarcolemme qui cheminent au contact du \cle{réticulum sarcoplasmique}, un réseau de saccules intracellulaires qui stockent les ions calcium \ce{Ca^{2+}}.

\begin{propriete}[Rôle déclencheur du calcium]
L'arrivée du potentiel d'action au niveau des tubules T provoque l'ouverture des canaux calciques (récepteurs ryanodine) du réticulum sarcoplasmique : les ions \ce{Ca^{2+}} sont massivement libérés dans le cytoplasme (le sarcoplasme) de la fibre. Le \ce{Ca^{2+}} est le \textbf{signal intracellulaire déclencheur} de la contraction : sans élévation de la concentration sarcoplasmique en \ce{Ca^{2+}}, aucune contraction n'est possible, même en présence d'ATP.
\end{propriete}

Comment le calcium agit-il ? Au repos, les sites de fixation de la myosine sur les filaments d'actine sont \textbf{masqués} par un complexe de protéines régulatrices : la \cle{tropomyosine} (long filament enroulé dans le sillon de l'hélice d'actine) et la \cle{troponine} (complexe de trois sous-unités). Lorsque le \ce{Ca^{2+}} se fixe sur la sous-unité TnC de la troponine, celle-ci change de conformation et \emph{déplace la tropomyosine} : les sites actifs de l'actine sont alors \textbf{démasqués} et accessibles aux têtes de myosine.

\courssub{Le cycle des ponts actine-myosine : le moteur moléculaire}

Le filament épais de myosine porte des centaines de \textbf{têtes} globulaires (domaines S1) orientées vers l'extérieur. Ces têtes possèdent deux propriétés essentielles : elles peuvent se \emph{fixer sur l'actine} et elles possèdent une \emph{activité enzymatique ATPasique} (elles hydrolysent l'ATP). La contraction résulte de la répétition, par des millions de têtes, d'un cycle en quatre étapes :

\begin{enumerate}
\item \textbf{Fixation.} Les sites de l'actine étant démasqués par le \ce{Ca^{2+}}, la tête de myosine (chargée des produits de l'hydrolyse de l'ATP, ADP et Pi) se fixe sur l'actine : c'est la formation du \cle{pont actine-myosine}.
\item \textbf{Basculement (coup de rame).} La libération de l'ADP et du Pi provoque le basculement de la tête de myosine, comme le coup de rame d'un canotier : ce mouvement \textbf{tire le filament d'actine vers le centre du sarcomère} sur une distance d'environ \SI{10}{\nano\meter}. C'est l'étape qui produit la force et le glissement.
\item \textbf{Détachement.} Une nouvelle molécule d'\textbf{ATP se fixe sur la tête de myosine}, ce qui provoque le \textbf{détachement} du pont actine-myosine.
\item \textbf{Réactivation.} La tête hydrolyse cet ATP (\ce{ATP -> ADP + Pi}) : l'énergie libérée la « réarme », c'est-à-dire la ramène en position initiale (conformation « tendue »), prête pour un nouveau cycle.
\end{enumerate}

Tant que le calcium est présent (sites démasqués) et que l'ATP est disponible, le cycle se répète (environ 5 fois par seconde et par tête à température ambiante, plus vite \textit{in vivo}) : les filaments glissent progressivement et le sarcomère se raccourcit.

\begin{popfigure}
\begin{center}
\begin{tikzpicture}[scale=0.92]
% étape 1 : fixation
\begin{scope}
\draw[line width=1.5pt,poporange] (-1.2,0) -- (1.6,0);
\fill[poporange!60] (0.6,0) circle (0.12);
\draw[fill=popblue!70] (0.2,0.55) circle (0.35);
\draw[line width=2pt,popblue] (0.2,0.9) -- (0.2,1.7);
\node[font=\scriptsize] at (0.2,0.55) {ADP+Pi};
\node[font=\small\bfseries,popdark] at (0.2,2.1) {1. Fixation};
\end{scope}
% étape 2 : basculement
\begin{scope}[xshift=3.4cm]
\draw[line width=1.5pt,poporange] (-1.2,0) -- (1.6,0);
\fill[poporange!60] (0.6,0) circle (0.12);
\draw[fill=popblue!70] (0.45,0.5) circle (0.35);
\draw[line width=2pt,popblue] (0.45,0.85) -- (1.15,1.6);
\draw[->,popgreen,thick] (-0.9,0.35) -- (-0.2,0.35);
\node[font=\scriptsize,popgreen] at (-0.55,0.6) {glissement};
\node[font=\small\bfseries,popdark] at (0.4,2.1) {2. Basculement};
\node[font=\scriptsize] at (0.4,-0.5) {ADP et Pi libérés};
\end{scope}
% étape 3 : détachement
\begin{scope}[xshift=6.9cm]
\draw[line width=1.5pt,poporange] (-1.2,0) -- (1.6,0);
\fill[poporange!60] (0.6,0) circle (0.12);
\draw[fill=popblue!70] (1.0,0.85) circle (0.35);
\draw[line width=2pt,popblue] (1.0,1.2) -- (1.0,1.9);
\node[font=\scriptsize] at (1.0,0.85) {ATP};
\node[font=\small\bfseries,popdark] at (0.4,2.1) {3. Détachement};
\node[font=\scriptsize] at (0.4,-0.5) {fixation de l'ATP};
\end{scope}
% étape 4 : réactivation
\begin{scope}[xshift=10.3cm]
\draw[line width=1.5pt,poporange] (-1.2,0) -- (1.6,0);
\fill[poporange!60] (0.6,0) circle (0.12);
\draw[fill=popblue!70] (1.0,0.85) circle (0.35);
\draw[line width=2pt,popblue] (1.0,1.2) -- (0.5,1.9);
\node[font=\scriptsize] at (1.0,0.85) {ADP+Pi};
\node[font=\small\bfseries,popdark] at (0.4,2.1) {4. Réactivation};
\node[font=\scriptsize] at (0.4,-0.5) {hydrolyse de l'ATP};
\end{scope}
% flèches du cycle
\draw[->,thick,popmuted] (1.7,1.0) -- (2.1,1.0);
\draw[->,thick,popmuted] (5.2,1.0) -- (5.6,1.0);
\draw[->,thick,popmuted] (8.6,1.0) -- (9.0,1.0);
\draw[->,thick,popmuted] (11.2,-0.9) .. controls (6,-1.8) .. (0.4,-1.0);
\node[font=\scriptsize,popmuted] at (6,-1.75) {nouveau cycle tant que \ce{Ca^{2+}} et ATP sont présents};
\end{tikzpicture}
\end{center}
\legende{Le cycle du pont actine-myosine. Le filament d'actine est figuré en orange (point : site actif démasqué par le \ce{Ca^{2+}}), la tête de myosine en bleu. L'hydrolyse de l'ATP (étape 4) réactive la tête (conformation tendue) ; la fixation d'un nouvel ATP (étape 3) provoque le détachement du pont.}
\end{popfigure}

\begin{attention}[L'ATP ne « fait pas contracter » le muscle]
Erreur très fréquente au Baccalauréat : écrire que « l'ATP fait se contracter le muscle » ou que « l'énergie de l'ATP déclenche la contraction ». En réalité :
\begin{itemize}
\item le \textbf{déclencheur} de la contraction est le \ce{Ca^{2+}} libéré par le réticulum sarcoplasmique, pas l'ATP ;
\item l'hydrolyse de l'ATP sert à \textbf{détacher} les têtes de myosine de l'actine (étape 3) et à les \textbf{réactiver} (réarmer, étape 4) pour le cycle suivant ;
\item l'ATP sert aussi au \textbf{repompage du calcium} vers le réticulum (via les ATPases \ce{Ca^{2+}}), donc à la relaxation.
\end{itemize}
Formulation correcte attendue : « la contraction musculaire \emph{consomme} de l'ATP, indispensable au fonctionnement cyclique des ponts actine-myosine et au retour du calcium dans le réticulum ».
\end{attention}

\begin{experience}[Preuve expérimentale du rôle de l'ATP : la rigidité en absence d'ATP]
\textbf{Protocole :} on place des myofibrilles isolées dans deux milieux : le milieu A contient du \ce{Ca^{2+}} et de l'ATP ; le milieu B contient du \ce{Ca^{2+}} mais \emph{pas} d'ATP.

\textbf{Observations :} dans le milieu A, les myofibrilles se raccourcissent puis se relâchent normalement. Dans le milieu B, les ponts actine-myosine se forment (le calcium a démasqué les sites) mais \textbf{restent figés} : les myofibrilles demeurent rigides, incapables de se relâcher ni de poursuivre le cycle.

\textbf{Interprétation :} sans ATP, les têtes de myosine ne peuvent ni se détacher de l'actine ni être réactivées. L'ATP est donc indispensable au \emph{cycle} des ponts, et en particulier à leur rupture.

\textbf{Conclusion :} la contraction est un phénomène \textbf{consommateur d'ATP} ; l'ATP est nécessaire au détachement et à la réactivation des ponts.
\end{experience}

\begin{savaistu}[La rigidité cadavérique : une preuve « naturelle » du rôle de l'ATP]
Quelques heures après la mort, le corps se raidit : c'est le \emph{rigor mortis}. L'explication est exactement celle de l'expérience précédente. Après la mort, la respiration cellulaire s'arrête : l'ATP n'est plus renouvelé et sa concentration s'effondre. Or les membranes du réticulum sarcoplasmique deviennent perméables et relâchent leur calcium dans le sarcoplasme : les ponts actine-myosine se forment massivement. En l'absence d'ATP, ces ponts \textbf{ne peuvent plus se rompre} : les muscles restent figés en état de contraction, d'où la rigidité cadavérique. Celle-ci disparaît après un à trois jours, lorsque les protéines contractiles se dégradent par autolyse. Les médecins légistes utilisent d'ailleurs l'état du rigor mortis pour estimer l'heure du décès — un exemple concret où la biologie moléculaire éclaire la pratique médico-légale.
\end{savaistu}

\courssub{La relaxation : un processus actif, lui aussi consommateur d'ATP}

Lorsque l'influx nerveux cesse, les canaux calciques du réticulum se referment et des \textbf{pompes à calcium} (ATPases \ce{Ca^{2+}} du réticulum sarcoplasmique, ou SERCA) rejettent activement le \ce{Ca^{2+}} du sarcoplasme vers l'intérieur du réticulum sarcoplasmique. Ce transport se fait \emph{contre le gradient de concentration} : il \textbf{consomme de l'ATP} (hydrolyse d'un ATP pour 2 \ce{Ca^{2+}} transportés).

La concentration sarcoplasmique en \ce{Ca^{2+}} chute : le calcium se défixe de la troponine, la tropomyosine reprend sa place et \textbf{masque à nouveau} les sites actifs de l'actine. Les ponts ne peuvent plus se former : le muscle se \cle{relâche}. La relaxation est donc, comme la contraction, un phénomène \emph{actif} qui exige de l'ATP — d'où la nécessité absolue de renouveler l'ATP, que nous étudierons dans la section suivante.

\courssub{Schéma d'interprétation du mécanisme de la contraction (savoir-faire exigible)}

\begin{methode}[Construire le schéma d'interprétation]
Pour réaliser un schéma d'interprétation correct et complet :
\begin{enumerate}
\item identifier la \textbf{cause initiale} (le message nerveux) et la \textbf{conséquence finale} (le raccourcissement du sarcomère) ;
\item enchaîner les étapes \textbf{intermédiaires dans l'ordre causal strict} : libération de \ce{Ca^{2+}} $\rightarrow$ démasquage des sites $\rightarrow$ cycle des ponts $\rightarrow$ glissement des filaments ;
\item relier chaque étape à la suivante par une \textbf{flèche causale} légendée si nécessaire ;
\item mentionner les \textbf{consommations d'ATP} aux étapes concernées (cycle des ponts, repompage du calcium).
\end{enumerate}
\end{methode}

\begin{popfigure}
\begin{center}
\begin{tikzpicture}[node distance=0.5cm,
etape/.style={rectangle,rounded corners,draw=popblue,fill=popblueL,thick,align=center,minimum height=0.9cm,inner sep=4pt,font=\small},
energie/.style={rectangle,draw=poporange,fill=poporangeL,thick,align=center,font=\scriptsize,inner sep=3pt}]
\node[etape, align=center] (n1){Message nerveux au niveau\\de la plaque motrice};
\node[etape,below=of n1, align=center] (n2){Potentiel d'action propagé\\par le sarcolemme et les tubules T};
\node[etape,below=of n2, align=center] (n3){Libération de \ce{Ca^{2+}} par le\\réticulum sarcoplasmique};
\node[etape,below=of n3, align=center] (n4){Fixation du \ce{Ca^{2+}} sur la troponine :\\démasquage des sites de l'actine};
\node[etape,below=of n4, align=center] (n5){Cycle des ponts actine-myosine\\(fixation, basculement, détachement, réactivation)};
\node[etape,below=of n5, align=center] (n6){Glissement des filaments d'actine\\vers le centre du sarcomère};
\node[etape,fill=popgreenL,draw=popgreen,below=of n6, align=center] (n7){Raccourcissement du sarcomère\\= contraction du muscle};
\draw[->,thick,popdark] (n1) -- (n2);
\draw[->,thick,popdark] (n2) -- (n3);
\draw[->,thick,popdark] (n3) -- (n4);
\draw[->,thick,popdark] (n4) -- (n5);
\draw[->,thick,popdark] (n5) -- (n6);
\draw[->,thick,popdark] (n6) -- (n7);
\node[energie,right=1.2cm of n5, align=center] (e1){Consommation d'ATP\\(détachement et réactivation\\des têtes de myosine)};
\draw[->,thick,poporange] (e1) -- (n5);
\node[energie,right=1.2cm of n3, align=center] (e2){Relaxation : repompage du \ce{Ca^{2+}}\\vers le réticulum (consomme de l'ATP)};
\draw[->,thick,poporange,dashed] (e2) -- (n3);
\end{tikzpicture}
\end{center}
\legende{Schéma d'interprétation du mécanisme de la contraction musculaire : chaîne causale du message nerveux au raccourcissement du sarcomère. Les encadrés orange rappellent les étapes consommatrices d'ATP.}
\end{popfigure}

\begin{exoresolu}[Interpréter une expérience sur le rôle du calcium et de l'ATP]
\textbf{Énoncé.} On prépare trois lots de myofibrilles de grenouille placés dans les milieux suivants : lot 1 : ATP seul ; lot 2 : \ce{Ca^{2+}} seul ; lot 3 : ATP + \ce{Ca^{2+}}. On mesure la longueur des sarcomères après 5 minutes : lot 1 : \SI{2,5}{\micro\meter} ; lot 2 : \SI{2,5}{\micro\meter} ; lot 3 : \SI{1,6}{\micro\meter}. Interpréter ces résultats.
\tcblower
\textbf{Solution détaillée.}

\emph{Analyse des résultats.} Les sarcomères ne se sont raccourcis (de \SI{2,5}{\micro\meter} à \SI{1,6}{\micro\meter}) que dans le lot 3, qui contient à la fois l'ATP et le calcium. Dans les lots 1 et 2, la longueur reste inchangée : aucune contraction n'a eu lieu.

\emph{Interprétation.}
\begin{itemize}
\item Lot 1 (ATP sans \ce{Ca^{2+}}) : l'énergie est disponible, mais en l'absence de calcium, la tropomyosine masque les sites actifs de l'actine : les ponts ne peuvent pas se former, il n'y a pas de contraction. L'ATP seul ne suffit donc pas.
\item Lot 2 (\ce{Ca^{2+}} sans ATP) : le calcium démasque les sites et les ponts actine-myosine peuvent se former (liaison forte rigor), mais sans ATP les têtes de myosine ne peuvent ni se détacher ni être réactivées : le cycle est bloqué dès le premier pas, il n'y a pas de glissement durable des filaments.
\item Lot 3 (ATP + \ce{Ca^{2+}}) : le calcium démasque les sites et l'ATP permet le fonctionnement cyclique des ponts : les filaments d'actine glissent entre les filaments de myosine, les sarcomères se raccourcissent.
\end{itemize}

\emph{Conclusion.} La contraction musculaire exige la présence \textbf{simultanée} du calcium (signal déclencheur) et de l'ATP (source d'énergie indispensable au cycle des ponts). Cette expérience illustre la double exigence énergétique et ionique du mécanisme contractile.
\end{exoresolu}

\begin{aretenir}[L'essentiel du mécanisme de la contraction]
\begin{itemize}
\item La contraction repose sur le \textbf{glissement des filaments d'actine entre les filaments de myosine}, sans raccourcissement des filaments (théorie de Huxley) : les bandes I et la zone H rétrécissent, la bande A reste constante, les lignes Z se rapprochent.
\item Le \textbf{message nerveux} (plaque motrice, potentiel d'action, tubules T) provoque la libération de \ce{Ca^{2+}} par le réticulum sarcoplasmique : le calcium est le \textbf{déclencheur}.
\item Le \ce{Ca^{2+}} démasque les sites de l'actine (via troponine et tropomyosine), permettant le \textbf{cycle des ponts actine-myosine} : fixation (ADP+Pi), basculement (libération ADP+Pi), détachement (fixation ATP), réactivation (hydrolyse ATP).
\item La contraction \textbf{consomme de l'ATP} (cycle des ponts) ; la relaxation aussi (repompage du calcium par les SERCA). En absence d'ATP, les ponts restent figés : c'est la rigidité cadavérique.
\end{itemize}
\end{aretenir}

Ainsi, toute contraction musculaire — du sprinteur qui s'élance sur la piste du stade Ahmadou Ahidjo au porteur qui soulève un sac de cacao — puise dans un stock d'ATP très limité, épuisé en quelques secondes. D'où la question qui fera l'objet de la section suivante : comment la cellule musculaire \emph{renouvelle-t-elle} en permanence son ATP ?

\coursec{IV. La nécessité du renouvellement de l'ATP : les voies de restauration}

\courssub{Le constat expérimental : l'ATP, une monnaie d'échange à flux tendu}

\emph{Rappel de la section précédente.} Nous avons vu que le mécanisme de la contraction (théorie des filaments glissants) repose sur le cycle des ponts croisés : la tête de myosine hydrolyse une molécule d'ATP pour se détacher de l'actine, se reculer (amorçage) et former un nouveau pont plus loin. Chaque cycle de glissement consomme \textbf{une molécule d'ATP}. Lors d'un effort intense, des millions de cycles se produisent chaque seconde dans chaque fibre. La question centrale est donc : \textbf{comment la cellule maintient-elle la concentration en ATP alors que la consommation est massive ?}

\begin{experience}[Biopsies musculaires et dosage des métabolites énergétiques]
\textbf{Matériel :} Échantillons de muscle \emph{vasto-latéral} (cuisse) prélevés par biopsie percutanée (aiguille de Bergström) chez des volontaires sains, au repos, puis après des exercices de durée et d'intensité contrôlées (ex. : sprint de 6 s, exercice à 100 \% \emph{VO}$_{2\text{max}}$ pendant 30 s, 1 min, 3 min). Dosage enzymatique spectrophotométrique de l'ATP, de la phosphocréatine (PCr), du glycogène, du lactate.
\textbf{Protocole :} Congélation instantanée des biopsies dans l'azote liquide (pour figer l'état métabolique), extraction acide (acide perchlorique), neutralisation, dosage.
\textbf{Observations (résultats types) :}
\begin{itemize}
    \item Au repos : [ATP] $\approx$ \SI{25}{\milli\mole\per\kilo\gram} de muscle sec ; [PCr] $\approx$ \SI{80}{\milli\mole\per\kilo\gram}.
    \item Après 6 s de sprint maximal : [ATP] $\approx$ \SI{22}{\milli\mole\per\kilo\gram} (légère baisse) ; [PCr] $\approx$ \SI{30}{\milli\mole\per\kilo\gram} (chute drastique de 60 \%).
    \item Après 30 s d'effort maximal : [ATP] $\approx$ \SI{20}{\milli\mole\per\kilo\gram} ; [PCr] $\approx$ \SI{10}{\milli\mole\per\kilo\gram} (quasiment épuisée) ; apparition massive de lactate.
    \item Après 3 min d'effort intense : [ATP] $\approx$ \SI{18}{\milli\mole\per\kilo\gram} (toujours maintenu) ; [PCr] $\approx$ \SI{5}{\milli\mole\per\kilo\gram} ; glycogène fortement diminué.
\end{itemize}
\textbf{Interprétation :} Le stock d'ATP reste \textbf{quasi constant} (variation $< 30\%$) alors que la phosphocréatine s'effondre et que le glycogène diminue. Cela prouve que l'ATP hydrolysée est \textbf{immédiatement resynthétisée} à partir d'autres réserves. L'ATP n'est pas un stock, c'est un \textbf{intermédiaire métabolique à turnover élevé}.
\textbf{Conclusion :} La contraction musculaire ne puise pas directement dans un « réservoir d'ATP » ; elle déclenche des \textbf{voies de restauration de l'ATP} (ou filières énergétiques) qui utilisent des substrats de réserve (phosphocréatine, glycogène, lipides) pour reformer l'ATP à partir de l'ADP et du P$_i$.
\end{experience}

\begin{popfigure}
\begin{tikzpicture}
\begin{axis}[
    width=\linewidth, height=0.55\linewidth,
    xlabel={Durée de l'effort intense (s)}, ylabel={Concentration relative (\% valeur repos)},
    xmin=0, xmax=180, ymin=0, ymax=110,
    xtick={0,30,60,90,120,150,180}, ytick={0,20,40,60,80,100},
    legend style={at={(0.98,0.98)}, anchor=north east, draw=popline, fill=popcream},
    grid=major, grid style={popline!30},
]
\addplot[popblue, very thick, mark=*, mark size=2.5] coordinates {
    (0,100) (6,92) (30,85) (60,80) (120,75) (180,72)
};
\addlegendentry{ATP}
\addplot[poporange, very thick, mark=square*, mark size=2.5] coordinates {
    (0,100) (6,40) (30,15) (60,8) (120,5) (180,5)
};
\addlegendentry{Phosphocréatine (PCr)}
\addplot[popgreen, very thick, mark=triangle*, mark size=2.5] coordinates {
    (0,100) (6,95) (30,80) (60,60) (120,35) (180,20)
};
\addlegendentry{Glycogène}
\end{axis}
\end{tikzpicture}
\legende{Évolution des concentrations relatives d'ATP, de phosphocréatine et de glycogène dans le muscle \emph{vasto-latéral} au cours d'un exercice intense continu (données compilées d'après Harris et al., 1976 ; Spriet et al., 1987). L'ATP est maintenu (homéostasie) au prix de l'épuisement successif de la PCr puis du glycogène.}
\end{popfigure}

\begin{attention}[Piège classique au Bac : « Le stock d'ATP s'épuise »]
\emph{Erreur fréquente :} Écrire « L'ATP diminue rapidement pendant l'effort » ou « Le muscle manque d'ATP ».
\emph{Réalité expérimentale :} La concentration en ATP chute très peu (homéostasie stricte). Si l'ATP chutait vraiment, la rigidité cadavérique s'installerait (pas d'ATP pour détacher les ponts myosine-actine). C'est la \textbf{phosphocréatine} qui chute en premier (secondes), puis le \textbf{glycogène} (minutes). Au Bac, précisez toujours : « Le taux d'ATP reste quasi constant grâce à l'activation immédiate des voies de resynthèse. »
\end{attention}

\begin{aretenir}[Pourquoi renouveler l'ATP ?]
\begin{itemize}
    \item \textbf{Stock initial ridicule :} $\sim$ \SI{25}{\milli\mole\per\kilo\gram} $\approx$ \SI{5}{\milli\mole\per\litre} de volume cellulaire. Pour une fibre de \SI{100}{\micro\metre} de diamètre, cela ne couvre que \textbf{1 à 3 secondes} d'effort maximal.
    \item \textbf{Fonctions vitales de l'ATP :} Cycle des ponts croisés (contraction) + Pompes $\text{Ca}^{2+}$-ATPase (relâchement, rétention $\text{Ca}^{2+}$ dans le RS) + Pompe $\text{Na}^+/\text{K}^+$-ATPase (potentiel de membrane, excitabilité).
    \item \textbf{Principe :} $\Delta G$ de l'hydrolyse ATP $\rightarrow$ ADP + P$_i$ maintenu très négatif ($\approx -55$ à $-60$ \SI{}{\kilo\joule\per\mole}) grâce à la resynthèse permanente. C'est ce $\Delta G$ qui fournit l'énergie libre utile.
\end{itemize}
\end{aretenir}

\courssub{Les trois voies de restauration de l'ATP : présentation unifiée}

La cellule musculaire dispose de trois systèmes enzymatiques pour reformer l'ATP à partir de l'ADP. Ils diffèrent par leur \textbf{vitesse maximale (puissance)}, leur \textbf{capacité (durée totale)}, leur \textbf{substrat}, leur \textbf{localisation}, leur \textbf{rendement en ATP} et leurs \textbf{produits finaux}.

\begin{definition}[Voie anaérobie alactique (ou filière phosphagène)]
C'est la resynthèse de l'ATP par \textbf{transfert de groupe phosphate} depuis la \textbf{phosphocréatine (PCr)} vers l'ADP, catalysée par la \textbf{créatine kinase (CK)}, \textbf{sans consommation d'oxygène} et \textbf{sans production de lactate}.
\[
\ce{PCr^{2-} + ADP^{3-} + H+ ->[Creatine\ Kinase] Cr + ATP^{4-}} \qquad \Delta G^{\circ\prime} \approx -12,6\ \text{kJ/mol}
\]
\textbf{Caractères :} Extrêmement rapide (vitesse max $\approx$ \SI{4}{\milli\mole\per\kilo\gram\per\second}), puissance maximale, capacité très faible (stock PCr $\approx$ 3-4 $\times$ stock ATP), durée $\approx$ \SIrange{5}{10}{s} d'effort maximal. Aucun déchet métabolique (la créatine est recyclée au repos).
\end{definition}

\begin{definition}[Voie anaérobie lactique (ou glycolyse anaérobie)]
C'est la dégradation du \textbf{glucose} (issu du glycogène musculaire ou de la glycémie) en \textbf{2 pyruvate}, puis réduction du pyruvate en \textbf{lactate} par la \textbf{lactate déshydrogénase (LDH)}, avec production nette de \textbf{2 ATP} par glucose (niveau substrat), \textbf{sans oxygène}.
\[
\ce{Glycogen_n + Pi -> Glycogen_{n-1} + Glucose-1-P} \quad (\text{glycogène phosphorylase})
\]
\[
\ce{Glucose + 2 ADP + 2 Pi + 2 NAD+ -> 2 Lactate- + 2 ATP + 2 H+ + 2 NAD+} \quad (\text{bilan global simplifié pour le glucose})
\]
\textbf{Précision :} Si le substrat est le glycogène, le rendement net est de \textbf{3 ATP} (économie de l'hexokinase).
\textbf{Caractères :} Rapide (mise en route $\approx$ quelques secondes, puissance élevée), capacité moyenne (stock glycogène $\approx$ \SIrange{300}{600}{\milli\mole\per\kilo\gram} sec), durée efficace $\approx$ \SIrange{10}{180}{s}. Limité par l'acidose (chute du pH $\rightarrow$ inhibition de la phosphofructokinase PFK et de la contraction).
\end{definition}

\begin{propriete}[Voie aérobie (oxydation complète) - Savoir admis au Bac]
C'est l'oxydation complète du \textbf{pyruvate} (issu de la glycolyse), des \textbf{acides gras} (beta-oxydation) et des \textbf{acides aminés} dans les \textbf{mitochondries} (cycle de Krebs + chaîne respiratoire), \textbf{nécessitant l'oxygène} comme accepteur final d'électrons.
\[
\ce{C6H12O6 + 6 O2 -> 6 CO2 + 6 H2O} \quad (\Delta G^{\circ\prime} \approx -2870\ \text{kJ/mol})
\]
\textbf{Rendement :} $\approx$ \textbf{30 à 32 ATP} par glucose (souvent arrondi à \textbf{36 ATP} dans les manuels anciens ; la valeur moderne consensuelle est 30-32 ATP compte tenu du coût du transport mitochondrial NADH et ADP/ATP). Soit \textbf{15 à 18 fois plus} que la glycolyse seule.
\textbf{Caractères :} Lente à se mettre en place (cinétique $\text{VO}_2$ : $\tau \approx$ \SIrange{20}{30}{s}), puissance modérée, capacité quasi illimitée (stocks lipides $\gg$ glucides), durée : heures. Déchets : $\text{CO}_2$ (éliminé par poumons) et $\text{H}_2\text{O}$.
\end{propriete}

\begin{popfigure}
\begin{tikzpicture}[
    node distance=1.2cm and 1.5cm,
    box/.style={draw=popdark, thick, rounded corners, fill=popcream, minimum height=1.8cm, text width=3.2cm, align=center, font=\small},
    arrow/.style={-Stealth, thick, popdark},
    labelbox/.style={font=\footnotesize, text width=2.5cm, align=center}
]
% Nœuds centraux
\node[box, fill=popblueL] (adp) {ADP + P$_i$};
\node[box, fill=poppinkL, right=of adp] (atp) {ATP};
\node[box, fill=popgreenL, below=of adp, align=center] (pcr){Phosphocréatine (PCr)\\(Stock : $\sim$ 80 mmol/kg)\\\emph{Créatine Kinase}};
\node[box, fill=poporangeL, below=of atp, align=center] (glyc){Glycogène / Glucose\\(Stock : $\sim$ 400 mmol/kg)\\\emph{Glycolyse cytoplasmique}};
\node[box, fill=poppurpleL, below=of glyc, align=center] (mito){Pyruvate + Acides Gras\\(Stocks quasi illimités)\\\emph{Mitochondries : Krebs + Ch. Resp.}};

% Flèches et annotations
\draw[arrow] (pcr.north east) -- node[above, sloped, labelbox, align=center]{\textbf{Anaérobie Alactique}\\\SI{~4}{mmol/kg/s}\\\SI{5-10}{s}\\\textit{0 ATP/glucose\\ (transfert direct)}} (adp.east);
\draw[arrow] (glyc.north east) -- node[above, sloped, labelbox, align=center]{\textbf{Anaérobie Lactique}\\\SI{~2-3}{mmol/kg/s}\\\SI{10-180}{s}\\\textit{2-3 ATP/glucose\\ Lactate + H+}} (atp.west);
\draw[arrow] (mito.north east) -- node[above, sloped, labelbox, align=center]{\textbf{Aérobie}\\\SI{~1}{mmol/kg/s}\\\textit{Heures}\\\textit{$\sim$30-32 ATP/glucose\\ CO2 + H2O}} (atp.west);

% Légendes substrats
\node[labelbox, left=0.2cm of pcr] {Substrat : PCr};
\node[labelbox, left=0.2cm of glyc] {Substrat : Glucides};
\node[labelbox, left=0.2cm of mito] {Substrat : Glucides + Lipides};
\end{tikzpicture}
\legende{Schéma récapitulatif des trois voies de restauration de l'ATP musculaire. Chaque voie est caractérisée par son substrat, son enzyme/système clé, sa vitesse maximale (puissance), sa durée efficace (capacité) et son rendement en ATP. La flèche épaisse vers ADP+P$_i$ symbolise le flux de phosphate à haut transfert d'énergie.}
\end{popfigure}

\courssub{Analyse comparative et dominance des filières selon l'effort}

\begin{experience}[Contribution relative des filières : expérience classique (Gastin, 2001) et modélisation]
\textbf{Document :} Mesure de la dépense énergétique totale (calorimétrie directe/indirecte) et partitionnement anaérobie (dette O$_2$ + accumulation lactate + déplétion PCr) vs aérobie ($\text{VO}_2$ mesurée) lors d'efforts exhaustifs de durées variées (10 s, 30 s, 60 s, 2 min, 5 min, 30 min, 60 min).
\textbf{Observations (courbe ci-dessous) :}
\begin{itemize}
    \item \textbf{0-10 s :} Filière alactique $> 70\%$ (sprinter 100 m, haltérophile).
    \item \textbf{10-30 s :} Transition alactique $\rightarrow$ lactique. Pic de puissance lactique $\approx$ 30 s.
    \item \textbf{30 s - 2 min :} Filière lactique dominante ($50-70\%$) (400 m, 800 m).
    \item \textbf{2-5 min :} Crossover : filière aérobie devient majoritaire ($> 50\%$).
    \item \textbf{$> 10$ min :} Filière aérobie $> 90\%$ (marathon, cyclisme sur route, marche).
\end{itemize}
\textbf{Interprétation :} La \textbf{puissance} (débit d'ATP/s) dicte le choix initial : alactique $>$ lactique $>$ aérobie. La \textbf{capacité} (stock substrat) dicte la durée : PCr s'épuise en secondes, glycogène en minutes, lipides en heures. L'intensité relative (\% $\text{VO}_{2\text{max}}$) module ce partage : à 130\% $\text{VO}_{2\text{max}}$, l'aérobie ne peut pas suivre, l'anaérobie compense.
\end{experience}

\begin{popfigure}
\begin{tikzpicture}
\begin{axis}[
    width=\linewidth, height=0.55\linewidth,
    xlabel={Durée de l'effort maximal (échelle logarithmique)}, ylabel={Contribution relative (\% de l'ATP total)},
    xmode=log, log basis x=10,
    xmin=1, xmax=3600, ymin=0, ymax=100,
    xtick={1,2,5,10,20,30,60,120,300,600,1800,3600},
    xticklabels={1s,2s,5s,10s,20s,30s,1min,2min,5min,10min,30min,1h},
    ytick={0,20,40,60,80,100},
    legend style={at={(0.5,-0.25)}, anchor=north, legend columns=3, draw=popline, fill=popcream},
    grid=major, grid style={popline!30},
    stack plots=y, area style
]
% Alactique : forte au début, chute vite
\addplot[popblue!80, fill=popblue!30] coordinates {
    (1,85) (2,75) (5,55) (10,30) (20,15) (30,8) (60,3) (120,1) (300,0) (600,0) (1800,0) (3600,0)
} \closedcycle;
\addlegendentry{Anaérobie Alactique (PCr)}

% Lactique : monte, pic vers 30-60s, descend
\addplot[poporange!80, fill=poporange!30] coordinates {
    (1,5) (2,10) (5,25) (10,45) (20,55) (30,60) (60,55) (120,40) (300,20) (600,10) (1800,3) (3600,1)
} \closedcycle;
\addlegendentry{Anaérobie Lactique (Glycolyse)}

% Aérobie : lent départ, domine après 2-3 min
\addplot[popgreen!80, fill=popgreen!30] coordinates {
    (1,10) (2,15) (5,20) (10,25) (20,30) (30,32) (60,42) (120,59) (300,80) (600,90) (1800,97) (3600,99)
} \closedcycle;
\addlegendentry{Aérobie (Oxydation)}
\end{axis}
\end{tikzpicture}
\legende{Modélisation de la contribution relative (\% de l'ATP total fourni) des trois filières énergétiques en fonction de la durée d'un effort \emph{maximal} continu (données synthétiques d'après Gastin, 2001 ; Spencer \& Gastin, 2001). Notez l'échelle logarithmique en abscisse. La zone hachurée (non montrée ici mais réelle) correspond à la variabilité inter-individuelle (entraînement, type de fibres).}
\end{popfigure}

\begin{exemplebox}[Application concrète : Le 400 m et le marathon au Cameroun]
\begin{itemize}
    \item \textbf{400 m (effort intense $\sim$ 45-50 s) :} Un sprinter comme ceux formés au Centre d'Excellence de Yaoundé mobilise : \textbf{Alactique} (premiers 5-6 s, départ blocs) $\rightarrow$ \textbf{Lactique dominante} (30-40 s suivantes, accumulation lactate $\approx$ 15-20 mmol/L, pH musculaire $\approx$ 6,5). L'aérobie ne contribue qu'à $\sim 20-30\%$. L'entraînement vise à repousser le seuil lactique et tamponner l'acidose (bicarbonate, carnosine).
    \item \textbf{Marathon (effort long $\sim$ 2 h 10 min pour l'élite) :} Coureur de fond (ex. : région de l'Ouest, altitude modérée). \textbf{Aérobie $> 97\%$}. Substrats : glucides (glycogène hépatique/musculaire + gels glucosés aux ravitos) \textbf{puis} lipides (après 1h30-2h, « mur du marathon » si glycogène épuisé). L'entraînement augmente la densité mitochondriale, la capillarisation et l'oxydation des lipides pour épargner le glycogène.
\end{itemize}
\end{exemplebox}

\begin{exoresolu}[Analyse de biopsies : interprétation de données chiffrées]
\textbf{Énoncé :} On réalise des biopsies du muscle \emph{vasto-latéral} chez un cycliste avant (T0), après 30 s (T1) et après 20 min (T2) d'un exercice intense constant (85\% $\text{VO}_{2\text{max}}$). Résultats (mmol/kg sec) :
\begin{center}
\begin{tabular}{|l|ccc|}\hline
Métabolite & T0 (Repos) & T1 (30 s) & T2 (20 min) \\ \hline
ATP & 24,5 & 22,1 & 21,8 \\
Phosphocréatine & 78,2 & 12,4 & 45,6 \\
Glycogène & 420 & 380 & 180 \\
Lactate & 5,1 & 65,3 & 18,2 \\ \hline
\end{tabular}
\end{center}
\textbf{Questions :}
\begin{enumerate}
    \item Commenter l'évolution du taux d'ATP. Que prouve-t-il ?
    \item Identifier la filière majoritaire à T1 et justifier par deux arguments chiffrés.
    \item Expliquer la remontée de la PCr à T2 par rapport à T1.
    \item Calculer la quantité d'ATP fournie par la glycolyse entre T0 et T1 (en mmol/kg sec), sachant que 1 glucose $\rightarrow$ 2 ATP + 2 Lactate.
\end{enumerate}
\textbf{Solution détaillée :}
\begin{enumerate}
    \item \textbf{ATP :} Variation faible : $(24,5-22,1)/24,5 \approx 9,8\%$ de baisse à T1 ; quasi stable à T2. \textbf{Preuve :} L'ATP est un intermédiaire à flux tendu, maintenu constant par resynthèse immédiate. Pas d'épuisement $\rightarrow$ pas de rigor mortis.
    \item \textbf{T1 (30 s) : Filière anaérobie lactique majoritaire.}
        \textbf{Arg 1 :} Chute massive PCr : $78,2 \rightarrow 12,4$ mmol/kg (perte 65,8 mmol/kg). La filière alactique a fourni $\approx$ 65,8 mmol ATP/kg (1 PCr $\rightarrow$ 1 ATP) en $\sim$ 10-15 premières secondes, puis s'est épuisée.
        \textbf{Arg 2 :} Explosion lactate : $5,1 \rightarrow 65,3$ mmol/kg (gain 60,2 mmol/kg). La glycolyse a produit $\approx$ 60,2 mmol lactate/kg $\rightarrow$ 30,1 mmol ATP/kg (car 2 lactate = 2 ATP). Puissance lactique $\approx$ 1 mmol ATP/kg/s, typique anaérobie lactique.
    \item \textbf{Remontée PCr à T2 (45,6 vs 12,4) :} À 20 min, l'effort est sous-maximal (85\% $\text{VO}_{2\text{max}}$). La filière \textbf{aérobie} fournit l'essentiel de l'ATP. L'excédent d'ATP mitochondrial (via créatine kinase mitochondriale) permet de \textbf{rephosphoryler la créatine} en PCr pendant l'effort même (récupération intra-exercice). Preuve que l'aérobie « recharge » la batterie alactique.
    \item \textbf{Calcul ATP glycolyse T0-T1 :} $\Delta$[Lactate] $= 60,2$ mmol/kg. $2$ Lactate $\leftrightarrow 2$ ATP $\Rightarrow$ 1 Lactate $\leftrightarrow$ 1 ATP. ATP glycolytique $= \mathbf{60,2\ mmol/kg\ sec}$.
        \emph{Vérification :} ATP total consommé $\approx$ (PCr utilisée) + (ATP glycolytique) + (ATP aérobie). PCr $\approx 65,8$ ; Glycolyse $\approx 60,2$ ; Aérobie (estimation $\text{VO}_2$ 30 s $\approx$ 1,5 L $\rightarrow$ $\sim$ 20 mmol ATP/kg). Total $\approx 146$ mmol/kg. Cohérent pour 30 s max.
\end{enumerate}
\end{exoresolu}

\courssub{Régulation fine : comment la cellule « choisit » la filière ?}

\begin{methode}[Régulation métabolique de la resynthèse ATP (niveau Terminale D)]
\begin{enumerate}
    \item \textbf{Signal déclencheur :} Augmentation [ADP], [AMP], [P$_i$], [Ca$^{2+}$] cytosolique $\downarrow$ [ATP]/[ADP] ratio, $\downarrow$ potentiel de phosphorylation.
    \item \textbf{Activation immédiate (ms-s) :} \textbf{Créatine Kinase} (équilibre proche) $\rightarrow$ tamponne ATP via PCr. \textbf{Phosphofructokinase (PFK)} activée par AMP, ADP, P$_i$, $\text{NH}_4^+$ ; inhibée par ATP, citrate, $\text{H}^+$ (pH). \textbf{Phosphorylase a} (glycogène) activée par $\text{Ca}^{2+}$ (via phosphorylase kinase) et AMP.
    \item \textbf{Activation différée (min) :} \textbf{Ca$^{2+}$ mitochondrial} active déshydrogénases du Krebs (pyruvate, isocitrate, $\alpha$-kétoglutarate). \textbf{ADP/ATP mitochondrial} contrôle vitesse chaîne respiratoire (contrôle respiratoire). \textbf{Transcription (heures/jours) :} PGC-1$\alpha$ $\rightarrow$ biogenèse mitochondriale, isoformes lentes (entraînement endurance).
\end{enumerate}
\end{methode}

\begin{savaistu}[Le « déficit en oxygène » et la dette O$_2$ (Historique : Hill \& Lupton, 1923 ; Margaria, 1960)]
Lors d'un effort brutal, la $\text{VO}_2$ met $\sim$ 20-30 s à atteindre son palier (cinétique lente). L'écart entre l'énergie demandée et l'énergie fournie par l'aérobie est le \textbf{déficit en O$_2$}. Il est comblé par l'anaérobie (PCr + Glycolyse). Après l'effort, la $\text{VO}_2$ reste élevée : \textbf{dette O$_2$} (EPOC - \emph{Excess Post-exercise Oxygen Consumption}). Elle sert à : 1) Reconstituer PCr (70\% dette O$_2$ alactique), 2) Oxider le lactate (Cycle de Cori foie, oxydation muscle/cœur), 3) Restaurer $\text{O}_2$ myoglobine, 4) Thermorégulation, 5) Ventilation/cœur élevés. Au Cameroun, en altitude (Dschang, Bafoussam $\approx$ 1400 m), le déficit O$_2$ est majoré (hypoxie) $\rightarrow$ plus grande part anaérobie pour même intensité relative.
\end{savaistu}

\begin{attention}[Erreurs de rédaction à bannir pour le Bac]
\begin{itemize}
    \item \textbf{Confondre « voie » et « substrat » :} « La voie anaérobie utilise le glycogène » $\rightarrow$ \textbf{Faux}. La \emph{glycolyse anaérobie} utilise le glycogène/glucose. La voie alactique utilise la PCr.
    \item \textbf{Oublier le rendement :} Dire « L'aérobie produit beaucoup d'ATP » sans chiffre. \textbf{Exigez :} « $\approx$ 30-32 ATP/glucose vs 2 ATP/glucose en anaérobie lactique ».
    \item \textbf{Dire « L'anaérobie ne produit pas d'ATP » :} Faux, elle en produit (2 ATP/glucose), mais \textbf{peu et vite}.
    \item \textbf{Mélanger lactate et acide lactique :} À pH physiologique, c'est l'ion \textbf{lactate}$^-$ + $\text{H}^+$. L'acidose vient des $\text{H}^+$ (hydrolyse ATP + glycolyse), pas de l'anion lactate qui est un carburant (navette lactate).
    \item \textbf{Négliger les lipides :} En effort long ($> 1$ h), l'oxydation des \textbf{acides gras} (beta-oxydation $\rightarrow$ Acétyl-CoA $\rightarrow$ Krebs) devient majoritaire. Rendement : $\approx$ 106 ATP / palmitate (C16).
\end{itemize}
\end{attention}

\begin{aretenir}[Synthèse pour le Bac : Les 3 piliers de la restauration ATP]
\begin{center}
\begin{tabular}{|l|c|c|c|}\hline
\textbf{Critère} & \textbf{Anaérobie Alactique} & \textbf{Anaérobie Lactique} & \textbf{Aérobie} \\ \hline
\textbf{Réaction clé} & $\ce{PCr + ADP -> Cr + ATP}$ & $\ce{Glucose -> 2 Lactate + 2 ATP}$ & $\ce{C6H12O6 + 6 O2 -> 6 CO2 + 6 H2O}$ \\ \hline
\textbf{Localisation} & Cytoplasme (myofibrilles) & Cytoplasme & Mitochondries \\ \hline
\textbf{O$_2$ requis} & Non & Non & \textbf{Oui} \\ \hline
\textbf{Vitesse (Puissance)} & \textbf{Très élevée} ($\sim$ 4-5) & Élevée ($\sim$ 2-3) & Modérée ($\sim$ 1) \\ \hline
\textbf{Capacité (Durée)} & Très faible (\SI{5-10}{s}) & Moyenne (\SI{1-3}{min}) & \textbf{Quasi illimitée} (heures) \\ \hline
\textbf{Rendement ATP/Glc} & 0 (transfert direct) & 2 ATP (net) & \textbf{30-32 ATP} \\ \hline
\textbf{Déchet} & Aucun (Cr recyclée) & \textbf{Lactate + H$^+$} (acidose) & $\text{CO}_2$ (gaz) + $\text{H}_2\text{O}$ \\ \hline
\textbf{Effort type} & Sprint 100 m, saut, lancer & 200-800 m, combat, intervalles & Marathon, cyclisme, marche, récupération \\ \hline
\end{tabular}
\end{center}
\textbf{Mots-clés pour dissertation :} \cle{Turnover ATP}, \cle{Homéostasie énergétique}, \cle{Créatine kinase}, \cle{Phosphofructokinase (PFK)}, \cle{Seuil lactique}, \cle{Déficit/Dette O$_2$}, \cle{Filière dominante}.
\end{aretenir}

\begin{exercice}[Entraînement Bac : QCM et Question courte]
\textbf{1.} Lors d'un effort intense de 30 secondes, la concentration musculaire en ATP diminue de 10\% tandis que celle en phosphocréatine chute de 80\%. Laquelle des affirmations suivantes est correcte ?
A. L'ATP n'est pas utilisée pendant l'effort.
B. L'ATP est resynthétisée aussi vite qu'elle est hydrolysée, principalement via la phosphocréatine.
C. La glycolyse est la seule voie active.
D. L'oxydation des lipides compense la perte d'ATP.

\textbf{2.} Un marathonien kényan court à 20 km/h (environ 85\% $\text{VO}_{2\text{max}}$). Après 1h30 de course, son glycogène musculaire est très bas. Quelle adaptation métabolique majeure se produit pour maintenir l'ATP ?
\textit{Indication : substrat, localisation, enzyme clé limitante.}

\textbf{3.} \emph{Vrai ou Faux ?} « La production de lactate est responsable de la fatigue musculaire car il acidifie le milieu et inhibe les enzymes. » Justifiez en deux lignes (distinction lactate/$\text{H}^+$, rôle tampon, navette lactate).
\end{exercice}

\begin{corrige}[Exercice n° Entraînement Bac]
\textbf{1. Réponse B.} L'ATP est maintenue quasi constante (homéostasie). La chute de PCr prouve que la réaction $\ce{PCr + ADP -> ATP + Cr}$ (créatine kinase) est la source immédiate majeure pendant les premières secondes, relayée par la glycolyse.
\textbf{2.} \textbf{Oxydation des acides gras (lipolyse $\rightarrow$ beta-oxydation mitochondriale).} Substrat : triglycérides intramusculaires + acides gras plasmatiques (libérés par lipoprotéine lipase). Enzyme limitante : \textbf{Carnitine Palmitoyl Transférase I (CPT-I)}, inhibée par le malonyl-CoA (marqueur statut glucidique). Baisse glycogène/insuline $\rightarrow$ baisse malonyl-CoA $\rightarrow$ activation CPT-I $\rightarrow$ entrée acyles-CoA dans mitochondrie. Rendement élevé mais puissance plus faible $\rightarrow$ oblige à baisser l'allure (« hitting the wall »).
\textbf{3.} \textbf{Faux (formulation piège).} C'est l'accumulation d'$\text{H}^+$ (protons) qui acidifie (pH $\downarrow$), pas l'anion lactate$^-$. L'hydrolyse de l'ATP ($\ce{ATP -> ADP + Pi + H+}$) est la source majeure d'$\text{H}^+$. Le lactate$^-$ est un \textbf{carburant} (oxydé in situ ou exporté vers foie/cœur via transporteurs MCT). L'acidose inhibe PFK et la fixation $\text{Ca}^{2+}$ sur la troponine C.
\end{corrige}

\coursec{V. Adaptation du métabolisme au type d'effort}

Dans la section précédente, nous avons établi que la cellule musculaire dispose de \cle{trois voies de restauration de l'ATP} : la filière \cle{anaérobie alactique} (phosphocréatine), la filière \cle{anaérobie lactique} (glycolyse produisant du lactate) et la filière \cle{aérobie} (respiration cellulaire). Une question surgit alors naturellement : ces trois filières fonctionnent-elles toutes de la même façon, quel que soit l'effort ? L'observation du sportif camerounais donne un premier indice : le sprinteur qui court le \SI{100}{m} au stade de Bépanda n'a pas le même gabarit, ni la même stratégie de course, que le marathonien qui gravit les pentes du Mont Cameroun lors de la \emph{Course de l'Espoir}. Le premier explose en moins de \SI{10}{s} puis s'arrête, essoufflé mais sans douleur durable ; le second court plus de \SI{2}{h} à allure régulière. Cette différence de comportement traduit une différence profonde de \cle{métabolisme} : chaque type d'effort sollicite préférentiellement une filière énergétique et un type de fibres musculaires. C'est ce que nous allons démontrer dans cette section.

\courssub{V.1. Les efforts brefs et violents : la dominance de la filière anaérobie alactique}

\textbf{Observation.} Un haltérophile soulève une barre de \SI{150}{kg} en \SI{2}{s} ; un sprinteur parcourt \SI{100}{m} en moins de \SI{11}{s}. Pendant ces efforts, l'athlète peut retenir sa respiration : l'effort est terminé avant que l'apport d'oxygène supplémentaire n'ait eu le temps d'atteindre les muscles (le délai de mise en route de la filière aérobie est de l'ordre de \SI{1}{min} à \SI{3}{min}). D'où vient alors l'énergie ?

\begin{experience}[Biopsies musculaires au cours d'un sprint maximal]
Des chercheurs ont prélevé par biopsie du vaste externe (muscle de la cuisse) des échantillons de muscle avant et immédiatement après un sprint maximal de \SI{10}{s}. Résultats moyens (en \SI{}{mmol/kg} de muscle sec) :
\begin{center}
\begin{tabularx}{\linewidth}{|l|X|X|}
\hline
\rowcolor{popblueL} \textbf{Paramètre} & \textbf{Avant l'effort} & \textbf{Après \SI{10}{s} de sprint} \\
\hline
ATP & environ \SI{25}{} & environ \SI{20}{} (baisse de \SI{20}{\%}) \\
\hline
Phosphocréatine (PC) & environ \SI{80}{} & environ \SI{15}{} (baisse de plus de \SI{80}{\%}) \\
\hline
Lactate musculaire & environ \SI{5}{} & environ \SI{25}{} (augmentation modérée) \\
\hline
\end{tabularx}
\end{center}
\textbf{Interprétation.} L'ATP ne s'effondre presque pas : il est immédiatement régénéré. En revanche, la phosphocréatine chute de plus de \SI{80}{\%} : c'est elle qui a fourni l'énergie, selon la réaction $\ce{PC + ADP -> Creatine + ATP}$, catalysée par la créatine kinase. Le lactate n'augmente que modérément : la glycolyse anaérobie commence à peine à contribuer à la fin de l'effort.
\textbf{Conclusion.} Lors d'un effort bref et violent (moins de \SI{10}{s}), la filière dominante est la filière \cle{anaérobie alactique}.
\end{experience}

\begin{definition}[Effort bref et violent]
Un effort est dit \cle{bref et violent} lorsqu'il dure moins de \SI{10}{s} environ et met en jeu une puissance maximale (sprint de \SI{100}{m}, lancer de poids, haltérophilie, saut en longueur). Il est couvert essentiellement par la filière anaérobie alactique, qui utilise les stocks de phosphocréatine du muscle, sans consommation d'oxygène et \textbf{sans production d'acide lactique}.
\end{definition}

Ces efforts sollicitent préférentiellement les \cle{fibres blanches} (fibres rapides, de type II), riches en phosphocréatine et en enzymes glycolytiques, capables de développer une grande puissance mais peu résistantes à la fatigue. C'est pourquoi les sprinteurs de haut niveau présentent une proportion élevée de fibres blanches dans leurs muscles, souvent supérieure à \SI{70}{\%}.

\begin{attention}[Anaérobie alactique $\neq$ anaérobie lactique]
Erreur classique au Baccalauréat : confondre les deux filières anaérobies. Retiens bien : seule la filière \cle{anaérobie lactique} produit de l'\cle{acide lactique} (lactate). La filière \cle{anaérobie alactique} (préfixe \emph{a-} = sans) régénère l'ATP à partir de la phosphocréatine, sans acide lactique. Écrire « le sprinteur produit beaucoup d'acide lactique pendant son \SI{100}{m} » est une erreur : la production massive de lactate caractérise les efforts de \SI{30}{s} à \SI{3}{min}.
\end{attention}

\courssub{V.2. Les efforts intenses et soutenus : la dominance de la filière anaérobie lactique}

\textbf{Situation problème.} Un athlète court le \SI{400}{m} en \SI{50}{s}. À l'arrivée, il ressent une vive « brûlure » dans les cuisses et peine à marcher. Pourtant, l'effort a duré bien plus de \SI{10}{s} : les stocks de phosphocréatine sont épuisés depuis longtemps. Quelle filière a pris le relais, et d'où vient la douleur ?

\begin{definition}[Effort intense et soutenu]
Un effort \cle{intense et soutenu} dure de \SI{10}{s} à \SI{3}{min} environ (\SI{400}{m}, \SI{800}{m}, \SI{100}{m} nage libre, combat de lutte). L'intensité est trop élevée pour que la filière aérobie, lente à se mettre en route et limitée en puissance, suffise à couvrir les besoins. La filière dominante est alors la filière \cle{anaérobie lactique} : la glycolyse dégrade le glucose (issu du glycogène musculaire) en pyruvate, puis en \cle{lactate}, avec un rendement faible (\SI{2}{ATP} par glucose) mais une puissance élevée.
\end{definition}

L'équation bilan de la glycolyse anaérobie s'écrit :
\[ \ce{C6H12O6 -> 2 CH3-CHOH-COOH} \qquad \text{avec production nette de \SI{2}{ATP}.} \]

Le lactate s'accumule dans la fibre musculaire puis diffuse dans le sang : la lactatémie peut passer de \SI{1}{mmol/L} au repos à plus de \SI{15}{mmol/L} à l'arrivée d'un \SI{400}{m}. Cette accumulation d'ions H$^+$ (acidose musculaire) inhibe les enzymes de la glycolyse et perturbe la formation des ponts actine-myosine : c'est une cause majeure de la \cle{fatigue} et de la sensation de brûlure.

\begin{propriete}[Origines de la fatigue selon le type d'effort — établie par biopsies musculaires]
Les biopsies musculaires réalisées chez des athlètes ont établi que la cause dominante de la fatigue dépend du type d'effort :
\begin{itemize}
\item effort \textbf{bref et violent} ($< \SI{10}{s}$) : fatigue liée à l'\cle{épuisement de la phosphocréatine} ;
\item effort \textbf{intense et soutenu} (\SI{10}{s} à \SI{3}{min}) : fatigue liée à l'\cle{accumulation de lactate} (acidose musculaire) ;
\item effort \textbf{prolongé} ($> \SI{3}{min}$) : fatigue liée à l'\cle{épuisement des réserves de glycogène}.
\end{itemize}
\end{propriete}

\courssub{V.3. Les efforts de longue durée : la dominance de la filière aérobie}

\textbf{Observation.} Le marathonien court \SI{42,195}{km} en plus de \SI{2}{h}. Un tel effort ne peut être couvert ni par la phosphocréatine (épuisée en quelques secondes) ni par la filière lactique (l'accumulation de lactate bloquerait le muscle en quelques minutes). Seule une filière de grande capacité, utilisable pendant des heures, convient : la filière \cle{aérobie}.

\begin{definition}[Effort de longue durée]
Un effort de \cle{longue durée} dépasse \SI{3}{min} (\SI{1500}{m}, marathon, cyclisme, match de football). Il est couvert majoritairement par la filière \cle{aérobie} : l'oxydation complète du glucose et des acides gras dans les mitochondries, selon le bilan :
\[ \ce{C6H12O6 + 6 O2 -> 6 CO2 + 6 H2O} \qquad \text{avec production d'environ \SI{36}{ATP} par glucose.} \]
Cette filière a une puissance limitée mais une capacité considérable : elle peut fonctionner des heures tant que l'oxygène et les substrats (glucose, acides gras) sont disponibles.
\end{definition}

Trois caractéristiques métaboliques accompagnent ce type d'effort :

\begin{itemize}
\item \textbf{Les substrats évoluent au cours du temps.} En début d'effort, le muscle consomme surtout le \cle{glycogène} musculaire et le glucose sanguin. Au-delà de \SI{30}{min} à \SI{60}{min}, la part des \cle{lipides} (acides gras issus du tissu adipeux) augmente progressivement et peut devenir majoritaire, ce qui ménage le glycogène.
\item \textbf{Les fibres rouges sont sollicitées.} Les fibres de type I (lentes), riches en mitochondries, en myoglobine et en capillaires sanguins, sont adaptées à l'oxydation : les marathoniens de haut niveau possèdent souvent plus de \SI{80}{\%} de fibres rouges dans les muscles des membres inférieurs.
\item \textbf{La ventilation et la circulation s'adaptent.} La fréquence cardiaque passe d'environ \SI{60}{battements/min} au repos à \SIrange{160}{180}{battements/min} à l'effort ; le débit cardiaque peut être multiplié par 5 (de \SI{5}{L/min} à \SI{25}{L/min}) ; la ventilation passe d'environ \SI{6}{L/min} à plus de \SI{100}{L/min}. Ces adaptations assurent l'apport d'O$_2$ et de substrats aux muscles et l'élimination du CO$_2$.
\end{itemize}

\begin{exemplebox}[Le marathonien, une machine aérobie]
Un marathonien de niveau élite puise \textbf{plus de \SI{95}{\%} de son énergie} dans la filière aérobie. Pendant la course, il peut consommer \textbf{plus de \SI{500}{g} de glycogène}, soit la quasi-totalité de ses réserves musculaires et hépatiques. Sa consommation d'oxygène atteint \SIrange{4}{5}{L/min}, contre \SI{0,3}{L/min} au repos.
\end{exemplebox}

\courssub{V.4. Interprétation des données expérimentales : courbes et tableaux}

\textbf{Évolution des paramètres physiologiques en fonction de l'intensité de l'effort.} Lorsqu'on fait réaliser à un sujet un exercice dont la puissance augmente par paliers (test d'effort sur ergocycle), on mesure la fréquence cardiaque et la concentration sanguine en lactate. On obtient les courbes suivantes :

\begin{popfigure}
\begin{center}
\begin{tikzpicture}
\begin{axis}[
width=0.9\linewidth, height=7cm,
xlabel={Puissance de l'effort (W)},
ylabel={Fréquence cardiaque (battements/min)},
axis y line*=left, axis x line*=bottom,
xmin=0, xmax=320, ymin=50, ymax=200,
legend style={at={(0.03,0.97)}, anchor=north west, font=\small},
every axis plot/.append style={thick}
]
\addplot[popcurve, domain=0:300, samples=100] {60 + 0.42*x};
\addlegendentry{Fréquence cardiaque}
\addplot[poporange, domain=0:300, samples=100] {1 + 0.004*x + 0.00012*x^2};
\addlegendentry{Lactate sanguin (mmol/L)}
\end{axis}
\end{tikzpicture}
\end{center}
\legende{Évolution de la fréquence cardiaque (bleu) et de la concentration sanguine en lactate (orange) en fonction de la puissance de l'effort. La fréquence cardiaque augmente de façon quasi proportionnelle à la puissance. Le lactate reste stable (environ \SI{1}{mmol/L} à \SI{2}{mmol/L}) jusqu'à environ \SI{150}{W}, puis s'élève de plus en plus vite : c'est le \textbf{seuil d'accumulation du lactate} (ou seuil lactique), qui marque le passage d'un métabolisme essentiellement aérobie à un métabolisme de plus en plus anaérobie lactique.}
\end{popfigure}

\textbf{Interprétation.} En dessous du seuil, le lactate produit est oxydé au fur et à mesure : la filière aérobie suffit. Au-dessus du seuil, la production de lactate dépasse les capacités d'oxydation : il s'accumule dans le sang. Ce seuil est un paramètre essentiel de l'entraînement : un athlète d'endurance s'entraîne juste en dessous de son seuil pour pouvoir courir longtemps sans accumulation de lactate.

\textbf{Contribution des trois filières selon la discipline.} Le tableau et le diagramme ci-dessous rassemblent des données classiques d'exigences métaboliques de différents sports (valeurs approchées, en \% de l'énergie totale fournie) :

\begin{center}
\begin{tabularx}{\linewidth}{|l|X|X|X|}
\hline
\rowcolor{popblueL} \textbf{Discipline} & \textbf{Anaérobie alactique (\%)} & \textbf{Anaérobie lactique (\%)} & \textbf{Aérobie (\%)} \\
\hline
\SI{100}{m} (sprint) & 90 & 8 & 2 \\
\hline
\SI{400}{m} & 30 & 55 & 15 \\
\hline
\SI{1500}{m} & 10 & 25 & 65 \\
\hline
Marathon & 0 & 2 & 98 \\
\hline
Football (match) & 15 & 15 & 70 \\
\hline
\end{tabularx}
\end{center}

\begin{popfigure}
\begin{center}
\begin{tikzpicture}
\begin{axis}[
ybar stacked, bar width=0.55cm,
width=0.9\linewidth, height=7.5cm,
ylabel={Contribution (\% de l'énergie totale)},
symbolic x coords={100 m, 400 m, 1500 m, Marathon},
xtick=data, ymin=0, ymax=105,
legend style={at={(1.02,0.5)}, anchor=west, font=\small},
every axis plot/.append style={thick}
]
\addplot[fill=poporange] coordinates {(100 m,90) (400 m,30) (1500 m,10) (Marathon,0)};
\addplot[fill=poppink] coordinates {(100 m,8) (400 m,55) (1500 m,25) (Marathon,2)};
\addplot[fill=popgreen] coordinates {(100 m,2) (400 m,15) (1500 m,65) (Marathon,98)};
\legend{Anaérobie alactique, Anaérobie lactique, Aérobie}
\end{axis}
\end{tikzpicture}
\end{center}
\legende{Diagramme à bandes empilées montrant la contribution approximative des trois filières énergétiques en fonction de la durée de l'effort. Plus l'effort est bref et intense, plus les filières anaérobies dominent ; plus il est long, plus la filière aérobie domine. Les trois filières fonctionnent toujours simultanément, mais dans des proportions différentes.}
\end{popfigure}

\textbf{Interprétation.} On observe une transition continue : la filière alactique domine sous \SI{10}{s}, la filière lactique domine entre \SI{30}{s} et \SI{2}{min}, la filière aérobie domine au-delà de \SI{3}{min}. Aucune filière ne fonctionne seule : même au \SI{100}{m}, la filière aérobie contribue faiblement, et même au marathon, un sprint final fait appel aux filières anaérobies.

\textbf{Évolution du glycogène au cours d'un effort prolongé.} Les biopsies réalisées au cours d'un effort d'endurance montrent une décroissance régulière du glycogène musculaire :

\begin{popfigure}
\begin{center}
\begin{tikzpicture}
\begin{axis}[
width=0.9\linewidth, height=7cm,
xlabel={Durée de l'effort (min)},
ylabel={Glycogène musculaire (mmol/kg de muscle sec)},
xmin=0, xmax=130, ymin=0, ymax=110,
every axis plot/.append style={very thick}
]
\addplot[popcurve, domain=0:120, samples=100] {100*exp(-x/55)};
\addplot[poporange, dashed, domain=0:120] {20};
\draw[popdark, dashed] (axis cs:100,0) -- (axis cs:100,20);
\node[popdark, font=\small, align=center] at (axis cs:100,55) {« mur » : épuisement\\du glycogène};
\end{axis}
\end{tikzpicture}
\end{center}
\legende{Évolution de la teneur en glycogène du muscle au cours d'un effort de longue durée à intensité constante (environ \SI{75}{\%} de la puissance maximale aérobie). Le glycogène décroît régulièrement ; lorsqu'il descend sous un seuil critique (environ \SI{20}{mmol/kg}, ici vers \SI{100}{min}), le sujet ressent une fatigue brutale et intense : c'est le fameux « mur » du marathonien, qui survient classiquement vers le 30\textsuperscript{e} kilomètre.}
\end{popfigure}

\begin{methode}[Analyser un document sur les exigences métaboliques d'un sport]
Face à un tableau, une courbe ou un texte portant sur un sport, procède en quatre étapes :
\begin{enumerate}
\item \textbf{Repère la durée de l'effort} : moins de \SI{10}{s} ? entre \SI{10}{s} et \SI{3}{min} ? plus de \SI{3}{min} ?
\item \textbf{Repère l'intensité} : maximale et discontinue, ou sous-maximale et continue ?
\item \textbf{Déduis la filière dominante} : bref et violent $\rightarrow$ anaérobie alactique ; intense et soutenu $\rightarrow$ anaérobie lactique ; long $\rightarrow$ aérobie.
\item \textbf{Déduis le type de fibres sollicité et la cause de la fatigue} : fibres blanches et épuisement de la PC ; accumulation de lactate ; fibres rouges et épuisement du glycogène.
\end{enumerate}
Vérifie toujours la cohérence avec les données chiffrées du document (lactatémie élevée = filière lactique ; PC effondrée = filière alactique).
\end{methode}

\courssub{V.5. Évolution des paramètres physiologiques au cours de l'exercice et récupération}

\textbf{Au cours d'un exercice prolongé d'intensité modérée}, les biopsies et prélèvements sanguins montrent une évolution caractéristique :

\begin{itemize}
\item l'\cle{ATP} musculaire reste \textbf{quasi stable} : il est régénéré en permanence ; sa chute serait incompatible avec la vie de la cellule ;
\item la \cle{phosphocréatine} \textbf{s'effondre dans les premières secondes} (baisse de \SIrange{50}{80}{\%}), puis se stabilise à un niveau bas tant que dure l'effort ;
\item le \cle{glycogène} \textbf{décroît régulièrement} jusqu'à l'épuisement ;
\item le \cle{lactate} sanguin \textbf{varie selon l'intensité} : il reste stable et bas si l'effort est sous le seuil, il s'accumule si l'effort est intense.
\end{itemize}

\textbf{La récupération} est la période qui suit l'effort et pendant laquelle l'organisme restaure son état initial. Sa durée dépend du paramètre considéré :

\begin{itemize}
\item la \textbf{reconstitution des stocks de phosphocréatine} est \textbf{rapide} : elle est complète en \SIrange{3}{5}{min} de repos, grâce à l'ATP produit par la filière aérobie (c'est pourquoi le sprinteur peut répéter des sprints après quelques minutes de récupération) ;
\item l'\textbf{élimination du lactate} demande \SIrange{30}{120}{min} : le lactate est oxydé en CO$_2$ et H$_2$O par le muscle cardiaque et les fibres rouges, ou reconverti en glucose par le foie (cycle de Cori) — c'est l'une des raisons de la consommation d'oxygène élevée qui persiste après l'effort (\cle{dette d'oxygène}) ;
\item la \textbf{reconstitution des stocks de glycogène} est \textbf{lente} : elle exige \SIrange{24}{48}{h} et une alimentation riche en glucides (ignames, riz, manioc), ce qui explique l'importance de la nutrition dans les jours qui suivent une compétition d'endurance.
\end{itemize}

\begin{aretenir}[L'essentiel de la section]
\begin{itemize}
\item Effort bref et violent ($< \SI{10}{s}$) : filière \textbf{anaérobie alactique}, fibres blanches, fatigue par épuisement de la phosphocréatine.
\item Effort intense et soutenu (\SI{10}{s} à \SI{3}{min}) : filière \textbf{anaérobie lactique}, fatigue par accumulation de lactate.
\item Effort de longue durée ($> \SI{3}{min}$) : filière \textbf{aérobie}, fibres rouges, consommation de glycogène puis de lipides, fatigue par épuisement du glycogène.
\item Les trois filières fonctionnent toujours ensemble, mais l'une domine selon la durée et l'intensité de l'effort.
\item Récupération : PC en quelques minutes, lactate en \SIrange{1}{2}{h}, glycogène en \SIrange{24}{48}{h}.
\end{itemize}
\end{aretenir}

\begin{savaistu}[Le football, un sport au profil métabolique mixte]
Au football, un joueur comme ceux des Lions Indomptables répète des sprints brefs et explosifs (\SIrange{10}{30}{m}, filière \textbf{anaérobie alactique}) entrecoupés de tacles et de duels (filière \textbf{lactique}), le tout sur un fond d'effort continu de \SI{90}{min} pendant lequel il parcourt \SIrange{10}{12}{km} (filière \textbf{aérobie}, environ \SI{70}{\%} de l'énergie totale). Le profil métabolique du football est donc \textbf{mixte} : l'entraînement doit développer les trois filières, et la récupération entre les sprints (reconstitution de la PC en quelques minutes) est un facteur décisif de performance en fin de match.
\end{savaistu}

\begin{exoresolu}[Analyse des exigences métaboliques de deux disciplines]
\textbf{Énoncé.} Des biopsies du muscle quadriceps ont été réalisées chez deux athlètes, A et B, immédiatement après leur épreuve respective. Athlète A (épreuve de \SI{45}{s}) : phosphocréatine effondrée à \SI{10}{\%} de sa valeur initiale, lactate musculaire multiplié par 12, glycogène diminué de \SI{15}{\%}. Athlète B (épreuve de \SI{2}{h} \SI{10}{min}) : lactate musculaire proche de la valeur de repos, glycogène diminué de \SI{90}{\%}. Identifie la discipline probable de chaque athlète, la filière dominante et la cause de sa fatigue.
\tcblower
\textbf{Solution détaillée.}
\textbf{Athlète A.} L'épreuve dure \SI{45}{s}, donc entre \SI{10}{s} et \SI{3}{min} : c'est un effort intense et soutenu (probablement un \SI{400}{m}). La phosphocréatine est effondrée (elle a contribué au début de l'effort), mais surtout le lactate musculaire est multiplié par 12 : la filière dominante est la filière \textbf{anaérobie lactique}. Le glycogène n'a diminué que de \SI{15}{\%} : il n'est pas limitant. La fatigue est due à l'\textbf{accumulation de lactate} (acidose musculaire), qui inhibe la glycolyse et la contraction.

\textbf{Athlète B.} L'épreuve dure \SI{2}{h} \SI{10}{min} : c'est un effort de longue durée (probablement un marathon). Le lactate est resté bas : l'athlète est resté sous son seuil d'accumulation, la filière dominante est la filière \textbf{aérobie}. Le glycogène a diminué de \SI{90}{\%} : la fatigue est due à l'\textbf{épuisement des réserves de glycogène} (le « mur » du marathonien). La récupération complète du glycogène exigera \SIrange{24}{48}{h} et une alimentation glucidique abondante.
\end{exoresolu}

\begin{exercice}[Pour t'entraîner]
Un footballeur effectue, au cours d'un match, environ 40 sprints de \SIrange{15}{30}{m} et parcourt au total \SI{11}{km} en \SI{90}{min}. Sa lactatémie mesurée à la mi-temps est de \SI{4}{mmol/L}.
\begin{enumerate}
\item Quelle filière fournit l'essentiel de l'énergie sur l'ensemble du match ? Justifie.
\item Quelle filière est sollicitée lors de chaque sprint ? Pourquoi le joueur peut-il répéter les sprints tout au long du match ?
\item Interprète la valeur de la lactatémie à la mi-temps.
\end{enumerate}
\end{exercice}

\begin{corrige}[Exercice : le profil métabolique du footballeur]
\begin{enumerate}
\item La filière \textbf{aérobie} fournit l'essentiel de l'énergie (environ \SI{70}{\%}), car le match dure \SI{90}{min}, largement plus de \SI{3}{min}, et l'intensité moyenne est sous-maximale.
\item Chaque sprint dure moins de \SI{10}{s} : c'est la filière \textbf{anaérobie alactique} (phosphocréatine) qui domine. Le joueur peut répéter les sprints car les stocks de PC se reconstituent en \SIrange{3}{5}{min} pendant les phases de jeu de faible intensité, grâce à l'ATP fourni par la filière aérobie.
\item Une lactatémie de \SI{4}{mmol/L} est supérieure à la valeur de repos (\SI{1}{mmol/L}) mais reste modérée : elle traduit une production de lactate lors des sprints et des duels (filière lactique), largement compensée par son oxydation pendant les phases aérobies. Le joueur évolue globalement autour de son seuil lactique.
\end{enumerate}
\end{corrige}

\newpage

\coursec{Activité d'intégration}

\courssub{Objectifs de l'activité}

À l'issue de cette activité d'intégration, tu devras être capable de :

\begin{itemize}
\item interpréter les données fournies par des biopsies musculaires (taux d'\cle{ATP}, de \cle{phosphocréatine} et de \cle{glycogène}) avant et après des exercices de durées différentes ;
\item analyser l'évolution de paramètres physiologiques (\cle{lactate} sanguin, consommation d'\ce{O2}) au cours de l'effort ;
\item identifier la \cle{filière énergétique} dominante (anaérobie alactique, anaérobie lactique, aérobie) en fonction du type d'effort ;
\item relier le \cle{profil en fibres musculaires} (fibres I et II) d'un athlète à sa spécialité sportive ;
\item rédiger une conclusion argumentée expliquant la nécessité du renouvellement de l'ATP et l'adaptation du métabolisme à chaque type d'effort.
\end{itemize}

\courssub{Situation}

\textbf{De la biopsie au podium : quel métabolisme pour quel champion ?}

Lors des Jeux universitaires de Yaoundé, l'équipe de médecine du sport du Centre Hospitalier Universitaire (CHU) de Yaoundé souhaite préparer une conférence de sensibilisation destinée aux jeunes athlètes : « Pourquoi le muscle a-t-il sans cesse besoin de renouveler son ATP, et comment s'adapte-t-il à l'effort ? ». Tu fais partie du groupe d'élèves de Terminale D invité à exploiter les données expérimentales recueillies.

\textbf{Document 1 : biopsies du vaste externe (muscle de la cuisse) chez un même sujet entraîné.}

Les prélèvements sont réalisés au repos puis immédiatement après trois exercices : un sprint de \SI{10}{s}, une course intense de \SI{2}{min} et une course d'endurance de \SI{60}{min}. Les teneurs sont exprimées en mmol par kg de muscle sec.

\begin{center}
\begin{tabularx}{\linewidth}{|l|X|X|X|X|X|}
\hline
\rowcolor{popblueL} \textbf{Paramètre} & \textbf{Repos} & \textbf{Après sprint 10 s} & \textbf{Après course 2 min} & \textbf{Après course 60 min} \\
\hline
ATP (mmol/kg) & 25 & 22 & 21 & 23 \\
\hline
Phosphocréatine (mmol/kg) & 80 & 20 & 35 & 60 \\
\hline
Glycogène (mmol/kg) & 400 & 385 & 260 & 80 \\
\hline
\end{tabularx}
\end{center}

\textbf{Document 2 : paramètres physiologiques mesurés pendant les trois exercices.}

\begin{center}
\begin{tabularx}{\linewidth}{|l|X|X|X|}
\hline
\rowcolor{popblueL} \textbf{Paramètre} & \textbf{Sprint 10 s} & \textbf{Course 2 min} & \textbf{Course 60 min} \\
\hline
Lactate sanguin maximal (mmol/L) & 4 & 18 & 5 \\
\hline
Consommation d'\ce{O2} (en \% de la \ce{VO2}max) & faible ($<$ 40 \%) & élevée (90 \%) mais insuffisante & élevée et stable (75 \%) \\
\hline
\end{tabularx}
\end{center}

\textbf{Document 3 : biopsies histochimiques de deux athlètes camerounais.}

La coloration des ATPases myofibrillaires permet de distinguer les fibres de type I (lentes, oxydatives) des fibres de type II (rapides, glycolytiques) dans le vaste externe.

\begin{center}
\begin{tabularx}{\linewidth}{|l|X|X|}
\hline
\rowcolor{popblueL} \textbf{Athlète} & \textbf{Fibres de type I (\%)} & \textbf{Fibres de type II (\%)} \\
\hline
Athlète A & 25 & 75 \\
\hline
Athlète B & 80 & 20 \\
\hline
\end{tabularx}
\end{center}

L'un de ces athlètes est spécialiste du \SI{100}{m}, l'autre du marathon. Le médecin du CHU te demande de trancher.

\courssub{Consignes}

Travail en groupes de quatre élèves, suivi d'une restitution orale avec schéma d'interprétation au tableau.

\begin{enumerate}
\item Pour chacun des trois exercices, décris et interprète les variations des taux d'ATP, de phosphocréatine et de glycogène (document 1). Que remarques-tu concernant le taux d'ATP ?
\item En t'aidant du document 2, identifie la filière énergétique dominante pour chaque exercice. Justifie chaque réponse par au moins deux arguments chiffrés.
\item Explique pourquoi le taux de lactate sanguin est très élevé après la course de \SI{2}{min} mais faible après le sprint de \SI{10}{s}, alors que les deux efforts sont violents.
\item Attribue à chaque athlète (document 3) sa spécialité (100 m ou marathon) en justifiant par les propriétés des fibres de type I et de type II.
\item Propose un schéma d'interprétation récapitulatif reliant, pour chaque type d'effort, la filière mobilisée, les substrats consommés et le type de fibres préférentiellement recruté.
\item Rédige une conclusion d'une quinzaine de lignes répondant à la question du médecin : pourquoi le renouvellement de l'ATP est-il indispensable au cours de l'exercice musculaire, et comment le métabolisme s'adapte-t-il à chaque type d'effort ?
\end{enumerate}

\courssub{Ressources à mobiliser}

\begin{aretenir}[Ce qu'il faut avoir en tête]
\begin{itemize}
\item La contraction musculaire consomme de l'\cle{ATP} (glissement des filaments d'actine sur la myosine) ; les stocks d'ATP musculaire sont très faibles et ne couvrent que quelques secondes d'effort : l'ATP doit donc être \textbf{renouvelée en permanence}.
\item \textbf{Filière anaérobie alactique} : \ce{PCr + ADP -> Cr + ATP} (enzyme : créatine kinase). Très rapide, sans oxygène, sans lactate, mais épuisable en une dizaine de secondes. Dominante pour les efforts \textbf{brefs et violents}.
\item \textbf{Filière anaérobie lactique} : glycolyse du glycogène, \ce{C6H12O6 -> 2 lactate + 2 ATP}. Rapide, sans oxygène, produit du \cle{lactate}. Dominante pour les efforts \textbf{intenses et soutenus} (30 s à 3 min environ).
\item \textbf{Filière aérobie} : oxydation complète du glucose et des lipides dans les mitochondries, \ce{C6H12O6 + 6O2 -> 6CO2 + 6H2O + 36 ATP}. Lente à se mettre en place, nécessite \ce{O2}, très rentable. Dominante pour les efforts de \textbf{longue durée}.
\item \textbf{Fibres de type I} (lentes, rouges) : riches en mitochondries et myoglobine, métabolisme oxydatif, résistantes à la fatigue. \textbf{Fibres de type II} (rapides, blanches) : riches en glycogène et enzymes glycolytiques, puissantes mais fatigables.
\end{itemize}
\end{aretenir}

\courssub{Démarche et corrigé commenté}

\begin{exoresolu}[De la biopsie au podium : corrigé détaillé]
Mobiliser les trois documents pour répondre aux six consignes.
\tcblower
\textbf{Étape 1 : interprétation des biopsies (document 1).}

\begin{itemize}
\item \textbf{Sprint 10 s} : la phosphocréatine chute brutalement de 80 à \SI{20}{mmol/kg} ($-75\,\%$), le glycogène ne diminue presque pas (400 à 385) et l'ATP reste quasi stable (25 à 22). Le sujet a puisé dans son stock de phosphocréatine pour régénérer l'ATP : c'est la signature de la filière anaérobie alactique.
\item \textbf{Course 2 min} : le glycogène chute fortement (400 à \SI{260}{mmol/kg}, soit $-35\,\%$), la phosphocréatine est aussi entamée (35). Le muscle a massivement dégradé son glycogène par glycolyse.
\item \textbf{Course 60 min} : le glycogène est presque épuisé (\SI{80}{mmol/kg}, $-80\,\%$), la phosphocréatine est partiellement reconstituée (60) car l'effort modéré permet sa resynthèse. L'oxydation aérobie a dominé.
\item \textbf{Point capital} : dans les trois cas, le taux d'ATP ne varie que faiblement (21 à 25). L'ATP est consommée dès qu'elle est produite : sa concentration reste stable parce que sa \textbf{vitesse de régénération égale sa vitesse d'utilisation}. C'est la preuve expérimentale de la nécessité du renouvellement permanent de l'ATP.
\end{itemize}

\textbf{Étape 2 : identification des filières dominantes (documents 1 et 2).}

\begin{center}
\begin{tabularx}{\linewidth}{|l|X|X|}
\hline
\rowcolor{popblueL} \textbf{Exercice} & \textbf{Filière dominante} & \textbf{Arguments} \\
\hline
Sprint 10 s & Anaérobie alactique & PCr effondrée (80 à 20) ; lactate faible (4 mmol/L) ; \ce{O2} peu consommé \\
\hline
Course 2 min & Anaérobie lactique & Glycogène fortement entamé ; lactate très élevé (18 mmol/L) ; \ce{VO2} insuffisante malgré 90 \% de \ce{VO2}max \\
\hline
Course 60 min & Aérobie & Glycogène presque épuisé ; lactate faible (5 mmol/L) ; consommation d'\ce{O2} élevée et stable (75 \%) \\
\hline
\end{tabularx}
\end{center}

\textbf{Étape 3 : le paradoxe du lactate.}

Après le sprint, le lactate est faible car l'effort de 10 s est couvert presque exclusivement par la phosphocréatine, voie qui ne produit \textbf{pas} de lactate : la glycolyse n'a pas le temps de s'installer pleinement. Après 2 min d'effort intense, la demande en ATP dépasse la capacité de la filière aérobie (l'apport d'\ce{O2} est limité par le débit cardiaque) : la glycolyse anaérobie tourne à plein régime et accumule le lactate (18 mmol/L contre environ 1 mmol/L au repos).

\textbf{Étape 4 : attribution des profils de fibres.}

L'\textbf{athlète A} (75 \% de fibres II) est le \textbf{sprinteur du 100 m} : ses fibres rapides glycolytiques développent une puissance élevée sur un effort bref. L'\textbf{athlète B} (80 \% de fibres I) est le \textbf{marathonien} : ses fibres lentes oxydatives, riches en mitochondries, résistent à la fatigue sur un effort prolongé.

\textbf{Étape 5 : schéma d'interprétation.}

\begin{popfigure}
\begin{center}
\begin{tikzpicture}[font=\small, >=Stealth]
\node[draw, rounded corners, fill=poporangeL, text width=3.1cm, align=center] (e1) at (0,3) {\textbf{Effort bref et violent} (10 s)};
\node[draw, rounded corners, fill=poppinkL, text width=3.1cm, align=center] (e2) at (0,0) {\textbf{Effort intense soutenu} (2 min)};
\node[draw, rounded corners, fill=popgreenL, text width=3.1cm, align=center] (e3) at (0,-3) {\textbf{Effort longue durée} (60 min)};
\node[draw, rounded corners, fill=poporange!40, text width=3.4cm, align=center] (f1) at (5.2,3) {Anaérobie alactique\\ \ce{PCr + ADP -> ATP}};
\node[draw, rounded corners, fill=poppink!40, text width=3.4cm, align=center] (f2) at (5.2,0) {Anaérobie lactique\\ glycogène $\to$ lactate};
\node[draw, rounded corners, fill=popgreen!40, text width=3.4cm, align=center] (f3) at (5.2,-3) {Aérobie\\ glucose + lipides + \ce{O2}};
\node[draw, rounded corners, fill=popblueL, text width=2.8cm, align=center] (t1) at (10.2,3) {Fibres II\\ (sprinteur)};
\node[draw, rounded corners, fill=popblueL, text width=2.8cm, align=center] (t2) at (10.2,0) {Fibres II puis I};
\node[draw, rounded corners, fill=popblueL, text width=2.8cm, align=center] (t3) at (10.2,-3) {Fibres I\\ (marathonien)};
\draw[->, thick, popdark] (e1) -- (f1); \draw[->, thick, popdark] (f1) -- (t1);
\draw[->, thick, popdark] (e2) -- (f2); \draw[->, thick, popdark] (f2) -- (t2);
\draw[->, thick, popdark] (e3) -- (f3); \draw[->, thick, popdark] (f3) -- (t3);
\end{tikzpicture}
\end{center}
\legende{Schéma d'interprétation : à chaque type d'effort correspondent une filière énergétique dominante et un recrutement préférentiel de fibres musculaires.}
\end{popfigure}

\textbf{Étape 6 : conclusion rédigée (modèle).}

La contraction musculaire est assurée par l'hydrolyse de l'ATP, dont les stocks ne couvrent que quelques secondes d'effort : sans renouvellement permanent, tout mouvement s'arrêterait presque aussitôt. La stabilité du taux d'ATP observée dans les biopsies montre que le muscle régénère l'ATP aussi vite qu'il la consomme. Pour cela, il adapte son métabolisme au type d'effort : la filière anaérobie alactique (phosphocréatine) fournit instantanément l'ATP des efforts brefs et violents ; la filière anaérobie lactique (glycolyse du glycogène) prend le relais des efforts intenses et soutenus, au prix d'une accumulation de lactate ; enfin la filière aérobie, lente mais très rentable, oxyde glucose et lipides pour soutenir les efforts de longue durée. Cette spécialisation métabolique se double d'une spécialisation structurale : les fibres II dominent chez le sprinteur, les fibres I chez le marathonien. Couvrir les besoins énergétiques de l'activité physique, c'est donc choisir la bonne filière au bon moment.
\end{exoresolu}

\begin{attention}[Pièges fréquents]
\begin{itemize}
\item Ne dis jamais que « le taux d'ATP diminue fortement pendant l'effort » : il reste quasi constant, et c'est justement cette stabilité qui prouve le renouvellement.
\item Le sprint ne produit \textbf{pas} de lactate en quantité notable : n'attribue pas la filière lactique à un effort de 10 s.
\item Une filière « dominante » n'est jamais exclusive : les trois voies fonctionnent toujours simultanément, avec des contributions variables.
\item Les fibres de type I sont dites « lentes » car leur vitesse de contraction est faible, pas parce qu'elles seraient inactives.
\end{itemize}
\end{attention}

\courssub{Bilan de l'activité}

Cette activité t'a permis de mettre en évidence, à partir de données authentiques de biopsies et de mesures physiologiques, trois idées fondamentales de la leçon :

\begin{enumerate}
\item \textbf{Le renouvellement de l'ATP est une nécessité absolue} : les stocks étant dérisoires, la concentration d'ATP ne reste stable que parce que sa régénération compense exactement sa consommation.
\item \textbf{Le métabolisme musculaire s'adapte au type d'effort} : anaérobie alactique pour l'effort bref et violent, anaérobie lactique pour l'effort intense et soutenu, aérobie pour l'effort de longue durée, chaque filière laissant une « signature » mesurable (phosphocréatine, lactate, \ce{O2}, glycogène).
\item \textbf{La structure épouse la fonction} : la proportion de fibres I et II, en partie déterminée génétiquement et modulable par l'entraînement, oriente chaque athlète vers sa spécialité — du sprinteur au marathonien camerounais.
\end{enumerate}

Tu es désormais capable d'analyser et d'interpréter tout document relatif aux exigences métaboliques des différents sports, une compétence directement exigible au Baccalauréat.

\newpage

\coursec{Méthodes et exercices résolus}

\coursec{I. Structure et composition chimique de la fibre musculaire striée squelettique}

\begin{methode}[Structure ultrastructurale et composition chimique de la fibre musculaire]
\begin{enumerate}
\item \cle{Échelle macroscopique} : la fibre (cellule) est un syncytium multinucléé (noyaux périphériques), longue de quelques cm, diamètre \SI{10}{\micro m} à \SI{100}{\micro m}. Elle est entourée du \cle{sarcolemme} (membrane plasmique) invaginé en \cle{tubules T} (transverses) qui pénètrent en profondeur.
\item \cle{Échelle mésoscopique} : le \cle{hyaloplasme} (sarcoplasme) remplit l'espace entre les myofibrilles ; il contient : \textbf{glycogène} (granules de réserve), \textbf{lipides} (gouttelettes), \textbf{myoglobine} (pigment rouge fixant l'O$_2$), \textbf{mitochondries} (nombreuses chez les fibres rouges), \textbf{réticulum sarcoplasmique} (réseau tubulaire stockant les Ca$^{2+}$), enzymes glycolytiques et du cycle de Krebs.
\item \cle{Échelle microscopique} : les \cle{myofibrilles} (faisceaux parallèles) occupent \SI{80}{\%} du volume. Elles sont composées de \cle{sarcomères} en série (unité fonctionnelle, \SI{2}{\micro m} à \SI{3}{\micro m} au repos), délimités par deux \cle{stries Z} (lignes sombres). Chaque sarcomère montre une \cle{bande A} (sombre, filaments épais de \cle{myosine} + chevauchement actine/myosine), une \cle{zone H} (claire, myosine seule), et deux \cle{bandes I} (claires, actine seule).
\item \cle{Composition chimique} (pour \SI{100}{g} de muscle frais) : eau \SI{75}{g}, protéines \SI{20}{g} (myofibrillaires : myosine, actine, titine, nébuline, troponine, tropomyosine ; sarcoplasmiques : enzymes ; stromales : collagène), glucides \SI{1}{g} (glycogène), lipides \SI{1}{g} à \SI{5}{g}, ATP \SI{5}{mmol/kg}, phosphocréatine \SI{15}{mmol/kg} à \SI{20}{mmol/kg}, ions (K$^+$, Mg$^{2+}$, Ca$^{2+}$).
\end{enumerate}
\textbf{Réflexe Bac} : savoir dessiner une coupe transversale (myofibrilles polygones, noyaux périphériques, mitochondria intermyofibrillaires) et une coupe longitudinale (alternance stries claires/sombres = sarcomères).
\end{methode}

\courssub{Fiche méthode 1 : Schématiser une fibre musculaire striée au repos et en contraction}

\begin{methode}[Schématiser une portion de fibre musculaire striée squelettique]
Pour réussir ce schéma, exigible au Bac, applique la démarche suivante :
\begin{enumerate}
\item \cle{Identifier l'échelle} : on représente une \emph{myofibrille} (ou un sarcomère), c'est-à-dire la structure comprise entre deux stries Z consécutives.
\item \cle{Placer les repères fixes} : deux stries Z (lignes verticales), les filaments fins d'\cle{actine} accrochés aux stries Z, les filaments épais de \cle{myosine} au centre (bande A), la zone H (claire, myosine seule) et les bandes I (actine seule, de part et d'autre de la bande A).
\item \cle{Comparer repos et contraction} : au repos, la zone H et les bandes I sont larges ; en contraction, les filaments d'actine \emph{glissent} entre ceux de myosine : la zone H et les bandes I \textbf{rétrécissent}, la bande A garde une \textbf{longueur constante}, le sarcomère se raccourcit.
\item \cle{Soigner la légende} : chaque structure est désignée par un nombre ou une flèche, avec une légende précise placée sous le schéma (titre + numéros).
\item \cle{Ne jamais} représenter les filaments qui se raccourcissent : c'est un \emph{glissement}, pas un repliement (théorie du glissement de Huxley).
\end{enumerate}
\textbf{Réflexe Bac} : la constance de la longueur de la bande A est la preuve expérimentale du glissement des filaments.
\end{methode}

\begin{exoresolu}[Schéma d'un sarcomère au repos et en contraction]
\textbf{Énoncé.} Des micrographies électroniques d'un muscle de grenouille montrent qu'au cours de la contraction, la longueur de la bande A reste égale à \SI{1,6}{\micro m} tandis que la zone H passe de \SI{0,5}{\micro m} à \SI{0,2}{\micro m} et le sarcomère de \SI{2,5}{\micro m} à \SI{2,0}{\micro m}. (1) Réalise les schémas légendés du sarcomère au repos et en contraction. (2) Interprète ces données.
\tcblower
\textbf{Solution.}
(1) \emph{Schémas} : on trace deux rectangles représentant le sarcomère délimité par deux stries Z. Au repos : filaments fins (actine) partant de chaque strie Z sans se rejoindre au centre, laissant une zone H large ; filaments épais (myosine) centraux. En contraction : les filaments d'actine ont glissé vers le centre, leurs extrémités libres se rapprochent (voire se chevauchent) : zone H et bandes I réduites, bande A inchangée, stries Z rapprochées.
\begin{popfigure}
\begin{tikzpicture}[scale=1]
% Repos
\draw[line width=2pt,popdark] (0,0)--(0,2.2); \draw[line width=2pt,popdark] (5,0)--(5,2.2);
\node[below] at (0,0) {\footnotesize Z}; \node[below] at (5,0) {\footnotesize Z};
\foreach \y in {0.4,0.8,1.2,1.6}{
  \draw[popblue,thick] (0,\y)--(1.9,\y);
  \draw[popblue,thick] (5,\y)--(3.1,\y);
  \draw[poporange,line width=2.5pt] (2.1,\y)--(2.9,\y);}
\draw[<->,popmuted] (2.1,2.35)--(2.9,2.35); \node[above] at (2.5,2.35) {\footnotesize H};
\node[below] at (2.5,-0.5) {\footnotesize \textbf{Repos} : sarcomère \SI{2,5}{\micro m}};
% Contraction
\begin{scope}[xshift=7.5cm]
\draw[line width=2pt,popdark] (0,0)--(0,2.2); \draw[line width=2pt,popdark] (4,0)--(4,2.2);
\node[below] at (0,0) {\footnotesize Z}; \node[below] at (4,0) {\footnotesize Z};
\foreach \y in {0.4,0.8,1.2,1.6}{
  \draw[popblue,thick] (0,\y)--(1.85,\y);
  \draw[popblue,thick] (4,\y)--(2.15,\y);
  \draw[poporange,line width=2.5pt] (1.6,\y)--(2.4,\y);}
\draw[<->,popmuted] (1.85,2.35)--(2.15,2.35); \node[above] at (2,2.35) {\footnotesize H};
\node[below] at (2,-0.5) {\footnotesize \textbf{Contraction} : sarcomère \SI{2,0}{\micro m}};
\end{scope}
\end{tikzpicture}
\legende{Sarcomère au repos (gauche) et en contraction (droite). Traits verticaux épais = stries Z ; traits bleus fins = filaments d'actine ; traits orange épais = filaments de myosine (bande A) ; double flèche = zone H.}
\end{popfigure}
(2) \emph{Interprétation} : la bande A (myosine) garde une longueur constante de \SI{1,6}{\micro m} : les filaments épais ne se raccourcissent pas. La diminution de la zone H ($0{,}5 \to 0{,}2$ \si{\micro m}) et celle du sarcomère ($2{,}5 \to 2{,}0$ \si{\micro m}) s'expliquent par le \cle{glissement} des filaments fins d'actine entre les filaments épais de myosine, vers le centre du sarcomère. \textbf{Conclusion} : la contraction musculaire résulte du glissement des filaments d'actine vers le centre du sarcomère, entraînant le raccourcissement de celui-ci sans raccourcissement des filaments eux-mêmes.
\end{exoresolu}

\coursec{II. Les principaux types de fibres musculaires}

\begin{methode}[Identifier et comparer les principaux types de fibres musculaires squelettiques]
\begin{enumerate}
\item \cle{Critères de classification} : vitesse de contraction (lente/rapide), voie métabolique dominante (oxydative/glycolytique), couleur (rouge/blanche), isoformes de myosine (MHC I, IIa, IIx).
\item \cle{Type I : Fibres lentes oxydatives (Rouges, SO)} : riches en \cle{myoglobine} (couleur rouge), nombreuses \cle{mitochondries}, dense réseau \cle{capillaire}, réserves de \cle{glycogène} modérées, lipides abondants, activité ATPase myosinique faible. \textbf{Contraction lente}, \textbf{très résistantes à la fatigue}. Recrutées en premier (principe de Henneman : motoneurones de petit diamètre, seuil bas). \emph{Ex.} : muscles posturaux (soléaire), marathonien.
\item \cle{Type IIa : Fibres rapides oxydo-glycolytiques (Intermédiaires, FOG)} : myoglobine et mitochondria modérés, glycogène élevé. Vitesse rapide, résistance intermédiaire. Recrutées en second. \emph{Ex.} : 800 m, sports collectifs.
\item \cle{Type IIx (ou IIb) : Fibres rapides glycolytiques (Blanches, FG)} : pauvres en myoglobine (couleur blanche), peu de mitochondria, peu vascularisées, \textbf{très riches en glycogène} et en \textbf{enzymes glycolytiques} (PFK, PK). \textbf{Contraction très rapide et puissante}, \textbf{très fatigables} (accumulation H$^+$, Pi). Recrutées en dernier (motoneurones gros diamètre, seuil haut). \emph{Ex.} : sprinteur (100 m), haltérophile.
\end{enumerate}
\begin{center}
\begin{tabularx}{\linewidth}{|l|X|X|X|}\hline
\rowcolor{popblueL} Caractère & Type I (Lente oxydative) & Type IIa (Rapide oxydo-glycolytique) & Type IIx (Rapide glycolytique) \\ \hline
Couleur & Rouge & Rose & Blanche \\ \hline
Myoglobine / Mitochondries / Capillaires & +++ / +++ / +++ & ++ / ++ / ++ & + / + / + \\ \hline
Glycogène / Enzymes glycolytiques & + / + & +++ / ++ & +++ / +++ \\ \hline
Vitesse de contraction & Lente & Rapide & Très rapide \\ \hline
Résistance à la fatigue & Excellente & Moyenne & Faible \\ \hline
Voie ATP dominante & Aérobie (lipides/glucides) & Aérobie + Anaérobie lactique & Anaérobie alactique + lactique \\ \hline
Sport type & Marathon, fond & 800 m, football (phases longues) & 100 m, saut, force \\ \hline
\end{tabularx}
\end{center}
\textbf{À retenir} : un muscle contient un mélange de types de fibres (hétérogénéité) ; l'entraînement peut modifier les propriétés (ex. IIx $\to$ IIa par endurance) mais pas le type I $\leftrightarrow$ II.
\end{methode}

\begin{exoresolu}[Profil de fibres et performance]
\textbf{Énoncé.} Une biopsie du vaste latéral montre chez un sprinteur 75 \% de fibres de type IIx et 20 \% de type IIa, tandis qu'un marathonien présente 80 \% de fibres de type I. Interprète.
\tcblower
\textbf{Solution.} Le sprinteur possède un profil \emph{glycolytique} : fibres blanches (IIx) majoritaires, contractant vite et fort grâce à l'ATP issu de la phosphocréatine et de la glycolyse anaérobie, mais s'épuisant vite (acidose). Le marathonien a un profil \emph{oxydatif} : fibres rouges (I) majoritaires, contractant lentement mais indéfiniment grâce à l'ATP aérobie (lipides + glucides), riches en myoglobine et mitochondria. \textbf{Conclusion} : la répartition typologique, en partie génétique, oriente la prédisposition sportive ; l'entraînement affine les propriétés métaboliques (hybridation IIx/IIa).
\end{exoresolu}

\coursec{III. Le mécanisme de la contraction musculaire}

\courssub{Fiche méthode 2 : Réaliser le schéma d'interprétation du mécanisme de la contraction}

\begin{methode}[Schéma fonctionnel du mécanisme de la contraction musculaire]
Le mécanisme se résume en une chaîne d'événements qu'il faut savoir ordonner :
\begin{enumerate}
\item \cle{Arrivée de l'influx nerveux} au niveau de la jonction neuromusculaire (plaque motrice) : libération d'\cle{acétylcholine}.
\item \cle{Dépolarisation} du sarcolemme, propagée dans les tubules T.
\item \cle{Libération des ions Ca$^{2+}$} par le réticulum sarcoplasmique dans le hyaloplasme (ouverture des canaux RyR).
\item Les Ca$^{2+}$ se fixent sur la \cle{troponine C} : changement de conformation du complexe troponine-tropomyosine, qui \emph{démasque} les sites actifs de l'actine (fixation myosine).
\item Fixation des \cle{têtes de myosine} (état ADP·Pi) sur l'actine : formation des \cle{ponts acto-myosiniques}. Libération du Pi $\Rightarrow$ \cle{basculement} des têtes (power stroke) $\Rightarrow$ glissement de l'actine vers le centre (raccourcissement sarcomère). Fixation d'un nouvel ATP $\Rightarrow$ détachement de la tête. Cycle répété tant que [Ca$^{2+}$] et [ATP] sont suffisants.
\item Fin de la stimulation : \emph{pompage actif} des Ca$^{2+}$ vers le réticulum sarcoplasmique (SERCA, consomme 1 ATP par 2 Ca$^{2+}$) $\Rightarrow$ chute [Ca$^{2+}$] $\Rightarrow$ tropomyosine rebouche les sites $\Rightarrow$ \cle{relaxation}.
\end{enumerate}
\textbf{À retenir} : l'ATP intervient deux fois : (1) pour le \emph{détachement/basculement} des têtes de myosine (énergie du power stroke) et (2) pour le \emph{pompage du calcium} (relaxation). Sans ATP, les ponts restent fixés rigidement : \emph{rigidité cadavérique}.
\end{methode}

\begin{exoresolu}[Rôle du calcium et de l'ATP dans la contraction]
\textbf{Énoncé.} On réalise trois expériences sur des fibres musculaires de grenouille maintenues en survie :
\begin{center}
\begin{tabularx}{\linewidth}{|c|X|X|}\hline
\rowcolor{popblueL} Expérience & Conditions & Résultat \\ \hline
1 & Fibre stimulée électriquement, milieu normal & Contraction \\ \hline
2 & Fibre stimulée, milieu sans Ca$^{2+}$ & Pas de contraction \\ \hline
3 & Fibre stimulée, milieu avec Ca$^{2+}$ mais ATP détruit par une enzyme (apyrase) & Pas de contraction \\ \hline
\end{tabularx}
\end{center}
(1) Interprète chaque résultat. (2) Construis le schéma d'interprétation du mécanisme de la contraction.
\tcblower
\textbf{Solution.}
(1) \emph{Expérience 1} (témoin) : la stimulation électrique déclenche la contraction : la fibre est fonctionnelle. \emph{Expérience 2} : en l'absence de Ca$^{2+}$ extracellulaire (et sans libération du RS si dépolarisation bloquée), pas de contraction : le calcium est \textbf{indispensable} ; il est le \emph{signal intracellulaire} qui, en se fixant sur la troponine, démasque les sites actifs de l'actine. \emph{Expérience 3} : sans ATP, pas de contraction malgré la présence de Ca$^{2+}$ : l'ATP fournit l'\textbf{énergie} nécessaire au cycle des ponts (basculement + détachement) et au pompage du Ca$^{2+}$.
(2) \emph{Schéma d'interprétation} :
\begin{popfigure}
\begin{tikzpicture}[node distance=0.35cm, every node/.style={font=\footnotesize},
  boite/.style={draw=popblue, fill=popblueL, rounded corners, align=center, text width=3.8cm, inner sep=3pt}]
\node[boite] (a) {Influx nerveux (motoneurone $\alpha$)};
\node[boite, below=of a] (b) {Libération d'acétylcholine (plaque motrice)};
\node[boite, below=of b] (c) {Dépolarisation du sarcolemme $\to$ Tubules T};
\node[boite, below=of c] (d) {Ouverture canaux RyR $\to$ Sortie Ca$^{2+}$ du Réticulum Sarcoplasmique};
\node[boite, below=of d] (e) {Ca$^{2+}$ + Troponine C $\to$ Déplacement Tropomyosine $\to$ Sites actine démasqués};
\node[boite, below=of e] (f) {Cycle des ponts : Myosine-ADP-Pi + Actine $\to$ Power stroke (glissement) $\to$ ATP $\to$ Détachement};
\node[boite, below=of f] (g) {Arrêt stimulation $\to$ Pompe SERCA (ATP) $\to$ Ca$^{2+}$ dans RS $\to$ Relaxation};
\draw[->, poporange, thick] (a)--(b); \draw[->, poporange, thick] (b)--(c);
\draw[->, poporange, thick] (c)--(d); \draw[->, poporange, thick] (d)--(e);
\draw[->, poporange, thick] (e)--(f); \draw[->, poporange, thick] (f)--(g);
\end{tikzpicture}
\legende{Schéma d'interprétation de la contraction musculaire : de l'influx nerveux au glissement des filaments (cycle des ponts) et à la relaxation.}
\end{popfigure}
\textbf{Conclusion} : la contraction musculaire est déclenchée par l'influx nerveux, relayée par les ions Ca$^{2+}$ et rendue possible par l'hydrolyse de l'ATP ; c'est un phénomène actif, cyclique et réversible.
\end{exoresolu}

\coursec{IV. La nécessité du renouvellement de l'ATP : les voies de restauration}

\begin{definition}[ATP musculaire et phosphocréatine]
La concentration en ATP dans le muscle au repos est faible (\SIrange{4}{6}{mmol/kg} frais), ne permettant que quelques contractions maximales (\SIrange{1}{2}{s}). L'ATP est \textbf{hydrolysé} en ADP + Pi par l'ATPase myosinique (contraction) et la pompe SERCA (relaxation). Sa concentration doit être maintenue constante (\cle{homéostasie}) par une \textbf{resynthèse immédiate} à partir de trois systèmes (voies) métaboliques.
\end{definition}

\begin{propriete}[Les trois voies de restauration de l'ATP]
\begin{enumerate}
\item \cle{Via anaérobie alactique (Phosphagène)} : \ce{Phosphocreatine + ADP <=>[Creatine Kinase] Creatine + ATP}. \textbf{Vitesse :} maximale (quelques ms). \textbf{Capacité :} faible (stock PC $\approx$ \SI{20}{mmol/kg}). \textbf{Produits :} pas de lactate, pas d'O$_2$. Domine efforts $<$ \SI{10}{\second}.
\item \cle{Via anaérobie lactique (Glycolyse rapide)} : \ce{Glycogen ->[Glycogen phosphorylase] Glucose-1-P -> ... -> 2 Pyruvate + 2 ATP + 2 H+}. Pyruvate $\xrightarrow{\text{LDH}}$ \cle{Lactate} + NAD$^+$ (régénération cofacteur). \textbf{Vitesse :} élevée. \textbf{Capacité :} moyenne (stock glycogène $\approx$ \SI{400}{mmol/kg} glucosyl). \textbf{Produits :} Lactate, H$^+$ (acidose). Domine efforts \SIrange{10}{180}{\second}.
\item \cle{Via aérobie (Respiration cellulaire)} : \ce{Pyruvate + NAD+ ->[Pyruvate DH] Acetyl-CoA ->[Cycle de Krebs] CO2 + NADH/FADH2 ->[Chaine respiratoire] H2O + 34 ATP} (environ). Oxydation aussi des \cle{acides gras} ($\beta$-oxydation). \textbf{Vitesse :} lente (mise en route \SI{1}{min} à \SI{3}{min}). \textbf{Capacité :} quasi illimitée (réserves lipidiques). \textbf{Produits :} CO$_2$, H$_2$O. Domine efforts $>$ \SI{3}{min}.
\end{enumerate}
\textbf{Démonstration exigible au Bac} : bilan global de la glycolyse (2 ATP nets) et de la respiration (environ 36-38 ATP/glucose) ; rôle du lactate (régénération NAD$^+$) et de la dette d'oxygène.
\end{propriete}

\courssub{Fiche méthode 3 : Interpréter les variations de la composition chimique du muscle}

\begin{methode}[Analyser l'évolution des paramètres chimiques (ATP, phosphocréatine, glycogène, lactate)]
\begin{enumerate}
\item \cle{Décrire d'abord les courbes ou le tableau} : pour chaque paramètre, indique le sens de variation (diminution/augmentation), les valeurs initiales et finales avec \emph{unités}, et la cinétique (rapide/lente).
\item \cle{Comparer} les vitesses de variation : la phosphocréatine (PC) chute en quelques secondes, l'ATP reste presque constant au début, le glycogène diminue sur des minutes, le lactate s'accumule lors des efforts intenses.
\item \cle{Relier les variations entre elles} : la chute de la PC compense la consommation d'ATP (réaction : \ce{PC + ADP -> Creatine + ATP}) ; la baisse du glycogène et l'apparition du lactate traduisent la glycolyse anaérobie.
\item \cle{Conclure} en reliant au type de voie métabolique dominante selon la durée et l'intensité de l'effort.
\end{enumerate}
\textbf{Réflexe} : un paramètre qui \emph{reste constant} (ATP) alors que le muscle travaille signifie qu'il est \emph{renouvelé} aussi vite qu'il est consommé : c'est le cœur de la leçon.
\end{methode}

\begin{exoresolu}[Variations de l'ATP et de la phosphocréatine au cours d'un effort maximal]
\textbf{Énoncé.} Des biopsies du quadriceps sont réalisées chez un sujet effectuant un exercice maximal sur bicyclette. Résultats (en mmol par kg de muscle frais) :
\begin{center}
\begin{tabularx}{\linewidth}{|l|X|X|X|X|}\hline
\rowcolor{popblueL} Durée de l'effort (s) & 0 & 10 & 30 & 60 \\ \hline
ATP (mmol/kg) & 5,0 & 4,6 & 4,0 & 3,2 \\ \hline
Phosphocréatine (mmol/kg) & 17,0 & 8,0 & 3,0 & 1,5 \\ \hline
\end{tabularx}
\end{center}
(1) Compare l'évolution des deux paramètres. (2) Explique pourquoi l'ATP diminue beaucoup moins que la phosphocréatine. (3) Calcule le pourcentage de phosphocréatine consommée après 30 s.
\tcblower
\textbf{Solution.}
(1) \emph{Comparaison} : l'ATP passe de 5,0 à 3,2 mmol/kg en 60 s, soit une baisse de 1,8 mmol/kg ($-36\,\%$), lente et progressive. La phosphocréatine passe de 17,0 à 1,5 mmol/kg, soit une chute de 15,5 mmol/kg ($-91\,\%$), très rapide dès les 10 premières secondes (déjà $-53\,\%$ à 10 s). La PC varie donc beaucoup plus vite et plus fortement que l'ATP.
(2) \emph{Explication} : l'ATP est immédiatement \cle{resynthétisé} à partir de la phosphocréatine selon la réaction réversible catalysée par la créatine kinase :
\[ \ce{Phosphocreatine + ADP <=> Creatine + ATP} \]
Tant que le stock de PC n'est pas épuisé, la concentration d'ATP est maintenue presque constante : c'est la \emph{voie anaérobie alactique}. Quand la PC s'épuise (au-delà de 30 s), l'ATP commence à baisser nettement car sa resynthèse ne suit plus la consommation.
(3) \emph{Calcul} :
\[ \%\,\text{PC consommée} = \frac{17{,}0 - 3{,}0}{17{,}0}\times 100 = \frac{14{,}0}{17{,}0}\times 100 \approx 82{,}4\,\% \]
\textbf{Conclusion} : après 30 s d'effort maximal, environ 82 \% de la phosphocréatine est consommée ; la PC constitue la réserve énergétique immédiate qui assure le renouvellement quasi instantané de l'ATP lors des efforts brefs et violents.
\end{exoresolu}

\courssub{Fiche méthode 4 : Interpréter les données de biopsies selon la durée de l'exercice}

\begin{methode}[Relier les données de biopsie aux voies de restauration de l'ATP]
\begin{enumerate}
\item \cle{Repérer la durée et l'intensité} de l'exercice : bref et violent (quelques secondes), intense et soutenu (30 s à quelques minutes), longue durée (au-delà de quelques minutes).
\item \cle{Associer chaque situation à sa voie dominante} :
\begin{itemize}
\item effort $<$ 10--20 s : voie \cle{anaérobie alactique} (phosphocréatine) ;
\item effort intense de 20 s à 2--3 min : voie \cle{anaérobie lactique} (glycolyse $\to$ lactate) ;
\item effort prolongé : voie \cle{aérobie} (respiration cellulaire, glucose + acides gras).
\end{itemize}
\item \cle{Exploiter quantitativement} les données : quantité d'ATP produite, rendement (2 ATP par glucose en anaérobie contre environ 36--38 ATP en aérobie), rapidité de mise en jeu.
\item \cle{Conclure} sur la complémentarité des voies : elles fonctionnent toujours ensemble, mais l'une domine selon l'effort.
\end{enumerate}
\end{methode}

\begin{exoresolu}[Biopsies lors d'exercices de durées différentes]
\textbf{Énoncé.} Trois athlètes réalisent des exercices différents. Les biopsies musculaires donnent :
\begin{center}
\begin{tabularx}{\linewidth}{|l|X|X|X|}\hline
\rowcolor{popblueL} Paramètre & Sujet A (sprint 100 m) & Sujet B (course 800 m) & Sujet C (marathon) \\ \hline
Phosphocréatine & Très fortement diminuée & Diminuée & Peu diminuée \\ \hline
Glycogène & Peu diminué & Fortement diminué & Très fortement diminué \\ \hline
Lactate sanguin (mmol/L) & 8 & 18 & 4 \\ \hline
\end{tabularx}
\end{center}
(1) Attribue à chaque sujet la voie métabolique dominante, en justifiant. (2) Explique la valeur du lactate chez le sujet C.
\tcblower
\textbf{Solution.}
(1) \emph{Sujet A (100 m, environ 10 s)} : effort bref et violent ; la phosphocréatine est très fortement diminuée alors que le glycogène est peu touché : la voie dominante est la voie \cle{anaérobie alactique} (utilisation de la PC), la plus rapide mais de faible capacité. \emph{Sujet B (800 m, environ 2 min)} : effort intense et soutenu ; forte baisse du glycogène et lactate sanguin très élevé (18 mmol/L) : la voie dominante est la voie \cle{anaérobie lactique} (glycolyse rapide produisant du lactate). \emph{Sujet C (marathon, plus de 2 h)} : effort de longue durée ; glycogène très fortement diminué, PC peu diminuée, lactate faible : la voie dominante est la voie \cle{aérobie} (oxydation complète du glucose puis des acides gras en présence d'O$_2$), lente mais de très grande capacité.
(2) \emph{Lactate du sujet C} : 4 mmol/L est une valeur proche de la normale (repos : 1--2 mmol/L). En effort aérobie, le pyruvate issu de la glycolyse est \emph{oxydé} dans les mitochondries (cycle de Krebs, chaîne respiratoire) au lieu d'être réduit en lactate ; le peu de lactate produit est réutilisé (oxydation in situ, cycle de Cori hépatique). \textbf{Conclusion} : les données de biopsie confirment que le muscle adapte la voie de resynthèse de l'ATP à la durée et à l'intensité de l'effort : PC pour l'instantanéité, glycolyse lactique pour l'intensité, respiration aérobie pour l'endurance.
\end{exoresolu}

\courssub{Fiche méthode 6 : Interpréter l'évolution des paramètres physiologiques au cours de l'exercice}

\begin{methode}[Exploiter une courbe physiologique (glycogène, ATP, PC, lactate, O$_2$) en fonction du temps]
\begin{enumerate}
\item \cle{Lire les axes} : grandeur et unité en abscisse et en ordonnée ; repérer les phases de l'expérience (repos, effort, récupération).
\item \cle{Découper la courbe en phases} et décrire chaque phase avec des valeurs chiffrées précises (ex. « le glycogène passe de 90 à 40 mmol/kg entre 0 et 60 min »).
\item \cle{Interpréter chaque phase} par un mécanisme : consommation du glycogène par la glycolyse et la respiration pendant l'effort ; resynthèse de la PC et élimination du lactate pendant la récupération (qui nécessitent de l'O$_2$ : « dette d'oxygène »).
\item \cle{Relier les courbes entre elles} : la restauration de la PC à la récupération se fait grâce à l'ATP produit par la voie aérobie.
\item \cle{Conclure} par une phrase reliant l'ensemble des observations à la nécessité du renouvellement permanent de l'ATP.
\end{enumerate}
\end{methode}

\begin{exoresolu}[Évolution du glycogène et du lactate pendant et après un exercice]
\textbf{Énoncé.} Chez un sujet réalisant un exercice intense de 10 min suivi de 30 min de récupération, on mesure :
\begin{center}
\begin{tabularx}{\linewidth}{|l|X|X|X|X|X|}\hline
\rowcolor{popblueL} Temps (min) & 0 & 5 & 10 & 20 & 40 \\ \hline
Glycogène musculaire (mmol/kg) & 90 & 70 & 55 & 53 & 52 \\ \hline
Lactate sanguin (mmol/L) & 1,5 & 8 & 12 & 6 & 2 \\ \hline
\end{tabularx}
\end{center}
(1) Décris et interprète l'évolution du glycogène. (2) Décris et interprète l'évolution du lactate pendant l'effort puis pendant la récupération.
\tcblower
\textbf{Solution.}
(1) \emph{Glycogène} : pendant l'effort (0 à 10 min), le glycogène musculaire diminue fortement, de 90 à 55 mmol/kg, soit une consommation de 35 mmol/kg ; il se stabilise ensuite autour de 52--53 mmol/kg pendant la récupération. \emph{Interprétation} : pendant l'exercice intense, le glycogène est hydrolysé en glucose-1-P puis dégradé par la \cle{glycolyse} pour resynthétiser l'ATP (voies lactique et aérobie) ; à l'arrêt de l'effort, la consommation cesse, d'où la stabilisation (la resynthèse du glycogène est lente, nécessite un apport alimentaire glucidique et de l'ATP).
(2) \emph{Lactate} : il augmente pendant l'effort, de 1,5 à 12 mmol/L (multiplié par 8), car la glycolyse anaérobie, très sollicitée, produit du pyruvate en excès par rapport à la capacité oxydative, réduit en \cle{lactate} par la LDH (régénérant NAD$^+$) qui diffuse dans le sang. Pendant la récupération (10 à 40 min), le lactate diminue de 12 à 2 mmol/L : il est \emph{éliminé} grâce à l'oxygène (oxydation en CO$_2$ et H$_2$O dans le muscle cardiaque et squelettique, reconversion en glucose par le foie — cycle de Cori). Cette consommation supplémentaire d'O$_2$ après l'effort constitue la \emph{dette d'oxygène}. \textbf{Conclusion} : l'exercice intense puise dans les réserves de glycogène et s'accompagne d'une production de lactate, éliminé à la récupération grâce au métabolisme aérobie : toutes ces voies convergent vers un même objectif, le renouvellement de l'ATP.
\end{exoresolu}

\coursec{V. Adaptation du métabolisme au type d'effort}

\courssub{Fiche méthode 5 : Analyser les exigences métaboliques des différents sports}

\begin{methode}[Analyser un document sur les exigences métaboliques d'un sport]
\begin{enumerate}
\item \cle{Identifier les caractéristiques de l'effort} : durée, intensité, caractère continu ou intermittent (alternance effort/repos).
\item \cle{En déduire la ou les voies sollicitées} et le \emph{type de fibres} majoritairement recruté : fibres rapides glycolytiques (blanches, type IIx) pour les efforts explosifs ; fibres lentes oxydatives (rouges, type I) pour l'endurance ; fibres IIa pour l'intermédiaire.
\item \cle{Croiser avec les données} du document (pourcentages de filiation énergétique, lactatémie, fréquence cardiaque, consommation d'O$_2$ $\dot{V}O_2$).
\item \cle{Conclure} sur le profil métabolique du sport et, le cas échéant, sur les orientations d'entraînement ou d'alimentation (ex. ration glucidique pré-compétition pour le marathonien, créatine pour le sprinteur).
\end{enumerate}
\textbf{Astuce} : un sport intermittent (football, handball, basketball) sollicite \emph{successivement} les trois voies : ne dis jamais qu'une seule voie fonctionne ; parle de \emph{dominance} aérobie avec \emph{sollicitations répétées} des filières anaérobies.
\end{methode}

\begin{exoresolu}[Profils métaboliques du sprinteur et du footballeur]
\textbf{Énoncé.} Le tableau ci-dessous indique la contribution approximative des trois voies énergétiques dans deux activités :
\begin{center}
\begin{tabularx}{\linewidth}{|l|X|X|}\hline
\rowcolor{popblueL} Voie & Sprint 60 m & Match de football (90 min) \\ \hline
Anaérobie alactique & 80 \% & 10 \% \\ \hline
Anaérobie lactique & 15 \% & 20 \% \\ \hline
Aérobie & 5 \% & 70 \% \\ \hline
\end{tabularx}
\end{center}
(1) Analyse ces données. (2) Précise le type de fibres musculaires dominant chez le sprinteur et chez le footballeur en endurance de jeu.
\tcblower
\textbf{Solution.}
(1) \emph{Analyse} : au sprint 60 m (effort de 6--7 s, maximal), la voie anaérobie alactique fournit 80 \% de l'énergie : l'ATP est resynthétisé quasi exclusivement à partir de la phosphocréatine, voie la plus rapide (puissance maximale). Au football, la voie aérobie domine (70 \%) car le match dure 90 min à intensité moyenne (marche, jogging), mais les voies anaérobies interviennent lors des sprints, duels, accélérations (30 \% cumulés) : c'est un sport \emph{intermittent} qui sollicite les trois voies de manière répétée.
(2) \emph{Types de fibres} : le sprinteur recrute majoritairement des fibres \cle{rapides glycolytiques} (fibres blanches, type IIx) : riches en glycogène et en enzymes glycolytiques, pauvres en mitochondries et en myoglobine, puissantes mais fatigables. Le footballeur, pendant les phases d'endurance (jogging, récupération active), recrute surtout des fibres \cle{lentes oxydatives} (fibres rouges, type I) : riches en mitochondries, myoglobine et capillaires, résistantes à la fatigue ; lors des sprints, il recrute les fibres IIa et IIx. \textbf{Conclusion} : la contribution des voies énergétiques dépend directement de la structure de l'effort (durée, intensité, intermittence) ; elle conditionne le profil de fibres recrutées et explique la spécialisation musculaire des athlètes.
\end{exoresolu}

\begin{attention}[Les 5 erreurs qui coûtent des points au Bac]
\begin{enumerate}
\item \cle{Confondre glissement et raccourcissement des filaments} : les filaments d'actine et de myosine gardent une longueur constante ; c'est leur \emph{glissement} relatif qui raccourcit le sarcomère (la bande A ne change jamais de taille).
\item \cle{Dire que l'ATP « est créée » par la contraction} : c'est l'inverse, la contraction \emph{consomme} l'ATP, qui doit être \emph{resynthétisée} en permanence par les trois voies. Écris toujours « renouvellement » ou « restauration » de l'ATP.
\item \cle{Attribuer une seule voie à un sport} : les trois voies fonctionnent simultanément ; il faut parler de voie \emph{dominante} et justifier par la durée et l'intensité de l'effort.
\item \cle{Oublier les unités et les valeurs chiffrées} dans l'analyse de documents : une description sans « de 17 à 3 mmol/kg en 30 s » est insuffisante ; cite toujours les données du tableau ou de la courbe.
\item \cle{Confondre lactate et « toxine » ou cause directe des courbatures} : le lactate est un produit normal de la glycolyse anaérobie, éliminé en quelques dizaines de minutes à la récupération ; il n'est pas responsable des courbatures (micro-lésions des fibres et inflammation retardée).
\end{enumerate}
\end{attention}

\newpage

\coursec{Exercices}

\courssub{Je vérifie mes connaissances}

\begin{exercice}[QCM : structure de la fibre musculaire]
Pour chaque question, choisis la (ou les) bonne(s) réponse(s) en justifiant brièvement ton choix.
\begin{enumerate}
\item La fibre musculaire striée squelettique est :
\begin{enumerate}
\item une cellule uninucléée de petite taille ;
\item une cellule géante plurinucléée résultant de la fusion de cellules embryonnaires ;
\item un syncytium dépourvu de mitochondries ;
\item une cellule entourée d'une membrane appelée sarcolemme.
\end{enumerate}
\item Le sarcomère est délimité par :
\begin{enumerate}
\item deux stries A ; \item deux stries I ; \item deux stries Z ; \item deux bandes H.
\end{enumerate}
\item Les filaments épais du myofibrille sont constitués de :
\begin{enumerate}
\item actine ; \item myosine ; \item troponine ; \item tropomyosine.
\end{enumerate}
\item Lors de la contraction, la longueur qui diminue est celle :
\begin{enumerate}
\item de la bande A ; \item de la strie I et de la bande H ; \item des filaments de myosine ; \item du sarcomère.
\end{enumerate}
\end{enumerate}
\end{exercice}

\begin{exercice}[Vrai ou faux justifié]
Indique si chacune des affirmations suivantes est vraie ou fausse, puis justifie ta réponse en une ou deux phrases.
\begin{enumerate}
\item Le taux d'ATP dans une fibre musculaire chute brutalement dès les premières secondes d'un exercice maximal.
\item La phosphocréatine est une molécule qui stocke de l'énergie et permet de régénérer rapidement l'ATP sans oxygène.
\item La filière anaérobie lactique produit de l'ATP avec un rendement supérieur à celui de la respiration cellulaire.
\item Les fibres rouges (lentes) sont riches en myoglobine et en mitochondries et sont adaptées aux efforts d'endurance.
\item Le lactate produit pendant l'effort est un déchet toxique définitivement inutilisable par l'organisme.
\end{enumerate}
\end{exercice}

\begin{exercice}[Définitions et vocabulaire]
Définis de manière précise et concise les termes suivants, en illustrant chacun par un exemple concret :
\begin{enumerate}
\item myofibrille ; \item sarcomère ; \item phosphocréatine ; \item filière anaérobie alactique ; \item fibre musculaire intermédiaire.
\end{enumerate}
\end{exercice}

\begin{exercice}[Calculs directs]
\begin{enumerate}
\item L'hydrolyse d'une mole d'ATP libère environ \SI{30,5}{kJ}. Un sprinteur hydrolyse \num{0,40} mol d'ATP pendant un 100 m. Calcule l'énergie totale libérée.
\item Une fibre musculaire contient \SI{5}{mmol} d'ATP et \SI{17}{mmol} de phosphocréatine (PC) par kilogramme de muscle frais. Sachant qu'une molécule de PC régénère une molécule d'ATP, calcule la quantité totale d'ATP mobilisable par la filière anaérobie alactique dans \SI{25}{kg} de muscle.
\item La glycolyse produit 2 ATP par molécule de glucose, la respiration complète en produit environ 36. Calcule le rapport des rendements et commente le résultat en une phrase.
\end{enumerate}
\end{exercice}

\courssub{Je m'entraîne}

\begin{exercice}[Schéma d'un sarcomère]
À partir des données suivantes mesurées sur une microphotographie de myofibrille au repos : longueur du sarcomère = \SI{2,5}{\micro m}, bande A = \SI{1,6}{\micro m}, strie I = \SI{0,45}{\micro m} de part et d'autre de la bande A.
\begin{enumerate}
\item Calcule la longueur de la bande H au repos.
\item Réalise un schéma légendé du sarcomère au repos à l'échelle (1 cm pour \SI{0,25}{\micro m}), en représentant les filaments d'actine et de myosine, les stries Z, la bande A et la bande H.
\item En contraction maximale, la strie I mesure \SI{0,15}{\micro m} et la bande H a disparu. Calcule la nouvelle longueur du sarcomère et représente le sarcomère contracté.
\item Que constates-tu concernant la longueur de la bande A ? Quelle conséquence en tires-tu sur le mécanisme de la contraction ?
\end{enumerate}
\end{exercice}

\begin{exercice}[Évolution de la phosphocréatine]
Lors d'un exercice maximal sur ergocycle, on mesure la teneur en phosphocréatine (PC) d'un muscle de la cuisse :
\begin{center}
\begin{tabular}{|l|c|c|c|c|c|c|c|}
\hline
\rowcolor{popblueL} Temps (s) & 0 & 10 & 20 & 30 & 40 & 50 & 60 \\
\hline PC (mmol/kg de muscle) & 17 & 12 & 8 & 5 & 3 & 2 & 2 \\
\hline
\end{tabular}
\end{center}
\begin{enumerate}
\item Construis la courbe d'évolution de la PC en fonction du temps (échelles : 1 cm pour 10 s ; 1 cm pour 2 mmol/kg).
\item Décris les variations observées et calcule le pourcentage de PC consommée au bout de 30 s.
\item Explique, à l'aide de tes connaissances, pourquoi la PC diminue alors que le taux d'ATP intramusculaire reste quasiment stable pendant tout l'exercice.
\item Pourquoi la courbe se stabilise-t-elle à une valeur faible après 40 s ?
\end{enumerate}
\end{exercice}

\begin{exercice}[Composition chimique du muscle à l'effort]
On prélève par biopsie des échantillons du quadriceps d'un sujet au repos puis après 3 minutes d'exercice intense :
\begin{center}
\begin{tabular}{|l|c|c|}
\hline
\rowcolor{popblueL} Paramètre & Repos & Après 3 min d'effort \\
\hline Glycogène (mmol/kg) & 90 & 55 \\
\hline ATP (mmol/kg) & 5,0 & 4,6 \\
\hline Lactate musculaire (mmol/kg) & 1,5 & 24 \\
\hline pH intramusculaire & 7,1 & 6,7 \\
\hline
\end{tabular}
\end{center}
\begin{enumerate}
\item Calcule la variation relative (en \%) du glycogène et du lactate.
\item Interprète l'évolution de chaque paramètre en identifiant la (ou les) filière(s) énergétique(s) dominante(s) pendant cet exercice.
\item Propose une explication de la baisse du pH et de son rôle possible dans l'apparition de la fatigue.
\end{enumerate}
\end{exercice}

\begin{exercice}[Choix de la filière énergétique]
Pour chacune des situations sportives suivantes, indique la filière énergétique dominante et justifie ton choix en deux lignes maximum :
\begin{enumerate}
\item un lancer de javelot (2 s) ;
\item une course de 400 m (50 s) ;
\item un match de football de 90 min ;
\item une montée du Mont Cameroun en course d'endurance (4 h) ;
\item une série de 10 répétitions en musculation (30 s).
\end{enumerate}
\end{exercice}

\courssub{Je me prépare au Bac}

\begin{exercice}[Type Bac : biopsies musculaires au cours d'un exercice maximal]
Lors d'un test d'effort maximal de 60 s réalisé sur cycloergomètre au laboratoire de physiologie de l'INJS, des biopsies du muscle vaste externe sont effectuées chez un sujet toutes les 10 s. Les résultats sont consignés ci-dessous :
\begin{center}
\begin{tabular}{|l|c|c|c|c|c|c|c|}
\hline
\rowcolor{popblueL} Temps (s) & 0 & 10 & 20 & 30 & 40 & 50 & 60 \\
\hline ATP (mmol/kg) & 5,0 & 4,8 & 4,7 & 4,6 & 4,5 & 4,4 & 4,3 \\
\hline PC (mmol/kg) & 17,0 & 11,5 & 7,5 & 4,5 & 2,5 & 1,5 & 1,0 \\
\hline Lactate sanguin (mmol/L) & 1,0 & 1,5 & 3,0 & 6,0 & 9,5 & 12,0 & 13,5 \\
\hline
\end{tabular}
\end{center}
\begin{enumerate}
\item Construis sur un même graphique les courbes d'évolution de l'ATP et de la PC en fonction du temps (échelles : 1 cm pour 10 s ; 1 cm pour 2 mmol/kg).
\item Compare les variations des deux molécules : calcule le pourcentage de diminution de chacune entre 0 et 60 s.
\item Explique pourquoi le taux d'ATP reste quasi stable alors que la PC s'effondre. Tu préciseras la réaction biochimique impliquant la PC et l'enzyme qui la catalyse.
\item Interprète l'évolution du lactate sanguin : quelle filière énergétique prend le relais de la PC et pourquoi ?
\item Le sujet ne peut pas prolonger l'exercice au-delà de 60 s à cette intensité. Propose une explication de la fatigue musculaire en mobilisant l'ensemble des données.
\end{enumerate}
\end{exercice}

\begin{exercice}[Type Bac : fibres musculaires et aptitude sportive]
Des biopsies du muscle gastrocnémien ont été réalisées chez deux athlètes camerounais, A (spécialiste du 100 m) et B (spécialiste du marathon). Les résultats sont les suivants :
\begin{center}
\begin{tabular}{|l|c|c|}
\hline
\rowcolor{popblueL} Paramètre & Athlète A & Athlète B \\
\hline Fibres rapides blanches (\%) & 78 & 25 \\
\hline Fibres lentes rouges (\%) & 22 & 75 \\
\hline Densité capillaire (capillaires/fibre) & 1,2 & 2,8 \\
\hline Activité myosine-ATPase & élevée & faible \\
\hline Contenu en myoglobine & faible & élevé \\
\hline
\end{tabular}
\end{center}
\begin{enumerate}
\item Rappelle les caractéristiques structurales et métaboliques des fibres rouges et des fibres blanches (vitesse de contraction, filière énergétique privilégiée, résistance à la fatigue).
\item À partir des données du tableau, attribue à chaque athlète sa spécialité (sprint ou endurance) en justifiant par trois arguments tirés du tableau.
\item Explique le lien entre la densité capillaire, le contenu en myoglobine et le métabolisme aérobie des fibres de l'athlète B.
\item Un entraîneur souhaite transformer l'athlète A en marathonien. À l'aide de tes connaissances, discute la possibilité et les limites d'une telle transformation par l'entraînement.
\end{enumerate}
\end{exercice}

\begin{exercice}[Type Bac : exigences métaboliques de différents sports]
Le tableau ci-dessous indique la contribution approximative (en \%) des trois filières énergétiques à la production d'ATP dans cinq disciplines sportives :
\begin{center}
\begin{tabular}{|l|c|c|c|}
\hline
\rowcolor{popblueL} Discipline & Anaérobie alactique & Anaérobie lactique & Aérobie \\
\hline 100 m (10 s) & 90 & 8 & 2 \\
\hline 400 m (50 s) & 30 & 55 & 15 \\
\hline Lutte (5 min) & 10 & 40 & 50 \\
\hline 5000 m (14 min) & 2 & 8 & 90 \\
\hline Marathon (2 h 15) & 0 & 1 & 99 \\
\hline
\end{tabular}
\end{center}
\begin{enumerate}
\item Représente graphiquement l'évolution de la contribution de la filière aérobie en fonction de la durée de l'effort.
\item Analyse les données du tableau : quel(s) facteur(s) détermine(nt) le choix de la filière énergétique dominante ?
\item Explique pourquoi la filière anaérobie alactique domine dans le 100 m mais devient négligeable au-delà d'une minute d'effort.
\item Justifie la contribution importante de la filière anaérobie lactique dans le 400 m et dans la lutte, en précisant ses avantages et ses limites.
\item En t'appuyant sur l'ensemble de tes connaissances, explique pourquoi un sprinter ne peut pas maintenir sa vitesse maximale au-delà d'une centaine de mètres.
\end{enumerate}
\end{exercice}

\newpage

\coursec{Corrigés des exercices}

\begin{corrige}[Exercice 1 — QCM : structure de la fibre musculaire]
\begin{enumerate}
\item \textbf{Réponses b et d.} La fibre musculaire striée squelettique est une cellule géante (jusqu'à plusieurs centimètres de long) et plurinucléée : elle résulte de la fusion de cellules embryonnaires appelées \cle{myoblastes}. Elle est entourée d'une membrane plasmique particulière, le \cle{sarcolemme}. Elle est riche en mitochondries (réponse c fausse) et n'est ni petite ni uninucléée (réponse a fausse).
\item \textbf{Réponse c.} Le \cle{sarcomère}, unité contractile de la myofibrille, est le segment compris entre deux \cle{stries Z} successives. Les stries A et I et la bande H sont des zones internes du sarcomère, non ses limites.
\item \textbf{Réponse b.} Les filaments épais sont constitués de \cle{myosine} ; les filaments fins contiennent l'actine, associée à la troponine et à la tropomyosine (protéines régulatrices).
\item \textbf{Réponses b et d.} Lors de la contraction, les filaments d'actine glissent entre les filaments de myosine : la strie I et la bande H se rétrécissent, donc le sarcomère raccourcit. La bande A (= longueur des filaments de myosine) reste constante, ce qui exclut les réponses a et c.
\end{enumerate}
\end{corrige}

\begin{corrige}[Exercice 2 — Vrai ou faux justifié]
\begin{enumerate}
\item \textbf{Faux.} Le taux d'ATP intramusculaire reste quasi stable pendant l'effort car il est immédiatement régénéré, d'abord par la phosphocréatine puis par la glycolyse et la respiration. C'est la phosphocréatine qui chute brutalement.
\item \textbf{Vrai.} La phosphocréatine transfère son groupement phosphate à l'ADP selon la réaction \ce{PC + ADP ->[CK] ATP + Cr}, catalysée par la \cle{créatine kinase} (CK). Cette réaction est très rapide et ne nécessite pas d'oxygène (filière anaérobie alactique).
\item \textbf{Faux.} La filière anaérobie lactique ne produit que 2 ATP par glucose, contre environ 36 ATP pour la respiration complète. Son rendement est donc environ 18 fois plus faible ; son avantage est sa rapidité, pas son rendement.
\item \textbf{Vrai.} Les fibres rouges (lentes, de type I) sont riches en myoglobine (d'où leur couleur), en mitochondries et en capillaires sanguins ; elles utilisent préférentiellement la filière aérobie, peu puissante mais très durable : elles sont adaptées à l'endurance.
\item \textbf{Faux.} Le lactate n'est pas un déchet : il est reconverti en glucose par le foie (cycle de Cori) ou oxydé directement par le muscle cardiaque et les fibres rouges pour produire de l'ATP. C'est un intermédiaire métabolique recyclable.
\end{enumerate}
\end{corrige}

\begin{corrige}[Exercice 3 — Définitions et vocabulaire]
\begin{enumerate}
\item \textbf{Myofibrille} : structure cylindrique contractile contenue en grand nombre (jusqu'à un millier) dans le cytoplasme (sarcoplasme) d'une fibre musculaire, formée d'une succession de sarcomères. \emph{Exemple :} les myofibrilles du biceps, alignées parallèlement, donnent au muscle son aspect strié.
\item \textbf{Sarcomère} : unité structurale et fonctionnelle de la myofibrille, segment délimité par deux stries Z successives, contenant les filaments d'actine et de myosine dont le glissement assure la contraction. \emph{Exemple :} au repos, un sarcomère mesure environ \SI{2,5}{\micro\meter}.
\item \textbf{Phosphocréatine (PC)} : molécule de réserve énergétique du muscle, constituée de créatine liée à un phosphate riche en énergie ; elle régénère rapidement l'ATP à partir de l'ADP grâce à la créatine kinase. \emph{Exemple :} elle fournit l'essentiel de l'énergie d'un sprint de 100 m.
\item \textbf{Filière anaérobie alactique} : voie de régénération de l'ATP sans oxygène et sans production de lactate, utilisant la phosphocréatine ; très rapide mais épuisable en quelques secondes. \emph{Exemple :} un lancer de poids ou un départ de sprint.
\item \textbf{Fibre musculaire intermédiaire (type IIa)} : fibre rapide mais relativement résistante à la fatigue, capable d'utiliser à la fois la glycolyse anaérobie et la filière aérobie. \emph{Exemple :} fibres dominantes chez un coureur de 400 m ou de 800 m.
\end{enumerate}
\end{corrige}

\begin{corrige}[Exercice 4 — Calculs directs]
\begin{enumerate}
\item Énergie libérée : $E = \num{0,40} \times 30,5 = \SI{12,2}{kJ}$. Le sprinteur libère donc environ \SI{12,2}{kJ} pendant son 100 m.
\item Quantité d'ATP mobilisable par kilogramme de muscle : $5 + 17 = \SI{22}{mmol/kg}$ (ATP initial + ATP régénéré par la PC). Pour \SI{25}{kg} de muscle : $22 \times 25 = \SI{550}{mmol} = \SI{0,55}{mol}$ d'ATP mobilisables par la filière anaérobie alactique.
\item Rapport des rendements : $\dfrac{36}{2} = 18$. La respiration cellulaire produit environ \textbf{18 fois plus} d'ATP par glucose que la seule glycolyse anaérobie : la filière aérobie est bien plus rentable, ce qui explique qu'elle soit privilégiée lors des efforts prolongés.
\end{enumerate}
\end{corrige}

\begin{corrige}[Exercice 5 — Schéma d'un sarcomère]
\begin{enumerate}
\item La bande A (\SI{1,6}{\micro\meter}) comprend la bande H centrale plus les zones de chevauchement actine-myosine. Les filaments d'actine pénètrent dans la bande A sur une longueur égale à la strie I de chaque côté, soit $2 \times 0,45 = \SI{0,90}{\micro\meter}$. Donc :
\[ H = 1,6 - 0,90 = \SI{0,70}{\micro\meter}. \]
\item À l'échelle 1 cm pour \SI{0,25}{\micro\meter} : sarcomère = 10 cm ; bande A = 6.4cm ; chaque strie I = 1.8cm ; bande H = 2.8cm. Le schéma représente deux stries Z verticales distantes de 10 cm, les filaments fins d'actine accrochés aux stries Z et s'étendant sur 1,8 cm vers le centre, les filaments épais de myosine au centre (6,4 cm), la zone centrale sans actine (bande H, 2,8 cm).
\begin{popfigure}
\begin{center}
\begin{tikzpicture}[scale=0.9]
% Stries Z
\draw[line width=1.5pt,popdark] (0,-0.9)--(0,0.9);
\draw[line width=1.5pt,popdark] (10,-0.9)--(10,0.9);
\node[below,popdark] at (0,-0.95) {\small Strie Z};
\node[below,popdark] at (10,-0.95) {\small Strie Z};
% Myosine (bande A) : de 1.8 à 8.2
\foreach \y in {-0.45,-0.15,0.15,0.45}{\draw[line width=2.2pt,poporange] (1.8,\y)--(8.2,\y);}
\foreach \x in {2.4,3.6,4.8,6.0,7.2}{\draw[poporange,line width=1pt] (\x,0.45)--(\x,0.75); \draw[poporange,line width=1pt] (\x,-0.45)--(\x,-0.75);}
% Actine
\foreach \y in {-0.3,0,0.3}{\draw[line width=1pt,popblue] (0,\y)--(1.8,\y); \draw[line width=1pt,popblue] (8.2,\y)--(10,\y);}
% Légendes de zones
\draw[<->,popmuted] (1.8,1.15)--(8.2,1.15); \node[above,popmuted] at (5,1.15) {\small Bande A (1,6 \si{\micro\meter})};
\draw[<->,popmuted] (3.6,-1.25)--(6.4,-1.25); \node[below,popmuted] at (5,-1.25) {\small Bande H (0,70 \si{\micro\meter})};
\draw[<->,popmuted] (0,1.15)--(1.8,1.15); \node[above,popmuted] at (0.9,1.15) {\small Strie I};
\end{tikzpicture}
\end{center}
\legende{Schéma du sarcomère au repos (échelle : 1 cm pour \SI{0,25}{\micro\meter}). En orange : filaments épais de myosine (avec ponts) ; en bleu : filaments fins d'actine ancrés sur les stries Z ; la bande H correspond à la zone centrale de myosine sans actine.}
\end{popfigure}
\item En contraction maximale : nouvelle longueur du sarcomère = bande A + 2 stries I = $1,6 + 2 \times 0,15 = \SI{1,9}{\micro\meter}$. Le schéma contracté montre les stries Z rapprochées (7,6 cm à l'échelle), les filaments d'actine se rejoignant presque au centre (bande H nulle).
\item La bande A garde une longueur constante (\SI{1,6}{\micro\meter}) au repos comme en contraction : les filaments de myosine ne raccourcissent pas. Conclusion : la contraction résulte du \cle{glissement des filaments fins d'actine entre les filaments épais de myosine}, qui rapproche les stries Z (théorie du glissement des filaments de Huxley).
\end{enumerate}
\end{corrige}

\begin{corrige}[Exercice 6 — Évolution de la phosphocréatine]
\begin{enumerate}
\item Courbe construite en reportant les points (0 ; 17), (10 ; 12), (20 ; 8), (30 ; 5), (40 ; 3), (50 ; 2), (60 ; 2), avec le temps en abscisse (1 cm = 10 s) et la PC en ordonnée (1 cm = 2 mmol/kg).
\begin{popfigure}
\begin{center}
\begin{tikzpicture}
\begin{axis}[width=0.85\linewidth,height=6cm,xlabel={Temps (s)},ylabel={PC (mmol/kg)},xmin=0,xmax=65,ymin=0,ymax=18,grid=both,grid style={popline!40},legend style={at={(0.97,0.97)},anchor=north east}]
\addplot[popcurve,mark=*,mark size=1.5pt] coordinates {(0,17)(10,12)(20,8)(30,5)(40,3)(50,2)(60,2)};
\legend{Phosphocréatine}
\end{axis}
\end{tikzpicture}
\end{center}
\legende{Évolution de la teneur en phosphocréatine musculaire au cours d'un exercice maximal de 60 s.}
\end{popfigure}
\item La PC diminue rapidement et de façon quasi linéaire pendant les 30 premières secondes (chute d'environ 4 mmol/kg par 10 s), puis la décroissance ralentit et la courbe se stabilise vers 2 mmol/kg. Pourcentage consommé à 30 s :
\[ \frac{17 - 5}{17} \times 100 = \frac{12}{17} \times 100 \approx 70,6\,\%. \]
Environ \textbf{71 \%} de la PC est consommée au bout de 30 s.
\item La PC sert de réservoir immédiat de phosphate énergétique : la réaction \ce{PC + ADP ->[CK] ATP + Cr} régénère l'ATP au fur et à mesure de son hydrolyse par la contraction. L'ATP est donc renouvelé instantanément : son taux reste stable tant que la PC n'est pas épuisée, tandis que la PC, elle, se consomme.
\item Après 40 s, le stock de PC est pratiquement épuisé (valeur résiduelle d'environ 2 mmol/kg, non utilisable efficacement car la vitesse de la réaction dépend de la concentration en PC). La filière anaérobie alactique devient insuffisante : ce sont alors la filière anaérobie lactique puis la filière aérobie qui prennent le relais pour régénérer l'ATP.
\end{enumerate}
\end{corrige}

\begin{corrige}[Exercice 7 — Composition chimique du muscle à l'effort]
\begin{enumerate}
\item Variation du glycogène : $\dfrac{55 - 90}{90} \times 100 \approx -38,9\,\%$, soit une baisse d'environ \textbf{39 \%}. Variation du lactate : $\dfrac{24 - 1,5}{1,5} \times 100 = +1500\,\%$, soit une multiplication par 16 (augmentation de \textbf{1500 \%}).
\item \textbf{Glycogène} : forte consommation ($-39\,\%$), car il est le substrat glucidique dégradé par la glycolyse pour régénérer l'ATP. \textbf{ATP} : quasi stable ($-8\,\%$ seulement), car il est continuellement régénéré. \textbf{Lactate} : accumulation massive ($\times 16$), signature d'une \cle{filière anaérobie lactique} très active : l'apport d'oxygène est insuffisant pour oxyder tout le pyruvate produit, qui est alors réduit en lactate. \textbf{pH} : baisse de 7,1 à 6,7, conséquence de l'accumulation de lactate (acide lactique) et des ions H$^+$. Conclusion : pendant ces 3 minutes intenses, la filière dominante est la \textbf{filière anaérobie lactique}, relayant la filière anaérobie alactique épuisée, avec une contribution croissante de la filière aérobie.
\item La baisse du pH (acidose musculaire) provient de la dissociation de l'acide lactique en lactate et ions H$^+$. Les ions H$^+$ inhibent les enzymes de la glycolyse (notamment la phosphofructokinase), perturbent la fixation du Ca$^{2+}$ sur la troponine et gênent la formation des ponts actine-myosine : la puissance de contraction diminue, c'est la \cle{fatigue musculaire}.
\end{enumerate}
\end{corrige}

\begin{corrige}[Exercice 8 — Choix de la filière énergétique]
\begin{enumerate}
\item \textbf{Lancer de javelot (2 s)} : filière \textbf{anaérobie alactique}. Effort bref et violent : l'ATP est régénéré instantanément par la phosphocréatine, sans oxygène ni lactate.
\item \textbf{Course de 400 m (50 s)} : filière \textbf{anaérobie lactique} dominante. Effort intense et soutenu : les stocks de PC s'épuisent en moins de 10 s, la glycolyse anaérobie fournit rapidement l'ATP avec production de lactate.
\item \textbf{Match de football (90 min)} : filière \textbf{aérobie} dominante, avec des phases anaérobies lors des sprints. Effort de longue durée entrecoupé d'accélérations : la respiration cellulaire couvre l'essentiel des besoins.
\item \textbf{Montée du Mont Cameroun (4 h)} : filière \textbf{aérobie} exclusivement ou presque. Effort prolongé d'intensité modérée : oxydation complète du glucose et des lipides, très rentable en ATP.
\item \textbf{Série de 10 répétitions en musculation (30 s)} : filière \textbf{anaérobie lactique} (après épuisement rapide de la PC). Effort intense de durée intermédiaire : la glycolyse anaérobie fournit l'ATP rapidement malgré l'accumulation de lactate.
\end{enumerate}
\end{corrige}

\begin{corrige}[Exercice 9 — Type Bac : biopsies musculaires au cours d'un exercice maximal]
\begin{enumerate}
\item Graphique : abscisse = temps (1 cm = 10 s), ordonnée = concentration (1 cm = 2 mmol/kg) ; on trace les deux courbes à partir des points du tableau.
\begin{popfigure}
\begin{center}
\begin{tikzpicture}
\begin{axis}[width=0.85\linewidth,height=6.5cm,xlabel={Temps (s)},ylabel={Concentration (mmol/kg)},xmin=0,xmax=65,ymin=0,ymax=18,grid=both,grid style={popline!40},legend style={at={(0.97,0.97)},anchor=north east}]
\addplot[popcurve,mark=*,mark size=1.5pt] coordinates {(0,17)(10,11.5)(20,7.5)(30,4.5)(40,2.5)(50,1.5)(60,1)};
\addplot[popcurve,popgreen,mark=square*,mark size=1.5pt] coordinates {(0,5)(10,4.8)(20,4.7)(30,4.6)(40,4.5)(50,4.4)(60,4.3)};
\legend{PC, ATP}
\end{axis}
\end{tikzpicture}
\end{center}
\legende{Évolution comparée des teneurs en ATP et en phosphocréatine (PC) du muscle vaste externe au cours d'un exercice maximal de 60 s.}
\end{popfigure}
\item Diminution de l'ATP : $\dfrac{5,0 - 4,3}{5,0} \times 100 = 14\,\%$. Diminution de la PC : $\dfrac{17,0 - 1,0}{17,0} \times 100 \approx 94,1\,\%$. Comparaison : la PC chute de façon spectaculaire (environ 94 \%) alors que l'ATP ne diminue que de 14 \% : les deux courbes sont quasi symétriques en sens inverse au début, ce qui suggère que la PC sert à régénérer l'ATP.
\item L'ATP hydrolysé par la contraction est immédiatement resynthétisé grâce à la phosphocréatine selon la réaction :
\[ \ce{PC + ADP ->[CK] ATP + Cr}. \]
Tant que la PC est disponible, l'ATP est maintenu à un niveau quasi constant : c'est donc la PC qui « paie » le coût énergétique à la place de l'ATP. Quand la PC s'épuise (après 40 s), l'ATP commence à décliner.
\item Le lactate sanguin reste faible pendant les 10 premières secondes (1,5 mmol/L) puis augmente fortement et régulièrement jusqu'à 13,5 mmol/L à 60 s. Cette accumulation traduit la mise en route de la \cle{filière anaérobie lactique} : quand la PC s'épuise et que l'apport en O$_2$ reste insuffisant pour l'intensité de l'effort, la glycolyse anaérobie prend le relais ; le pyruvate produit en excès est réduit en lactate, qui diffuse dans le sang.
\item La fatigue s'explique par la conjonction de plusieurs facteurs tirés des données : (a) épuisement quasi total de la PC (94 \% consommés), privant le muscle de sa source d'ATP la plus rapide ; (b) début de baisse de l'ATP, limitant directement le fonctionnement des ponts actine-myosine ; (c) accumulation massive de lactate (acidose) qui inhibe les enzymes glycolytiques et la contraction. Le muscle ne peut plus régénérer l'ATP à la vitesse exigée par un exercice maximal : l'intensité doit chuter.
\end{enumerate}
\textbf{Barème indicatif (sur 10)} : construction correcte du graphique (axes légendés, échelles, points) : 2 pts ; calculs des pourcentages : 1,5 pt ; réaction PC + ADP avec enzyme citée : 2 pts ; interprétation du lactate et filière relais : 2 pts ; explication de la fatigue mobilisant les trois paramètres : 2,5 pts.\par
\textbf{Points de vigilance} : citer obligatoirement la \cle{créatine kinase} et écrire l'équation de la réaction ; ne pas écrire que « l'ATP ne varie pas » mais « reste quasi stable grâce à sa régénération » ; relier explicitement chaque conclusion aux données chiffrées du tableau.
\end{corrige}

\begin{corrige}[Exercice 10 — Type Bac : fibres musculaires et aptitude sportive]
\begin{enumerate}
\item \textbf{Fibres rouges (lentes, type I)} : contraction lente, petit diamètre, riches en myoglobine, en mitochondries et entourées de nombreux capillaires ; filière privilégiée : \textbf{aérobie} ; très résistantes à la fatigue. \textbf{Fibres blanches (rapides, type IIb)} : contraction rapide et puissante, gros diamètre, pauvres en myoglobine et en mitochondries, riches en glycogène et en enzymes glycolytiques ; filière privilégiée : \textbf{anaérobie lactique} ; fatigables rapidement.
\item L'athlète A est le \textbf{sprinter} et l'athlète B le \textbf{marathonien}. Arguments : (a) A possède 78 \% de fibres blanches rapides, adaptées aux efforts brefs et violents, contre 75 \% de fibres rouges lentes chez B, adaptées à l'endurance ; (b) la densité capillaire de B (2,8 capillaires/fibre contre 1,2) garantit un apport élevé en O$_2$, indispensable au métabolisme aérobie des longues distances ; (c) le contenu en myoglobine élevé de B favorise le stockage et le transfert de l'O$_2$ vers les mitochondries, tandis que l'activité myosine-ATPase élevée de A traduit une hydrolyse rapide de l'ATP, donc une contraction rapide, typique du sprint.
\item La myoglobine capte l'O$_2$ apporté par le sang et le transfère aux mitochondries ; une densité capillaire élevée assure un débit sanguin important, donc un approvisionnement continu en O$_2$ et en substrats (glucose, acides gras) ainsi que l'évacuation du CO$_2$. Ces deux caractères permettent aux fibres rouges de B de maintenir une \textbf{respiration cellulaire} intense et prolongée : production abondante d'ATP (environ 36 ATP/glucose) sans accumulation de lactate, d'où une excellente résistance à la fatigue.
\item L'entraînement d'endurance peut améliorer chez A les caractères aérobies : augmentation de la densité capillaire, du nombre de mitochondries, des enzymes oxydatives, et transformation partielle des fibres intermédiaires (IIa) vers un profil plus oxydatif. \textbf{Limites} : la proportion de fibres rouges et blanches est en grande partie déterminée génétiquement ; les fibres blanches IIb ne se transforment pas en véritables fibres rouges. Avec 78 \% de fibres blanches, l'athlète A pourra progresser en endurance mais n'atteindra vraisemblablement pas le niveau d'un marathonien de haut niveau comme B.
\end{enumerate}
\textbf{Barème indicatif (sur 10)} : tableau comparatif des deux types de fibres : 2,5 pts ; attribution justifiée par trois arguments : 3 pts ; lien capillaires--myoglobine--aérobie : 2,5 pts ; discussion nuancée entraînement/déterminisme génétique : 2 pts.\par
\textbf{Points de vigilance} : justifier chaque attribution par des \textbf{données chiffrées du tableau} (pourcentages, densité capillaire) ; ne pas affirmer que l'entraînement transforme totalement le typage des fibres ; citer le rôle de la myoglobine dans le stockage de l'O$_2$.
\end{corrige}

\begin{corrige}[Exercice 11 — Type Bac : exigences métaboliques de différents sports]
\begin{enumerate}
\item On représente la contribution aérobie (2, 15, 50, 90, 99 \%) en fonction de la durée de l'effort (10 s, 50 s, 5 min, 14 min, 135 min) ; l'axe des durées pouvant être gradué en échelle ordinale croissante.
\begin{popfigure}
\begin{center}
\begin{tikzpicture}
\begin{axis}[width=0.85\linewidth,height=6cm,xlabel={Durée de l'effort (min, échelle non linéaire)},ylabel={Contribution aérobie (\%)},symbolic x coords={0.17,0.83,5,14,135},xtick=data,ymin=0,ymax=105,grid=both,grid style={popline!40},legend style={at={(0.03,0.97)},anchor=north west}]
\addplot[popcurve,mark=*,mark size=1.5pt] coordinates {(0.17,2)(0.83,15)(5,50)(14,90)(135,99)};
\legend{Filière aérobie}
\end{axis}
\end{tikzpicture}
\end{center}
\legende{Évolution de la contribution de la filière aérobie à la production d'ATP en fonction de la durée de l'effort (0,17 min = 10 s ; 0,83 min = 50 s ; 135 min = 2 h 15).}
\end{popfigure}
La contribution aérobie croît fortement avec la durée de l'effort, de 2 \% (10 s) à 99 \% (marathon).
\item Deux facteurs déterminent la filière dominante : la \textbf{durée} de l'effort (plus l'effort est long, plus l'aérobie domine) et son \textbf{intensité} (plus l'effort est violent, plus les filières anaérobies, rapides mais peu rentables, sont sollicitées). En effet, les filières anaérobies fournissent l'ATP très rapidement mais en quantité limitée, tandis que la filière aérobie, lente à se mettre en place, produit beaucoup d'ATP tant que l'apport en O$_2$ suffit.
\item Dans le 100 m, l'effort est maximal et très bref (10 s) : seule la filière anaérobie alactique (PC + créatine kinase) régénère l'ATP assez vite, sans attendre la glycolyse ni l'oxygénation. Mais le stock de PC est faible (environ 17 mmol/kg) et s'épuise en moins de 10 s d'effort maximal : au-delà d'une minute, sa contribution devient négligeable (0 à 2 \%).
\item Dans le 400 m et la lutte, l'effort est intense et dure de 50 s à 5 min : la PC est épuisée et l'apport en O$_2$ reste insuffisant pour couvrir seul les besoins. La filière anaérobie lactique prend le relais : \textbf{avantages} — mise en route rapide, puissance élevée, fonctionnement sans O$_2$ ; \textbf{limites} — faible rendement (2 ATP/glucose), épuisement du glycogène et accumulation de lactate responsable de l'acidose et de la fatigue, ce qui la limite à quelques minutes.
\item La vitesse maximale exige une production d'ATP extrêmement rapide, que seules les filières anaérobies assurent. Or la PC s'épuise en quelques secondes et la filière anaérobie lactique, qui prend le relais, provoque une accumulation de lactate et une acidose qui inhibent la contraction et la glycolyse elle-même. La filière aérobie, trop lente à fournir l'ATP, ne peut pas soutenir une intensité maximale. Le sprinter doit donc inévitablement ralentir au-delà d'une centaine de mètres.
\end{enumerate}
\textbf{Barème indicatif (sur 10)} : graphique correct : 2 pts ; identification des facteurs durée/intensité : 2 pts ; explication du 100 m (stock de PC limité) : 2 pts ; avantages et limites de la filière lactique : 2 pts ; synthèse sur la limitation de la vitesse maximale : 2 pts.\par
\textbf{Points de vigilance} : toujours associer chaque filière à son \textbf{rendement} (2 ou 36 ATP/glucose) et à sa \textbf{vitesse de production} ; citer l'épuisement du stock de PC comme cause du déclin de la filière alactique ; ne pas présenter le lactate comme un simple déchet.
\end{corrige}

\newpage

\coursec{Fiche bilan}

\begin{popfigure}
\begin{tikzpicture}[
  mindmap,
  grow cyclic,
  every node/.style={concept, text=white, font=\small\bfseries},
  root concept/.append style={concept color=popink, font=\large\bfseries},
  level 1/.append style={level distance=3.6cm, sibling angle=60, font=\footnotesize},
  level 2/.append style={level distance=2.4cm, font=\scriptsize},
]
\node[root concept]{ATP et exercice musculaire}
  child[concept color=popblue]{ node {Structure de la fibre}
    child { node {Myofibrilles : actine + myosine} }
    child { node {Sarcomère : striations} }
  }
  child[concept color=popgreen]{ node {Types de fibres}
    child { node {Rouges : lentes, aérobies} }
    child { node {Blanches : rapides, anaérobies} }
  }
  child[concept color=poporange]{ node {Contraction}
    child { node {Ponts actine--myosine} }
    child { node {ATP + Ca$^{2+}$ indispensables} }
  }
  child[concept color=poppurple]{ node {Voie anaérobie alactique}
    child { node {Phosphocréatine} }
    child { node {0--10 s, effort bref} }
  }
  child[concept color=poppink]{ node {Voie anaérobie lactique}
    child { node {Glycolyse $\to$ lactate} }
    child { node {10 s--3 min} }
  }
  child[concept color=popgold]{ node {Voie aérobie}
    child { node {Respiration cellulaire} }
    child { node {Efforts longs, 36 ATP} }
  };
\end{tikzpicture}
\legende{Carte mentale de la leçon : le renouvellement de l'ATP est au cœur de la performance musculaire ; chaque type d'effort mobilise une voie métabolique privilégiée.}
\end{popfigure}

\begin{bilan}
\textbf{1. Définitions clés}
\begin{itemize}
  \item \cle{Fibre musculaire striée squelettique} : cellule multinucléée contenant des \cle{myofibrilles} formées de filaments d'\cle{actine} (fins) et de \cle{myosine} (épais), organisés en \cle{sarcomères} (unité contractile entre deux stries Z).
  \item \cle{Contraction musculaire} : raccourcissement des sarcomères par glissement des filaments d'actine entre ceux de myosine, grâce aux \cle{ponts actine--myosine} ; elle consomme de l'ATP et nécessite du Ca$^{2+}$.
  \item \cle{Phosphocréatine (PC)} : molécule de réserve musculaire qui reconstitue rapidement l'ATP : $\text{PC} + \text{ADP} \rightarrow \text{créatine} + \text{ATP}$.
  \item \cle{Lactate} : produit de la glycolyse anaérobie ; son accumulation acidifie le muscle et participe à la fatigue.
\end{itemize}

\textbf{2. Bilans énergétiques à connaître par cœur}
\begin{center}
\begin{tabularx}{\linewidth}{|l|X|X|}
\hline
\rowcolor{popblueL} \textbf{Voie} & \textbf{Bilan simplifié} & \textbf{Rendement / durée} \\
\hline
Anaérobie alactique & $\text{PC} + \text{ADP} \rightarrow \text{C} + \text{ATP}$ & Très rapide ; 0 à \SI{10}{s} \\
\hline
Anaérobie lactique & \ce{C6H12O6 -> 2 lactate + 2 ATP} & Rapide ; \SI{10}{s} à \SI{3}{min} \\
\hline
Aérobie & \ce{C6H12O6 + 6O2 -> 6CO2 + 6H2O + 36 ATP} & Lente mais durable ; efforts longs \\
\hline
\end{tabularx}
\end{center}

\textbf{3. Types de fibres}
\begin{itemize}
  \item \cle{Fibres rouges (lentes)} : riches en myoglobine et mitochondries, métabolisme aérobie, résistantes à la fatigue (endurance, marathon).
  \item \cle{Fibres blanches (rapides)} : pauvres en mitochondries, métabolisme anaérobie, puissantes mais fatigables (sprint, lancer).
\end{itemize}

\textbf{4. Adaptation au type d'effort}
\begin{itemize}
  \item \cle{Effort bref et violent} (sprint, \SI{100}{m}) : voie anaérobie alactique (phosphocréatine).
  \item \cle{Effort intense et soutenu} (\SI{400}{m}--\SI{1500}{m}) : voie anaérobie lactique.
  \item \cle{Effort de longue durée} (marathon, football) : voie aérobie (glycogène puis lipides).
\end{itemize}

\textbf{5. Méthodes express}
\begin{itemize}
  \item Courbe ATP/PC en fonction du temps : une chute rapide de la PC sans chute de l'ATP $\Rightarrow$ la PC reconstitue l'ATP (voie alactique).
  \item Biopsie : comparer les taux avant/après effort ; identifier la voie dominante selon la durée de l'exercice.
  \item Schéma de contraction : toujours légender stries Z, filaments fins/épais, sens du glissement.
\end{itemize}
\end{bilan}

\begin{aretenir}[Checklist avant le Bac]
\begin{itemize}
  \item $\square$ Je sais décrire et schématiser un sarcomère au repos et en contraction (stries Z, actine, myosine).
  \item $\square$ Je connais le rôle de l'ATP et du Ca$^{2+}$ dans la formation des ponts actine--myosine.
  \item $\square$ Je sais que les stocks d'ATP musculaire sont très faibles et doivent être régénérés en permanence.
  \item $\square$ Je connais l'équation de la réaction phosphocréatine + ADP $\rightarrow$ créatine + ATP.
  \item $\square$ Je connais le bilan de la glycolyse anaérobie : 1 glucose $\rightarrow$ 2 lactate + 2 ATP.
  \item $\square$ Je connais le bilan de la respiration cellulaire : \ce{C6H12O6 + 6O2 -> 6CO2 + 6H2O + 36 ATP}.
  \item $\square$ Je sais associer chaque voie métabolique à une durée d'effort (alactique : 0--\SI{10}{s} ; lactique : \SI{10}{s}--\SI{3}{min} ; aérobie : au-delà).
  \item $\square$ Je sais distinguer fibres rouges (aérobies, endurance) et fibres blanches (anaérobies, puissance).
  \item $\square$ Je sais interpréter une courbe d'évolution de l'ATP, de la phosphocréatine ou du glycogène au cours d'un exercice.
  \item $\square$ Je sais exploiter des données de biopsies musculaires pour identifier la voie énergétique dominante.
  \item $\square$ Je soigne mes schémas : légende numérotée, titre, vocabulaire anatomique exact.
  \item $\square$ Je rédige mes interprétations en démarche scientifique : observation $\rightarrow$ hypothèse $\rightarrow$ exploitation des données $\rightarrow$ conclusion.
\end{itemize}
\end{aretenir}

\findecours

\end{document}
